\chapter{Monge--Amp\`ere energy and mixed measures}
\index{Monge--Ampere energy!BJ}\index{Monge--Ampere measure!mixed BJ}
\label{chap:bj-energy}

The envelope theorem changes the scope of the trivially valued theory, but
not its intersection-theoretic origin.  Energy is first an intersection
number on one test configuration.  Monotone regularization then extends it
to singular psh functions, and polarization produces mixed
Monge--Amp\`ere measures.  This order of construction is important: no
local wedge product of singular currents is hidden in the definition.

This is also the chapter in which the trivially valued theory becomes
usable.  \Crefrange{chap:bj-cartier-realization}{chap:bj-subintegral-normalization}
were forbidden to use the envelope property, because they were proving it;
\Cref{chap:bj-envelope-consequences} discharged that prohibition
(\Cref{imp:ch45-firewall-discharged}) and turned the obstacle envelope
$P_\theta$ into an operator with values in $\PSH_X(\theta)$
(\Cref{imp:ch45-obstacle-operator}).  From here to the end of
\Cref{part:bj} these results are available under their stated hypotheses.
The present chapter constructs energy and mixed measures: the passage from
a finitary intersection number to a functional on singular psh functions
uses the unconditional fixed-ample-class regularization of
\Cref{thm:ch34-fs-regularization}.  The envelope theorem enters later
through compactness, extremal functions, and obstacle calculus.
\Cref{sec:ch46-where-envelope} isolates those uses, and
\Cref{rem:ch46-without-envelope} records what remains without them.

\section{Orientation}
\label{sec:ch46-orientation}

\subsection{What is constructed, and out of what}
\label{sec:ch46-what}

The object constructed here is a single symmetric multilinear pairing of
$n+1$ pairs $(\theta,\phi)$, a numerical class together with a
$\theta$-psh function.  Everything else in the chapter is a specialization
or a polarization of it:

\begin{itemize}
\item taking all $n+1$ pairs equal to $(\omega,\phi)$ and dividing by
      $(n+1)V$ gives the \emph{Monge--Amp\`ere energy} $E_\omega(\phi)$;
\item taking the first pair to be $(0,f)$ with $f$ a PL function gives
      the \emph{mixed Monge--Amp\`ere measure} of the remaining $n$ pairs,
      through the Riesz representation theorem;
\item taking two pairs of the form $(0,f)$ and $(0,g)$ gives the
      \emph{Dirichlet form}, whose symmetry is integration by parts and
      whose sign is the Hodge index inequality.
\end{itemize}

The construction proceeds in two steps, and the two steps use disjoint
inputs.  The first step is finitary: on PL data the pairing is an
$(n+1)$-fold intersection number on a compactified test configuration,
and every property of it --- symmetry, multilinearity, the behaviour under
constants, monotonicity, the Lipschitz estimate, the atomic formula for the
associated measure --- is an identity or an inequality of intersection
numbers on a projective $(n+1)$-dimensional $k$-scheme.  The inputs are
\Cref{thm:ch33-compactification,thm:ch33-intersection-invariance} and the
Hodge index inequality for divisors supported in one fibre.  No analysis
occurs, and no envelope.

The second step is a monotone extension, and this is where analysis enters.
The extension is by an infimum over PL psh competitors lying above the
given psh functions.  One must prove that the competitor family is
nonempty, that the infimum reproduces the intersection number on PL data,
and that the resulting functional is continuous along decreasing nets.
Fixed-class PL regularization from \Cref{thm:ch34-fs-regularization},
together with a compact-space Dini argument, supplies these facts; the
envelope property is not needed for this foundational extension.

\subsection{The envelope theorem is available here}
\label{sec:ch46-envelope-available}

\begin{important}{Regularization and envelope inputs}
\label{imp:ch46-envelope-imported}
Throughout this chapter $X$ is integral projective and \emph{unibranch},
so that \Cref{thm:bj-final-envelope} applies to every real class.  The two
inputs below have different provenance, which matters for tracking what
survives without unibranchness.
\begin{enumerate}
\item\label{itm:ch46-imp:regularization}
      \emph{Fixed-class monotone approximation.}  For
      $\omega\in\Amp(X)$ and $u\in\PSH_X(\omega)$ there is a downward
      directed net in $\cH_{\mathrm{dom}}(\omega)$, hence of continuous
      PL $\omega$-psh functions, with $u_g\downarrow u$ pointwise on
      $X^{\An}$ (\Cref{thm:ch34-fs-regularization}).  This is
      unconditional and valid for every real ample class.  The envelope
      construction \Cref{thm:ch45-continuous-regularization} gives a second
      such net, indexed by continuous obstacles.
\item\label{itm:ch46-imp:compactness}
      \emph{Compactness.}  The normalized space
      $\PSH^{\sup}_X(\theta)$ of \eqref{eq:normalized-psh-space} is
      compact for pointwise convergence on $X^{\divis}$, for every real
      class $\theta$ (\Cref{cor:ch45-compactness}).  Before
      \Cref{chap:bj-envelope-consequences} this was the conditional
      statement of \Cref{warn:ch34-compactness-conditional}.
\end{enumerate}
In addition the calculus of $P_\theta$
(\Cref{prop:bj-final-obstacle-calculus,prop:ch45-obstacle-calculus-2}), the
extremal function $V_\theta=P_\theta(0)$ of \Cref{cor:ch45-extremal}, and
the monotone statements
\Cref{thm:ch45-decreasing-obstacle-nets,thm:ch45-increasing-obstacles,cor:ch45-uniform-obstacles}
are quoted rather than reproved.
\end{important}

The reader who wants to see the dependence made explicit should turn to
\Cref{sec:ch46-where-envelope}, where the inputs are separated by
provenance, and to \Cref{rem:ch46-without-envelope}, which
describes the theory that survives if the envelope theorem is withheld.
The short version is: the PL theory and the monotone energy extension to
unbounded psh functions survive intact; compactness, extremal envelopes,
and the variational existence package do not.

\subsection{Standing hypotheses and notation}
\label{sec:ch46-hypotheses}

Throughout this chapter $X$ is an integral projective variety of pure
dimension $n$ over an algebraically closed trivially valued field $k$,
and $X$ is unibranch. Statements involving two differential slots or
$n-1$ auxiliary classes assume $n\ge1$. We fix an ample numerical class

\[
                    \omega\in\Amp(X),\qquad V:=(\omega^n)>0.
\tag{46.1}\label{eq:bj-volume}
\]

Since $X$ is fixed we abbreviate $\PSH(\theta):=\PSH_X(\theta)$ and
$\PSH^{\sup}(\theta):=\PSH^{\sup}_X(\theta)$, and we write
$\cH_{\mathrm{dom}}(\omega)$ for the set of \Cref{def:ch34-hdom}, that is,
the PL functions $\phi_D$ with $D\in\operatorname{VCar}(\cX)_{\mathbf Q}$
on a test configuration $\cX$ dominating $X\times\mathbf A^1$ and
$\omega_{\cX}+D$ relatively ample.  The subspace of PL functions that are
$\theta$-psh is written

\[
 \PL^{\theta}(X):=\PL_{\mathbf R}(X)\cap\PSH(\theta),
\]

which by \Cref{prop:ch34-pl-consistency} is the same thing whether pshness
is tested by \Cref{def:bj-psh} or by the relative nefness condition
\eqref{eq:pl-relative-nefness} of \Cref{def:positive-pl-function}.  The
unibranch hypothesis is used only for the envelope-dependent inputs in
\Cref{imp:ch46-envelope-imported}; in particular it is not used for the
fixed-class regularization in item~\cref{itm:ch46-imp:regularization}.
Every statement about PL data holds for an arbitrary integral projective
$X$, and the wider scope of the monotone extension is recorded statement
by statement in \Cref{sec:ch46-hypothesis-use}.

\subsection{The normalization, declared}
\label{sec:ch46-normalization}

The results through \Cref{thm:bj-mixed-ma} are those of
\cite[Sections~3 and 7]{BJ22}; the envelope theorem of
\Crefrange{chap:bj-envelope-statement}{chap:bj-envelope-consequences} is
not needed for the finitary parts of them, and
\Cref{sec:ch46-where-envelope} says exactly where it is needed for the
rest.  BJ normalize the pure Monge--Amp\`ere measure to
mass $1$.  Thus in this chapter

\[
 \MA_\omega(\phi):=V^{-1}(\omega+\ddc\phi)^n,
 \qquad
 \MA^{\mathrm{u}}_\omega(\phi):=V\MA_\omega(\phi).
\tag{46.2}\label{eq:bj-ma-normalizations}
\]

The superscript ``u'' means unnormalized.  It is the second measure that
has total mass $V$, in agreement with the global convention of this book.

\begin{warning}{A declared departure from the global convention}
\label{warn:ch46-normalization}
The front matter fixes the total mass of a Monge--Amp\`ere measure to be
$(L^n)$ unless the word \emph{probability} occurs, and
\eqref{eq:appJ-model-ma}--\eqref{eq:appJ-prob-ma} record the two
normalizations side by side.  This chapter departs from that convention in
the single respect displayed in \eqref{eq:bj-ma-normalizations}: the
unadorned symbol $\MA_\omega$ denotes the \emph{probability} measure, and
the unnormalized measure of mass $V$ carries the superscript
$\mathrm u$.  Consequently the unadorned energy $E_\omega$ of
\Cref{def:bj-ma-energy} is the one that satisfies
$E_\omega(\phi+c)=E_\omega(\phi)+c$, not
$E_\omega(\phi+c)=E_\omega(\phi)+cV$.  Both conventions are legitimate and
both are recorded in \Cref{sec:appJ-energy}: \eqref{eq:appJ-energy} is the
unnormalized one used in \Cref{part:bfj}, and the sentence following
\eqref{eq:appJ-energy-translation} is precisely the licence used here.  The
translation rules

\[
 \begin{aligned}
  E_\omega(\phi+c)&=E_\omega(\phi)+c,\\
  V\cdot E_\omega(\phi)=E^{\mathrm u}_\omega(\phi)
   &:=\frac{1}{n+1}\sum_{j=0}^{n}
     \int_{X^{\An}}\phi\,(\omega+\ddc\phi)^j
       \wedge\omega^{\,n-j}
 \end{aligned}
\]

are the two tests that detect a mixture of the conventions immediately, and
they are applied in every worked example below.  \emph{The word
``probability'' is never used with a class of volume zero}, which is why
$\omega$ is required to be ample in \eqref{eq:bj-volume}; the pairing
itself is defined for arbitrary classes and is where zero-volume statements
are made.
\end{warning}

Two consistency checks of \eqref{eq:bj-ma-normalizations} against
\Cref{app:normalizations} are carried out in the text: the model formula
\eqref{eq:bj-model-ma-atomic} is \eqref{eq:appJ-cl-measure} with
$\varpi$-multiplicities $b_E$ and normalized divisorial points
$v_E=b_E^{-1}\ord_E$ exactly as there
(\Cref{def:ch33-component-valuation}); and the first-variation test
\eqref{eq:appJ-energy-variation} is verified in
\Cref{prop:ch46-first-variation-pl}, in the normalized form
$\frac{\rd}{\rd t}E_\omega(\phi_t)=\int\dot\phi_t\,\MA_\omega(\phi_t)$
with $\MA_\omega$ a probability measure.

\subsection{Extended-real ledger}
\label{sec:ch46-ledger}

The following entries of the extended-real ledger govern the limiting
arguments and differences in this chapter.

\begin{enumerate}
\item A decreasing net of extended-real functions has a pointwise limit in
      $\Rbar$; no hypothesis is needed.  This is what makes
      \eqref{eq:bj-energy-infimum} an infimum of a genuine family rather
      than a limit that must be shown to exist.
\item A $\theta$-psh function is bounded above and is finite on
      $X^{\divis}$, by \Cref{def:bj-psh}\,(2),
      \eqref{eq:normalized-psh-space} and \Cref{thm:ch34-psh-compactness};
      so a difference $\phi-\psi$ of $\theta$-psh functions is defined
      on $X^{\divis}$, but need not be bounded there. It may be
      undefined at a point where both functions are $-\infty$.
      A difference with a continuous function is defined everywhere and
      bounded above when the psh function is the first term. Differences
      of finite-energy functions in integrals are formed almost everywhere,
      using the integrability proved in \Cref{thm:bj-mixed-ma}.
\item Multiplication by a strictly positive real preserves
      $\Rbar$-suprema and infima, which is what makes the scaling action
      \eqref{eq:function-scaling-action} interact correctly with the
      pairing in \Cref{sec:ch46-homogeneous}.
\end{enumerate}

The prohibition on $(-\infty)-(-\infty)$ is the reason
\Cref{warn:ch46-no-subtraction} exists and is the reason the multilinear
extension of \Cref{sec:ch46-e1} is asserted on $\vec\cE^{\,1}(X)$ and
nowhere else.

\subsection{Relation to \texorpdfstring{\Cref{part:bfj}}{Part III}}
\label{sec:ch46-part3}

This chapter is the trivially valued counterpart of
\Cref{chap:bfj-ma,chap:bfj-energy}.  The parallel is close enough that the
differences are worth tabulating before they are met.

\begin{center}
\begin{tabular}{@{}p{.22\textwidth}p{.73\textwidth}@{}}
\toprule
feature & comparison with \Cref{part:bfj}\\
\midrule
ambient &
The discretely valued pairing uses a proper model over $\Spec R$;
here it uses a test configuration compactified over $\mathbf P^1$
(\Cref{thm:ch33-compactification}). Both have dimension $n+1$.\\
fibre class &
The special fibre $\cX_s$ there and $\overline\pi^*\{0\}$ here have
model function $1$ and pairing $(\theta_1\cdots\theta_n)$ with the
remaining horizontal classes; see \eqref{eq:ch20-pairing-as-integral}
and \eqref{eq:ch33-intersection-twist}.\\
PL energy &
The mixed vertical sum in \eqref{eq:model-energy} corresponds to the
compactified intersection in \Cref{def:bj-pl-energy-pairing} and
\eqref{eq:bj-energy-definition}. The conversion is exhibited in
\Cref{ex:ch46-same-number-twice}.\\
normalization &
Part III uses $E(\varphi+c)=E(\varphi)+cV$; here
$E(\phi+c)=E(\phi)+c$ (\Cref{warn:ch46-normalization}).\\
Dirichlet form &
Both forms are minus vertical intersection: compare
\Cref{prop:ch26-dirichlet-vertical,prop:ch46-dirichlet-vertical}.\\
scaling &
Here the trivial absolute value gives a scaling action on $X^{\An}$,
compatible with the $\mathbf G_m$-equivariant test-configuration
formalism (\Cref{sec:ch46-homogeneous}). Scaling valuations does not
preserve a fixed nontrivial absolute value on the discretely valued base.\\
approximation &
General limits use nets. Sequential reductions require a stated reason
(\Cref{rem:ch22-legitimate-sequences,warn:ch45-nets}); explicit curve
examples may still be sequences.\\
\bottomrule
\end{tabular}
\end{center}

Arguments using multilinearity, the fibre-class identity and vertical
Hodge index transfer with $\cX_s$ replaced by
$\overline\pi^*\{0\}$, after converting the normalization. The scaling
action supplies the additional identities proved in
\Cref{sec:ch46-homogeneous}. \Cref{ex:ch46-same-number-twice} computes
one energy in both valuation regimes.

\subsection{Roadmap}
\label{sec:ch46-roadmap}

\Cref{sec:ch46-model} constructs the PL pairing and proves its formal
properties.  \Cref{sec:ch46-pl-measure} converts it into a measure at PL
level, proves the atomic formula, integration by parts and the Hodge index
inequality, and identifies the energy of a PL metric with the first moment
of a Duistermaat--Heckman measure.  \Cref{sec:ch46-p1} computes a single
example on $\mathbf P^1$ in three independent ways.
\Cref{sec:ch46-extension} performs the monotone extension and proves the
theorem announced as \Cref{thm:bj-extended-energy-pairing}.
\Cref{sec:ch46-e1} treats finite energy and polarization,
\Cref{sec:ch46-mixed} the mixed measures, \Cref{sec:ch46-ibp} locality and
integration by parts, and \Cref{sec:ch46-homogeneous} the scaling action.
\Cref{sec:ch46-examples} contains the toric computation, the comparison
with \Cref{part:bfj}, and the discontinuity example that motivates
\Cref{chap:bj-strong-topologies}.  \Cref{sec:ch46-scope} audits the
hypotheses, \Cref{sec:ch46-dictionary} is the chapter dictionary, and
\Cref{sec:ch46-source} the source note.

\section{The model-level intersection number}
\label{sec:ch46-model}

\subsection{Common models}
\label{sec:ch46-common-models}

Let $(\cX_i,\cL_i)$ represent an $L_i$-psh PL metric
function $\phi_i$, and choose one integrally closed test configuration
$\cX$ dominating all $\cX_i$.  Write

\[
 \cL_i=\rho_i^*c_1(L_i{}_{\mathbf A^1})+D_i,
 \qquad D_i\in\operatorname{VCar}(\cX)_{\mathbf R}.
\tag{46.3}\label{eq:bj-common-model-pairs}
\]

Here and below an expression involving a vertical divisor is pulled to a
common dominating test configuration before being intersected.  Three
statements make this legitimate and are used without further comment.

First, the set of test configurations dominating $X\times\mathbf A^1$ is
directed under domination, by \Cref{prop:ch33-directed-poset}, so a common
$\cX$ exists for any finite family; and it may be taken integrally
closed by \Cref{prop:ch33-integral-closure-exists}, which changes neither
the associated PL functions nor the intersection numbers.  Second, the
correspondence $\phi\leftrightarrow D$ between PL functions and vertical
divisors on the limit is bijective, by \Cref{thm:tc-pl-correspondence};
this is what makes \eqref{eq:bj-common-model-pairs} a statement about the
functions $\phi_i$ rather than about a chosen presentation.  Third, a
test configuration is compactified canonically over $\mathbf P^1$, and
the polarizations extend canonically, by \Cref{thm:ch33-compactification};
this is what makes a degree available.  The compactification need not be a product:
\Cref{warn:ch33-compactification-not-product} is worth rereading here,
because the whole force of \eqref{eq:bj-trivial-energy-zero} below is that
the \emph{trivial} test configuration is the one case in which the
compactification is a product.

\begin{definition}[Energy pairing on PL pairs]
\label{def:bj-pl-energy-pairing}
For $\theta_i=c_1(L_i)\in N^1(X)$ and
$\phi_i=\phi_{D_i}\in\PL(X)_{\mathbf R}$, set

\[
 (\theta_0,\phi_0)\cdots(\theta_n,\phi_n)
 :=(\cL_0\cdots\cL_n)_{\overline{\cX}},
\tag{46.4}\label{eq:bj-pl-energy-intersection}
\]

where the test configuration is compactified over $\mathbf P^1$ using
the fixed trivialization away from the origin. For arbitrary real
$\theta_i$, choose real Cartier representatives and extend the formula
by real multilinearity. Thus $\cL_i=\theta_{i,\cX}+D_i$ is a real
Cartier class; numerical independence is proved below.
\end{definition}

The right-hand side does not depend on the chosen common model.  Indeed,
if $\pi:\cX'\to\cX$ is another dominating model, the
projection formula gives

\[
 (\pi^*\cL_0\cdots\pi^*\cL_n)_{\overline{\cX}' }
       =(\cL_0\cdots\cL_n)_{\overline{\cX}}.
\tag{46.5}\label{eq:bj-energy-projection}
\]

Thus \Cref{def:bj-pl-energy-pairing} is intrinsic to the PL
functions.  It is symmetric and multilinear in the pairs
$(\theta_i,\phi_i)$.  Translation in one metric variable gives

\[
 (\theta_0,\phi_0+c)\prod_{i=1}^n(\theta_i,\phi_i)
 = (\theta_0,\phi_0)\prod_{i=1}^n(\theta_i,\phi_i)
    +c(\theta_1\cdots\theta_n).
\tag{46.6}\label{eq:bj-energy-translation-pairing}
\]

The three assertions just made deserve proof, since everything in the
chapter rests on them.

\begin{proposition}
\label{prop:ch46-pairing-wellposed}
\Cref{def:bj-pl-energy-pairing} depends only on the pairs
$(\theta_i,\phi_i)\in N^1(X)_{\mathbf R}\times\PL_{\mathbf R}(X)$, not on
the model, the divisors or the line bundles chosen to present them.  The
resulting map
\[
 \bigl(N^1(X)_{\mathbf R}\times\PL_{\mathbf R}(X)\bigr)^{n+1}
 \longrightarrow\mathbf R
\]
is symmetric and $\mathbf R$-multilinear in the $n+1$ pairs, and
satisfies \eqref{eq:bj-energy-translation-pairing} and
\eqref{eq:bj-trivial-energy-zero}.
\end{proposition}

\begin{proof}
\emph{Independence of the model.}  Let $\cX'\to\cX$ be a domination and
$\pi:\overline{\cX'}\to\overline\cX$ its unique extension, which exists
by \Cref{thm:ch33-compactification}\,(3) and satisfies
$\overline{\pi^*\cL}=\pi^*\overline\cL$.  The morphism $\pi$ is
projective and birational, so $\pi_*[\overline{\cX'}]=[\overline\cX]$,
and the projection formula \Cref{thm:appF-projection} gives
\eqref{eq:bj-energy-projection}; this is
\eqref{eq:ch33-intersection-domination}.  Given two presentations of the
same data on models $\cX_1,\cX_2$, pass to a common domination by
\Cref{prop:ch33-directed-poset}; the pulled-back divisors agree because
the map $D\mapsto\phi_D$ is injective on the limit
(\Cref{thm:tc-pl-correspondence}, Step~2), and both numbers equal the one
computed upstairs.

\emph{Multilinearity.}  On a fixed $\overline\cX$ the number
$(\cL_0\cdots\cL_n)$ is the mixed intersection number of $n+1$
$\mathbf R$-Cartier classes, which is symmetric and multilinear by
\Cref{prop:appF-mixed-cartier}.  It remains to note that the map
$(\theta,\phi)\mapsto\cL=\theta_{\cX}+D_\phi$ is additive: for fixed
$\cX$, $\theta\mapsto\theta_{\cX}$ is the linear pullback
$N^1(X)_{\mathbf R}\to N^1(\overline\cX)_{\mathbf R}$ along
$\overline\cX\to X$, and $\phi\mapsto D_\phi$ is linear by
\Cref{thm:tc-pl-correspondence}.  Hence the composite is multilinear in
the pairs.

\emph{Constants.}  A rational constant $c$ is the PL function
$\phi_{c\,\cX_0}$, because $\phi_{\cX_0}=1$: indeed
$\phi_{\cX_0}(v_E)=b_E^{-1}\ord_E(\cX_0)=1$ for every component $E$ by
\eqref{eq:ch33-bE}, and a PL function is determined by its values at the
$v_E$ on a model determining it.  On the compactification,
$\cO(\cX_0)=\overline\pi^*\cO_{\mathbf P^1}(\{0\})$.  Adding $c$ to
$\phi_0$ therefore replaces $\cL_0$ by
$\cL_0+c\,\overline\pi^*\{0\}$, and
\[
\begin{aligned}
 \bigl((\cL_0+c\,\overline\pi^*\{0\})\cdot\cL_1\cdots\cL_n\bigr)
 &=(\cL_0\cdots\cL_n)
   +c\bigl(\overline\pi^*\{0\}\cdot\cL_1\cdots\cL_n\bigr)\\
 &=(\cL_0\cdots\cL_n)+c(\theta_1\cdots\theta_n),
\end{aligned}
\]
the last equality because $\overline\pi^*\{0\}$ is a fibre class:
it is numerically equivalent to $\overline\pi^*\{\infty\}$, the fibre at
infinity is $X\times\{\infty\}$, and $\cL_i$ restricts there to
$L_i$, so the intersection number is $(\theta_1\cdots\theta_n)_X$.
This is \eqref{eq:ch33-intersection-twist} in polarized form, and it proves
\eqref{eq:bj-energy-translation-pairing} for rational $c$; both sides are
affine in $c$, so it holds for real $c$.

\emph{The trivial metric.}  If all $\phi_i=0$ one may take
$\cX=X\times\mathbf A^1$, $\overline\cX=X\times\mathbf P^1$ and
$\cL_i=p_1^*L_i$.  Then
$(\cL_0\cdots\cL_n)=\bigl(p_1^*(\theta_0\cdots\theta_n)\bigr)$, and
$\theta_0\cdots\theta_n$ is a class of dimension $-1$ on the
$n$-dimensional $X$; concretely, the product of $n+1$ classes pulled
back from an $n$-dimensional variety vanishes.  This is
\eqref{eq:bj-trivial-energy-zero}.
\end{proof}

\begin{remark}[The fibre class is the whole difference with \Cref{part:bfj}]
\label{rem:ch46-fibre-class}
The proof of \eqref{eq:bj-energy-translation-pairing} used exactly one
property of $\overline\pi^*\{0\}$: that its model function is the
constant $1$ and that
$(\overline\pi^*\{0\}\cdot\cL_1\cdots\cL_n)=(\theta_1\cdots\theta_n)$.
Over a discretely valued base, \Cref{sec:ch20-pairing} records the same two
properties for $\cX_s$: $\varphi_{\cX_s}=1$, and taking
$D=\cX_s$ in \eqref{eq:ch20-pairing-as-integral} gives the total mass
$(\{\theta_1\}\cdots\{\theta_n\})$.  Every argument in
\Cref{chap:bfj-energy} that uses only multilinearity and these two
properties therefore transfers verbatim, and this is why so much of the
present section reads like \Cref{sec:ch26-model-energy}.  What does
\emph{not} transfer is that $\cX_s$ is the divisor of a uniformizer,
whereas $\overline\pi^*\{0\}$ is not principal on $\overline\cX$; the
compensating structure on the trivially valued side is the
$\mathbf G_m$-action, and it is exploited in
\Cref{sec:ch46-homogeneous}.
\end{remark}

\subsection{Monotonicity and the Hodge index inequality}
\label{sec:ch46-monotonicity}

The pairing is not merely multilinear: on psh data it is monotone, and the
mechanism is the negativity of vertical self-intersections.  We isolate
that mechanism first, because it recurs four times below --- in
monotonicity, in concavity of energy, in the tangent inequalities, and in
the Cauchy--Schwarz estimate.

\begin{lemma}
\label{lem:ch46-vertical-negativity}
Let $\cX$ be a test configuration dominating $X\times\mathbf A^1$, let
$F$ be a vertical $\mathbf R$-Cartier divisor on $\cX$, and let
$\cL_1,\dots,\cL_{n-1}$ be $\mathbf R$-Cartier classes on
$\overline\cX$ that are relatively nef over $\mathbf P^1$.  Then
\[
 \bigl(F^2\cdot\cL_1\cdots\cL_{n-1}\bigr)_{\overline\cX}\le0 .
\]
Consequently, for vertical $F,G$,
\[
 \bigl|(F\cdot G\cdot\cL_1\cdots\cL_{n-1})\bigr|^2
 \le
 \bigl(F^2\cdot\cL_1\cdots\cL_{n-1}\bigr)
 \bigl(G^2\cdot\cL_1\cdots\cL_{n-1}\bigr).
\]
\end{lemma}

\begin{proof}
\emph{Step 1: perturbation and base twist.}  Let $A$ be an ample class on
$\overline\cX$ and put $F_m:=\overline\pi^*\{0\}$.  Since $F$ is
supported in the fibre $\overline\pi^{-1}(0)$ and $F_m$ is numerically
the fibre over any other point, $F\cdot F_m=0$ as a cycle class; hence
for any real numbers $m_i$,
\[
 \bigl(F^2\cdot(\cL_1+m_1F_m)\cdots(\cL_{n-1}+m_{n-1}F_m)\bigr)
 =\bigl(F^2\cdot\cL_1\cdots\cL_{n-1}\bigr),
\]
every term containing an $F_m$ vanishing. For $\varepsilon>0$, the
classes $\cL_i+\varepsilon A$ are relatively ample. Each becomes ample
on $\overline\cX$ after adding a sufficiently large multiple of the
fibre class. Apply Step~2 to these twisted classes, remove the twists by
the displayed identity, and let $\varepsilon\downarrow0$. Continuity of
intersection numbers gives the assertion for the original relatively nef
classes. Relative nefness alone is not used to assert that a base twist
is globally nef.

\emph{Step 2: the Hodge index inequality for a divisor in one fibre.}  The
scheme $\overline\cX$ is projective and integral of dimension $n+1$ by
\Cref{lem:ch33-central-fibre-cartier}\,(3) and
\Cref{thm:ch33-compactification}\,(1), and $\overline\pi$ maps it onto
the smooth projective curve $\mathbf P^1$.  For a divisor supported in a
single fibre of such a morphism, and nef classes
$\cL_1,\dots,\cL_{n-1}$, the inequality
$(F^2\cdot\cL_1\cdots\cL_{n-1})\le0$ is the mixed relative Hodge index
inequality; over a discretely valued base it is
\Cref{thm:ch20-mixed-hodge}, \eqref{eq:ch20-mixed-hodge}, and the proof
given there uses only that the ambient scheme is projective over a
one-dimensional base, that $F$ is supported in one fibre, and that the
auxiliary classes are nef in the fibre directions.  In the present setting
it is the vertical negativity inequality of \cite[Lemma~3.24]{BJ22}.

\emph{Step 3: Cauchy--Schwarz.}  By Step~2 the symmetric bilinear form
$(F,G)\mapsto-(F\cdot G\cdot\cL_1\cdots\cL_{n-1})$ on the finite
dimensional space of vertical $\mathbf R$-Cartier divisors on $\cX$ is
positive semidefinite.  The Cauchy--Schwarz inequality for a positive
semidefinite form gives the second display.
\end{proof}

\begin{warning}{Relative nefness, not nefness, is what psh gives}
\label{warn:ch46-relative-nef}
The class $\cL_i=\theta_{i,\cX}+D_i$ attached to a $\theta_i$-psh PL
function is relatively nef over $\mathbf A^1$ by
\eqref{eq:pl-relative-nefness}, and it is relatively nef over
$\mathbf P^1$ because over $\mathbf P^1\setminus\{0\}$ it is the
pullback of $\theta_i$ from $X$ along a product.  It is \emph{not}
nef on $\overline\cX$ in general: adding a negative multiple of the
central fibre to a relatively ample class preserves relative ampleness
but can make its degree on a horizontal curve negative. The background
$\theta_i$ is nef, as required for a bounded psh function.
The perturbation and base twist in Step~1 of
\Cref{lem:ch46-vertical-negativity} is the standard device for
neutralizing this, and it works only because the divisor in the first two
slots is vertical.  An inequality with a horizontal class in one of the two
distinguished slots is false in general and is never used.
\end{warning}

\begin{proposition}\label{prop:bj-model-energy-monotonicity}
Let $\theta_0$ be arbitrary, let $\phi_0,\psi_0\in\PL_{\mathbf R}(X)$
with $\phi_0\leq\psi_0$, and let $(\theta_i,\phi_i)$ be psh PL
pairs for $1\le i\le n$. Then

\[
 (\theta_0,\phi_0)\prod_{i=1}^n(\theta_i,\phi_i)
 \leq
 (\theta_0,\psi_0)\prod_{i=1}^n(\theta_i,\phi_i).
\tag{46.7}\label{eq:bj-model-monotonicity}
\]
\end{proposition}

\begin{proof}
On a common model, $\psi_0-\phi_0$ corresponds to a vertical divisor
whose associated model function is nonnegative.  The intersection of that
divisor with the nef restrictions of the remaining model line bundles is
nonnegative.  This is the vertical negativity/positivity calculation of
\cite[Lemma~3.15]{BJ22}.  Multilinearity reduces the displayed inequality
to precisely that calculation.

We write the reduction out.  Choose an integrally closed test
configuration $\cX$ dominating $X\times\mathbf A^1$ and determining
$\phi_0,\psi_0,\phi_1,\dots,\phi_n$, and let $D_0,D_0',D_1,\dots,D_n$
be the corresponding vertical divisors, so that $\cL_i=\theta_{i,\cX}+D_i$
is relatively nef for $1\le i\le n$.  Put $F:=D_0'-D_0$.  By
multilinearity,
\[
 (\theta_0,\psi_0)\prod_{i=1}^n(\theta_i,\phi_i)
 -(\theta_0,\phi_0)\prod_{i=1}^n(\theta_i,\phi_i)
 =\bigl(F\cdot\cL_1\cdots\cL_n\bigr)_{\overline\cX}.
\]
The function $\phi_F=\psi_0-\phi_0$ is $\ge0$, so $F$ is an effective
divisor by the effectivity criterion \Cref{lem:ch33-effectivity}, applied
on the integrally closed $\cX$.  Write $F=\sum_Ea_EE$ with
$a_E=b_E\phi_F(v_E)\ge0$, the sum over the components of $\cX_0$.  Each
$E$ is a projective $k$-scheme contained in a fibre of
$\overline\pi$, so $\cL_i|_E$ is nef, and therefore
$(\cL_1|_E\cdots\cL_n|_E)\ge0$ by \Cref{prop:appF-mixed-cartier}.  Hence
\[
 \bigl(F\cdot\cL_1\cdots\cL_n\bigr)
 =\sum_Ea_E\bigl(\cL_1|_E\cdots\cL_n|_E\bigr)\ \ge0 .
\]
No positivity of $\theta_0$ was used, and no relation between
$\theta_0$ and the other classes; what was used is that the $n$
classes in the \emph{other} slots are relatively nef.
\end{proof}

\begin{corollary}
\label{cor:ch46-pl-lipschitz}
Let $(\theta_i,\phi_i)$ and $(\theta_0,\psi_0)$ be psh at PL level.
Then
\[
 \Bigl|(\theta_0,\psi_0)\prod_{i=1}^n(\theta_i,\phi_i)
      -(\theta_0,\phi_0)\prod_{i=1}^n(\theta_i,\phi_i)\Bigr|
 \le\lVert\psi_0-\phi_0\rVert_\infty\,(\theta_1\cdots\theta_n).
\]
In particular, for $\phi,\psi\in\PL^{\omega}(X)$,
$|E_\omega(\psi)-E_\omega(\phi)|\le\lVert\psi-\phi\rVert_\infty$.
\end{corollary}

\begin{proof}
Put $\delta:=\lVert\psi_0-\phi_0\rVert_\infty$, a finite number since PL
functions are bounded (\Cref{lem:ch34-pl-lattice}\,\cref{itm:ch34-pl:lattice}).
Then $\phi_0-\delta\le\psi_0\le\phi_0+\delta$, and all three functions
are $\theta_0$-psh at PL level, constants being harmless by
\Cref{prop:ch34-pl-operations}\,\cref{itm:ch34-plop:const}.  Apply
\Cref{prop:bj-model-energy-monotonicity} twice and then
\eqref{eq:bj-energy-translation-pairing}:
\[
 (\theta_0,\phi_0)\prod-\delta(\theta_1\cdots\theta_n)
 \le(\theta_0,\psi_0)\prod
 \le(\theta_0,\phi_0)\prod+\delta(\theta_1\cdots\theta_n).
\]
For the last assertion take all pairs equal, use the telescoping identity
\eqref{eq:ch46-telescope} below $n+1$ times --- or, more simply, note
that \Cref{def:bj-ma-energy} divides the pure pairing by $(n+1)V$ and
apply the first assertion in each of the $n+1$ slots, each contributing
at most $\lVert\psi-\phi\rVert_\infty V$.
\end{proof}

\begin{example}[The trivial metric]
\label{ex:ch46-trivial-metric}
If all $\phi_i=0$, then the compactified line bundles are pulled back
from $X$.  Since $n+1$ horizontal classes are intersected on
$X\times\mathbf P^1$,

\[
                    (\theta_0,0)\cdots(\theta_n,0)=0.
\tag{46.8}\label{eq:bj-trivial-energy-zero}
\]

Formula~\eqref{eq:bj-energy-translation-pairing} is therefore the correct
way in which constants enter the pairing.  Combining the two,
\[
 (\theta_0,c_0)\cdots(\theta_n,c_n)
 =\Bigl(\sum_{i=0}^nc_i\Bigr)(\theta^n)
 \quad\text{if }n\ge1\text{ and all }\theta_i=\theta;
\]
note that the right-hand factor is the $n$-fold product $(\theta^n)$,
not the $(n+1)$-fold product, which vanishes by
\eqref{eq:bj-trivial-energy-zero};
in general the expansion of \eqref{eq:bj-energy-translation-pairing} in all
$n+1$ slots gives
$\sum_{i}c_i(\theta_0\cdots\widehat{\theta_i}\cdots\theta_n)$, the hat
denoting omission.  Setting all $c_i=c$ and all $\theta_i=\omega$
recovers $E_\omega(c)=c$, which is the normalization test of
\Cref{warn:ch46-normalization}.
\end{example}

\subsection{The telescoping identity}
\label{sec:ch46-telescope}

The pure pairing $(\omega,\phi)^{n+1}$ is what the energy is built from,
and the identity that converts it into integrals is elementary
multilinearity.  We record it once, in the form used throughout.

\begin{lemma}
\label{lem:ch46-telescope}
For pairs $(\theta,\phi)$ and $(\theta,\psi)$ psh at PL level,
\[
 (\theta,\phi)^{n+1}-(\theta,\psi)^{n+1}
 =\sum_{j=0}^{n}(0,\phi-\psi)\,(\theta,\phi)^{j}(\theta,\psi)^{n-j}.
\tag{46.9}\label{eq:ch46-telescope}
\]
In particular, taking $\psi=0$ and using
\eqref{eq:bj-trivial-energy-zero},
\[
 (\omega,\phi)^{n+1}
 =\sum_{j=0}^{n}\int_{X^{\An}}\phi\;
   (\omega+\ddc\phi)^{j}\wedge\omega^{\,n-j},
\tag{46.10}\label{eq:ch46-pure-pairing-as-integrals}
\]
the integrals being those of \Cref{prop:ch46-pl-measure}.
\end{lemma}

\begin{proof}
Write $a:=(\theta,\phi)$ and $b:=(\theta,\psi)$ as formal symbols in
the multilinear pairing, and $f:=(0,\phi-\psi)=a-b$, the last equality
being multilinearity in the pair together with additivity of
$\phi\mapsto D_\phi$.  Then the identity
$a^{n+1}-b^{n+1}=(a-b)\sum_{j=0}^na^jb^{n-j}$ holds in the polynomial
ring generated by the pairs, hence after applying the symmetric
multilinear form.  The second display follows from
\eqref{eq:bj-trivial-energy-zero} and the definition
\eqref{eq:bj-mixed-ma-duality} of the mixed measure, which at PL level is
\Cref{prop:ch46-pl-measure}.
\end{proof}

\begin{remark}[The corresponding computation in \Cref{part:bfj}]
\label{rem:ch46-telescope-part3}
Over a discretely valued base the corresponding identity is the mixed
vertical energy difference, not a difference of absolute powers.  Namely,
if $u=\phi-\psi=\varphi_U$, where $U$ is the actual vertical divisor on
a common model representing $u$, then
\[
 E_\omega(\phi)-E_\omega(\psi)
 =\frac1{n+1}\sum_{j=0}^n
   \bigl(U\cdot\omega_\phi^{\,j}\cdot\omega_\psi^{\,n-j}\bigr).
\]
Each summand is the vertical-factor pairing of
\Cref{prop:ch20-intersection-pairing}; by
\eqref{eq:ch20-pairing-as-integral} this is the cocycle formula
\eqref{eq:energy-cocycle}, while \Cref{lem:ch26-telescope} supplies its
stepwise comparison.  Thus the two settings use the same telescoping
bookkeeping, but only the compactified test-configuration pairing in
\eqref{eq:ch46-telescope} is an absolute $(n+1)$-fold intersection.
\end{remark}

\section{The Monge--Amp\`ere measure of a PL metric}
\label{sec:ch46-pl-measure}

\subsection{The atomic formula}
\label{sec:ch46-atomic}

Fix psh PL data $(\theta_i,\phi_i)$, $1\le i\le n$, determined on an
integrally closed $\cX$, and consider the linear functional
\[
 \PL_{\mathbf R}(X)\longrightarrow\mathbf R,
 \qquad
 f\longmapsto(0,f)\prod_{i=1}^n(\theta_i,\phi_i).
\]
It is linear in $f$ by \Cref{prop:ch46-pairing-wellposed}.  The next
proposition identifies it with integration against an explicit finite
atomic measure, and is the model case of \Cref{thm:bj-mixed-ma}.

\begin{proposition}
\label{prop:ch46-pl-measure}
Let $\theta_1,\dots,\theta_n\in N^1(X)_{\mathbf R}$ and
$\phi_i=\phi_{D_i}\in\PL_{\mathbf R}(X)$ be determined on an integrally
closed test configuration $\cX$ dominating $X\times\mathbf A^1$, with
$\cL_i=\theta_{i,\cX}+D_i$.  Put
\[
 \mu:=\sum_{E\subset\cX_0}
      b_E\bigl(\cL_1|_E\cdots\cL_n|_E\bigr)\,\delta_{v_E},
\]
the sum over the irreducible components of the central fibre, with
$b_E$ and $v_E$ as in \Cref{def:ch33-component-valuation}.  Then
\begin{enumerate}
\item\label{itm:ch46-pl-ma:duality}
      $\displaystyle\int_{X^{\An}}f\,\rd\mu=(0,f)\prod_{i=1}^n(\theta_i,\phi_i)$
      for every $f\in\PL_{\mathbf R}(X)$;
\item\label{itm:ch46-pl-ma:indep}
      $\mu$ is independent of the model $\cX$ chosen, and is symmetric
      and multilinear in the $n$ pairs $(\theta_i,\phi_i)$;
\item\label{itm:ch46-pl-ma:mass}
      $\mu(X^{\An})=(\theta_1\cdots\theta_n)$;
\item\label{itm:ch46-pl-ma:pos}
      if every $\theta_i$ is nef and every $\phi_i$ is
      $\theta_i$-psh, then $\mu\ge0$; its support is a finite set of
      divisorial points, so $\mu$ gives no mass to $X^{\An}\setminus
      X^{\divis}$.
\end{enumerate}
\end{proposition}

\begin{proof}
\cref{itm:ch46-pl-ma:duality}.
Keep the model $\cX$ determining the measure slots fixed. Choose a
domination $\rho:\cX'\to\cX$ on which $f=\phi_F$ is determined.
The projection formula gives
\[
 (0,f)\prod_i(\theta_i,\phi_i)
  =(\rho_*[F]\cdot\cL_1\cdots\cL_n)_{\overline\cX}.
\]
The coefficient of $\rho_*[F]$ at a central component $E$ is
$b_Ef(v_E)$: this is the generic divisor order evaluated by the normalized
Gauss valuation of \Cref{def:ch33-component-valuation}.
Thus
\[
 \begin{aligned}
 \rho_*[F]&=\sum_Eb_Ef(v_E)[E],\\
 (0,f)\prod_i(\theta_i,\phi_i)
   &=\sum_Eb_Ef(v_E)(\cL_1|_E\cdots\cL_n|_E)\\
   &=\int f\,\rd\mu.
 \end{aligned}
\]
This is the fixed-model formula of \cite[Proposition~3.14(vi)]{BJ22}.
It allows the test function to live on a higher model without changing
the model defining the measure slots.

\cref{itm:ch46-pl-ma:indep}.
The functional just computed is intrinsic by
\Cref{prop:ch46-pairing-wellposed}. Since $\PL_{\mathbf R}(X)$ is
uniformly dense in $C^0(X^{\An})$ by \Cref{thm:ch34-pl-density}, two
finite signed Radon measures agreeing on every PL test agree on every
continuous test and hence coincide. Symmetry and multilinearity follow
from the same properties of the pairing.

\cref{itm:ch46-pl-ma:mass}.  Take $f=1$ in
\cref{itm:ch46-pl-ma:duality} and use
\eqref{eq:bj-energy-translation-pairing} with $\phi_0=0$, $c=1$, which
gives $(0,1)\prod(\theta_i,\phi_i)=(\theta_1\cdots\theta_n)$.
Equivalently, $\phi_1=\phi_{\cX_0}$ and $\cX_0=\overline\pi^*\{0\}$ is
a fibre class.

\cref{itm:ch46-pl-ma:pos}.  Each $E$ is projective over $k$ and lies in
a fibre of $\overline\pi$, and $\cL_i$ is relatively nef by
\eqref{eq:pl-relative-nefness}; hence $\cL_i|_E$ is nef and
$(\cL_1|_E\cdots\cL_n|_E)\ge0$ by \Cref{prop:appF-mixed-cartier}.  The
multiplicities $b_E$ are strictly positive integers by
\eqref{eq:ch33-bE}.  Finally $v_E\in X^{\divis}$ by
\Cref{lem:ch33-gauss-component}, and there are finitely many components.
\end{proof}

Combining \Cref{prop:ch46-pl-measure} with
\eqref{eq:bj-ma-normalizations}, the model formula announced in
\eqref{eq:bj-model-ma-atomic} is proved, and it is exactly
\eqref{eq:appJ-cl-measure} with the same pairing of the multiplicity
$b_E$ with the normalized point $v_E=b_E^{-1}\ord_E$.  The audit
demanded there --- ``using $\ord_E$ in place of $v_E$ without modifying
the coefficient moves the measure to a differently scaled valuation cone''
--- is a real one here, because the scaling action of
\Cref{sec:ch46-homogeneous} moves points along exactly those rays.

\begin{warning}{Two indexings of the central fibre}
\label{warn:ch46-two-indexings}
The components $E$ of $\cX_0$ index both the atoms of $\mu$ and the
vertices of the dual complex when $\cX$ is SNC.  The atom is placed at
$v_E=b_E^{-1}\ord_E|_{k(X)}$, not at $\ord_E$, and its mass carries the
compensating factor $b_E$.  Both conventions occur in the literature and
they differ by the scaling action; \Cref{warn:ch33-dh-normalizations} and
\Cref{rem:ch33-two-indexings} record the same issue for weights and flag
ideals.  In this chapter the pair $(b_E,v_E)$ is always used together.
\end{warning}

\subsection{Integration by parts and the Dirichlet form}
\label{sec:ch46-pl-ibp}

\begin{corollary}
\label{cor:ch46-pl-ibp}
For $f,g\in\PL_{\mathbf R}(X)$ and psh PL pairs
$(\theta_i,\phi_i)$, $1\le i\le n-1$,
\[
 \begin{aligned}
  \int_{X^{\An}}f\;\ddc g\wedge
    \bigwedge_{i=1}^{n-1}(\theta_i+\ddc\phi_i)
   &=(0,f)(0,g)\prod_{i=1}^{n-1}(\theta_i,\phi_i)\\
   &=\int_{X^{\An}}g\;\ddc f\wedge
     \bigwedge_{i=1}^{n-1}(\theta_i+\ddc\phi_i),
 \end{aligned}
\]
where $\ddc g$ denotes the pair $(0,g)$ in the appropriate slot.
\end{corollary}

\begin{proof}
Apply \Cref{prop:ch46-pl-measure}\,\cref{itm:ch46-pl-ma:duality} with the
$n$ pairs $(0,g),(\theta_1,\phi_1),\dots,(\theta_{n-1},\phi_{n-1})$ and
test function $f$, then with the roles of $f$ and $g$ exchanged.  The
middle term is the same in both cases because the pairing is symmetric
(\Cref{prop:ch46-pairing-wellposed}).
\end{proof}

The case $n=1$ of \Cref{cor:ch46-pl-ibp} deserves a separate statement,
because it is the exact counterpart of
\Cref{prop:ch26-dirichlet-vertical}, and because a curve is where every
sign in this chapter can be checked by hand.

\begin{proposition}
\label{prop:ch46-dirichlet-vertical}
\index{Dirichlet form!non-Archimedean}
Let $\cX$ be a test configuration dominating $X\times\mathbf A^1$ and
let $\cL_1,\dots,\cL_{n-1}$ be relatively nef.  For
$f=\phi_F,g=\phi_G\in\PL_{\mathbf R}(X)$ determined on $\cX$, set
\[
 \langle f,g\rangle_{\mathrm D}
 :=-\int_{X^{\An}}f\;\ddc g\wedge\cL_1\cdots\cL_{n-1}.
\]
Then:
\begin{enumerate}
\item\label{itm:ch46-dir:sym}
      $\langle f,g\rangle_{\mathrm D}
        =-\bigl(F\cdot G\cdot\cL_1\cdots\cL_{n-1}\bigr)_{\overline\cX}$
      is symmetric;
\item\label{itm:ch46-dir:pos}
      $\langle f,f\rangle_{\mathrm D}\ge0$, and
      $|\langle f,g\rangle_{\mathrm D}|^2
        \le\langle f,f\rangle_{\mathrm D}\langle g,g\rangle_{\mathrm D}$;
\item\label{itm:ch46-dir:curve}
      if $n=1$ no auxiliary classes occur and
      $\langle f,g\rangle_{\mathrm D}=-(F\cdot G)_{\overline\cX}$ is
      minus the intersection form on the divisors of $\overline\cX$
      supported in the fibre over $0$;
\item\label{itm:ch46-dir:energy}
      if $n=1$, $\deg\omega=d$ and $f$ is $\omega$-psh, then with
      $\mu_0=\MA^{\mathrm u}_\omega(0)=d\,\delta_{v_{\triv}}$,
      \[
       E^{\mathrm u}_\omega(f)
       =\int_{X^{\An}}f\,\rd\mu_0-\tfrac12\langle f,f\rangle_{\mathrm D},
       \qquad
       E_\omega(f)=\frac1d\Bigl(\int f\,\rd\mu_0
                    -\tfrac12\langle f,f\rangle_{\mathrm D}\Bigr).
      \]
\end{enumerate}
\end{proposition}

\begin{proof}
\cref{itm:ch46-dir:sym} is \Cref{cor:ch46-pl-ibp}, the middle term there
being $(0,f)(0,g)\prod(\theta_i,\phi_i)=(F\cdot G\cdot\cL_1\cdots\cL_{n-1})$
by \eqref{eq:bj-pl-energy-intersection}.
\cref{itm:ch46-dir:pos} is \Cref{lem:ch46-vertical-negativity}.
\cref{itm:ch46-dir:curve} is the case $n=1$, in which the product over
$i$ is empty.

\cref{itm:ch46-dir:energy}.  For $n=1$,
\eqref{eq:ch46-pure-pairing-as-integrals} reads
\[
 (\omega,f)^2=\int f\,\omega+\int f\,(\omega+\ddc f)
 =2\int f\,\rd\mu_0+\int f\,\ddc f,
\]
because $\omega=\MA^{\mathrm u}_\omega(0)=\mu_0$ by
\Cref{prop:ch46-pl-measure} applied to $\phi=0$, which puts the whole
mass $d$ at $v_{\triv}=v_{\widetilde X}$, the component of $\cX_0$
dominating $X$.  Dividing by $n+1=2$ and using
$\int f\,\ddc f=-\langle f,f\rangle_{\mathrm D}$ gives the first identity;
the second is \eqref{eq:bj-energy-definition} with $V=d$.
\end{proof}

\begin{remark}[Thuillier's form, and what changes over a trivial valuation]
\label{rem:ch46-thuillier}
On a curve, the graph formula identifies this form with the integral
of products of slopes. Its null space consists of constants: zero
Dirichlet norm forces every slope to vanish, and the valuation tree is
connected. This is the same curve form as
\Cref{prop:ch26-dirichlet-vertical}\,\cref{itm:ch26-dir:curve}, following
\cite[\S3.5]{Thu05}; see also \Cref{rem:ch26-thuillier-dirichlet}.
In both valuation regimes the energy is an affine functional minus half a
positive quadratic form. The affine term pairs the function with the
reference measure and is present in both regimes. The fact that the
special fibre is principal on a model over $R$ does not erase its degree
pairing supported on that fibre. On the compactified test configuration,
the corresponding fibre class is not principal
(\Cref{rem:ch46-fibre-class}). Thus energy on a curve cannot in general
be replaced by minus the Dirichlet norm alone.
\end{remark}

\subsection{First variation, and the normalization test}
\label{sec:ch46-first-variation}

\begin{proposition}
\label{prop:ch46-first-variation-pl}
Let $\phi,\psi\in\PL^{\omega}(X)$ and put
$\phi_t:=(1-t)\phi+t\psi$ for $t\in[0,1]$, which lies in
$\PL^{\omega}(X)$ because
\eqref{eq:pl-relative-nefness} defines a convex cone
(\Cref{prop:ch34-pl-operations}\,\cref{itm:ch34-plop:sum}).  Then
$t\mapsto E_\omega(\phi_t)$ is a polynomial of degree $\le n+1$, it is
concave, and
\[
 \frac{\rd}{\rd t}E_\omega(\phi_t)
 =\int_{X^{\An}}(\psi-\phi)\,\MA_\omega(\phi_t),
\]
with $\MA_\omega$ the \emph{probability} measure of
\eqref{eq:bj-ma-normalizations}.
\end{proposition}

\begin{proof}
Polynomiality is multilinearity: $(\omega,\phi_t)^{n+1}$ is a polynomial
of degree $\le n+1$ in $t$ because the pair $(\omega,\phi_t)$ is
affine in $t$.  Differentiating the expansion, and using symmetry,
\[
 \frac{\rd}{\rd t}(\omega,\phi_t)^{n+1}
 =(n+1)\,(0,\psi-\phi)(\omega,\phi_t)^{n}
 =(n+1)\int(\psi-\phi)\,\MA^{\mathrm u}_\omega(\phi_t),
\]
the last equality by
\Cref{prop:ch46-pl-measure}\,\cref{itm:ch46-pl-ma:duality}.  Dividing by
$(n+1)V$ gives the displayed formula, since
$\MA_\omega=V^{-1}\MA^{\mathrm u}_\omega$.

For concavity, differentiate once more.  Writing $h:=\psi-\phi$,
\[
 \begin{aligned}
 \frac{\rd^2}{\rd t^2}E_\omega(\phi_t)
 &=\frac{n}{V}\,(0,h)(0,h)(\omega,\phi_t)^{n-1}\\
 &=\frac nV\int h\,\ddc h
     \wedge(\omega+\ddc\phi_t)^{n-1}\\
 &=-\frac nV\langle h,h\rangle_{\mathrm D}\le0 .
 \end{aligned}
\]
by \Cref{prop:ch46-dirichlet-vertical}\,\cref{itm:ch46-dir:pos}, the
auxiliary classes being the relatively nef $\omega+\ddc\phi_t$.
\end{proof}

The displayed first variation is exactly
\eqref{eq:appJ-energy-variation} in the normalized convention, and it is
the second of the two tests announced in \Cref{warn:ch46-normalization}.
The two tests are not independent.  A candidate normalization is
$E=\lambda E^{\mathrm u}_\omega$ for some $\lambda>0$, and then
$E(\phi+c)=E(\phi)+\alpha c$ with $\alpha=\lambda V$ while
$\frac{\rd}{\rd t}E=\beta\int\dot\phi\,\MA^{\mathrm u}$ with
$\beta=\lambda$; so every admissible pair satisfies $\alpha=\beta V$, a
one-parameter family, and the unnormalized convention of \Cref{part:bfj} is
the member $(\alpha,\beta)=(V,1)$.  What pins the normalization down is
that relation \emph{together with} the translation rule
$E_\omega(\phi+c)=E_\omega(\phi)+c$, which selects $\alpha=1$ and hence
$\beta=V^{-1}$.

\begin{corollary}
\label{cor:ch46-pl-tangent}
For $\phi,\psi\in\PL^{\omega}(X)$:
\begin{enumerate}
\item\label{itm:ch46-tan:mono}
      $\phi\le\psi$ implies $E_\omega(\phi)\le E_\omega(\psi)$;
\item\label{itm:ch46-tan:concave}
      $E_\omega$ is concave on $\PL^{\omega}(X)$;
\item\label{itm:ch46-tan:tangent}
      $\displaystyle\int(\phi-\psi)\,\MA_\omega(\phi)
        \le E_\omega(\phi)-E_\omega(\psi)
        \le\int(\phi-\psi)\,\MA_\omega(\psi)$;
\item\label{itm:ch46-tan:cocycle}
      $\displaystyle
       E_\omega(\phi)-E_\omega(\psi)
       =\frac1{n+1}\sum_{j=0}^n\int(\phi-\psi)\,
        V^{-1}(\omega+\ddc\phi)^j\wedge(\omega+\ddc\psi)^{n-j}$.
\end{enumerate}
\end{corollary}

\begin{proof}
\cref{itm:ch46-tan:mono} is \Cref{prop:bj-model-energy-monotonicity} applied
in each of the $n+1$ slots in turn.
\cref{itm:ch46-tan:concave} is the second-derivative computation in the
proof of \Cref{prop:ch46-first-variation-pl}, applied to the segment
joining any two points.
\cref{itm:ch46-tan:cocycle} is \eqref{eq:ch46-telescope} divided by
$(n+1)V$, together with
\Cref{prop:ch46-pl-measure}\,\cref{itm:ch46-pl-ma:duality}.
\cref{itm:ch46-tan:tangent}: a concave differentiable function $h$ on
$[0,1]$ satisfies $h(1)-h(0)\le h'(0)$ and $h(1)-h(0)\ge h'(1)$;
apply this to $h(t)=E_\omega((1-t)\psi+t\phi)$ and use
\Cref{prop:ch46-first-variation-pl} at $t=0$ and $t=1$.
\end{proof}

\subsection{Energy is the first moment of a Duistermaat--Heckman measure}
\label{sec:ch46-dh}

The normalization \eqref{eq:bj-energy-definition} is not an arbitrary
choice: it is the one for which the energy of a PL psh metric is a moment
of a probability measure that was constructed, for entirely different
reasons, in \Cref{chap:bj-test-configurations}.

\begin{proposition}
\label{prop:ch46-energy-dh}
Let $L$ be an ample $\mathbf Q$-line bundle with
$\omega=c_1(L)$, let $(\cX,\cL)$ be an integrally closed test
configuration for $(X,L)$ with $\cL$ relatively ample, let
$\phi=\phi_{\cL}\in\cH_{\mathrm{dom}}(\omega)$ be its PL function, and
let $\mathcal F$ be the associated filtration of
\Cref{def:ch33-tc-filtration}.  Then
\[
 E_\omega(\phi)=\int_{\mathbf R}\lambda\,\rd\DHm(\mathcal F)(\lambda)
 =\lim_{m\to\infty}\frac{w_m}{m\,N_m},
\]
where $\DHm(\mathcal F)$ is the \emph{probability} Duistermaat--Heckman
measure of \Cref{thm:dh-limit-filtration}, and $w_m,N_m$ are the total
weight and the dimension of $H^0(\cX_0,m\cL_0)$.
\end{proposition}

\begin{proof}
By \Cref{def:bj-pl-energy-pairing} with all pairs equal to
$(\omega,\phi)$ and by \Cref{thm:ch33-compactification}\,(2),
$(\omega,\phi)^{n+1}=(\overline\cL^{\,n+1})$.  Hence, by
\eqref{eq:bj-energy-definition},
\[
 E_\omega(\phi)=\frac{(\overline\cL^{\,n+1})}{(n+1)V}.
\]
By \Cref{thm:ch33-dh-weight-count}\,(3) and the normalization
\eqref{eq:ch33-dh-unnormalized},
\[
 \int\lambda\,\rd\DHm^{\mathrm u}(\mathcal F)
 =\frac{(\overline\cL^{\,n+1})}{n+1},
 \qquad
 \DHm^{\mathrm u}(\mathcal F)=V\cdot\DHm(\mathcal F),
\]
so that $\int\lambda\,\rd\DHm(\mathcal F)
=(\overline\cL^{\,n+1})/((n+1)V)$.  The two right-hand sides agree.  The
last equality of the statement is
\Cref{thm:ch33-dh-weight-count}\,(1),(2).
\end{proof}

\begin{remark}[What the moment interpretation buys, and what it costs]
\label{rem:ch46-dh-scope}
Three consequences are worth naming.  First, the normalization test:
translating the rational linearization by $c\in\mathbf Q$ adds $c$ to the first moment and
$(n+1)cV$ to $(\overline\cL^{n+1})$, by
\eqref{eq:ch33-intersection-twist} and
\Cref{prop:ch33-dh-transformations}, and both changes are $+c$ in
$E_\omega$; this is \eqref{eq:bj-energy-basic-identities}.  Second,
degree-$d$ base change multiplies both by $d$, by
\eqref{eq:ch33-intersection-base-change}, which is the homogeneity proved
in \Cref{thm:ch46-energy-homogeneous} below.  Third, the moment formula
identifies $E_\omega(\phi)$ with an asymptotic weight average, which is
the form in which energy appears in stability theory
(\Cref{chap:bj-stability}).

The cost is that \Cref{prop:ch46-energy-dh} requires $\cL$ relatively
ample and $L$ rational and ample. The stated algebraic moment theorem
does not cover arbitrary real PL psh data, whereas
\eqref{eq:bj-energy-definition} does. The
proposition is therefore a bridge, not a definition, and
\Cref{warn:ch33-dh-normalizations} must be consulted before any moment is
quoted: the three normalizations there differ by the factors $V$ and
$n!$, and only the probability one gives $E_\omega$ on the nose.
\end{remark}

\section{A worked example on the projective line}
\label{sec:ch46-p1}

Before extending anything we compute one example completely, in three
independent ways, and check every normalization against it.  The example is
the deformation to the normal cone of a point, whose test configuration was
built in \Cref{ex:product-dnc-test-configurations}, whose dual complex was
computed in \Cref{ex:ch34-normal-cone-segment}, whose psh range was
determined in \Cref{ex:ch34-psh-check}, and whose Duistermaat--Heckman
measure was computed in \Cref{ex:ch33-dnc-p1}.

\begin{example}[Energy on $\mathbf P^1$, computed three ways]
\label{ex:ch46-p1-energy}
Let $X=\mathbf P^1_k$, $p\in X(k)$ a closed point with local parameter
$z$, $\theta=\omega=c_1(\cO(d))$ with $d\in\mathbf Z_{>0}$, so that
$n=1$ and $V=(\omega)=d$.  Let
\[
 \cY=\operatorname{Bl}_{(p,0)}\bigl(\mathbf P^1\times\mathbf A^1\bigr),
 \qquad
 \cY_0=\widetilde P+E,\quad b_{\widetilde P}=b_E=1,
\]
as in \Cref{ex:ch34-normal-cone-segment}, and let
$c\in\mathbf R$ with $0\le c\le d$.  Put $D:=-cE$ and
$\phi:=\phi_D$, so that by \eqref{eq:ch34-segment-function}
\[
 \phi(s\,\ord_q)=
 \begin{cases}
  -c\min\{s,1\}, & q=p,\\
  0, & q\ne p,
 \end{cases}
 \qquad
 \phi(v_{\triv})=0,\quad \phi(\ord_p)=-c .
\]
Equivalently $\phi=c\max\{\log|\mathfrak m_p|,-1\}$.  By
\eqref{eq:ch34-segment-psh-range}, $\phi$ is $\omega$-psh at PL level
exactly for $0\le c\le d$, which is the range fixed above.

\emph{First computation: intersection theory.}  The compactification is
$\overline\cY=\operatorname{Bl}_{(p,0)}(\mathbf P^1\times\mathbf P^1)$.
Write $H$ for the pullback of $\cO_{\mathbf P^1}(1)$ from the factor
$X$ and $F=\overline\pi^*\{0\}$ for the fibre class of the base.  Then
\[
 H^2=0,\quad F^2=0,\quad H\cdot F=1,\quad
 E^2=-1,\quad H\cdot E=F\cdot E=0,
\]
the last two because $E$ is contracted by both projections.  With
$\cL=\theta_{\overline\cY}+D=dH-cE$,
\[
 (\omega,\phi)^2=(\cL^2)=d^2H^2-2dc\,H\cdot E+c^2E^2=-c^2,
\]
so that by \eqref{eq:bj-energy-definition}
\[
 E_\omega(\phi)=\frac{(\omega,\phi)^2}{(n+1)V}=-\frac{c^2}{2d},
 \qquad
 E^{\mathrm u}_\omega(\phi)=-\frac{c^2}{2}.
\]

\emph{Second computation: from the measure.}  By
\Cref{prop:ch46-pl-measure}, with $\widetilde P=F-E$,
\[
 \MA^{\mathrm u}_\omega(\phi)
 =\bigl((dH-cE)\cdot\widetilde P\bigr)\delta_{v_{\triv}}
  +\bigl((dH-cE)\cdot E\bigr)\delta_{\ord_p}
 =(d-c)\,\delta_{v_{\triv}}+c\,\delta_{\ord_p},
\]
using $v_{\widetilde P}=v_{X,\triv}$ and $v_E=\ord_p$ from
\Cref{ex:ch34-normal-cone-segment}.  The total mass is $d=V$, confirming
\Cref{prop:ch46-pl-measure}\,\cref{itm:ch46-pl-ma:mass}, and the measure is
positive exactly for $0\le c\le d$, confirming
\cref{itm:ch46-pl-ma:pos} and matching
\eqref{eq:ch34-segment-psh-range}.  Also
$\MA^{\mathrm u}_\omega(0)=d\,\delta_{v_{\triv}}$.  Now
\eqref{eq:ch46-pure-pairing-as-integrals} with $n=1$ gives
\[
 (\omega,\phi)^2
 =\int\phi\,\MA^{\mathrm u}_\omega(0)+\int\phi\,\MA^{\mathrm u}_\omega(\phi)
 =d\cdot0+\bigl[(d-c)\cdot0+c\cdot(-c)\bigr]
 =-c^2 ,
\]
which is the number found above.

\emph{Third computation: Duistermaat--Heckman.}  For $d=1$ and
$0<c<1$ rational, \Cref{ex:ch33-dnc-p1} computes
$\DHm(\mathcal F)=(1-c)\delta_0+\mathbf 1_{[-c,0]}\rd\lambda$, whose
first moment is
$\int_0^1\min\{0,u-c\}\rd u=-c^2/2$.  By
\Cref{prop:ch46-energy-dh} this must equal $E_\omega(\phi)$, and indeed
$-c^2/(2\cdot1)=-c^2/2$. The intersection and measure computations
above apply also to real $c$; the moment formula extends to real $c$
by the explicit weak limit of these probability measures.

\emph{The unit tests.}  Adding a constant $e$ replaces $D$ by
$D+e\cY_0$, hence $\cL$ by $\cL+eF$, and
$(\cL+eF)^2=(\cL^2)+2e\,H\cdot F\cdot d=-c^2+2ed$; dividing by
$(n+1)V=2d$ gives $E_\omega(\phi)+e$, as
\eqref{eq:bj-energy-basic-identities} requires.  The Dirichlet form is
$\langle\phi,\phi\rangle_{\mathrm D}=-(D^2)=c^2\ge0$, as
\Cref{prop:ch46-dirichlet-vertical}\,\cref{itm:ch46-dir:pos} requires, and
\Cref{prop:ch46-dirichlet-vertical}\,\cref{itm:ch46-dir:energy} gives
$E^{\mathrm u}_\omega(\phi)=\int\phi\,\rd\mu_0-\tfrac12\langle\phi,\phi\rangle
=0-\tfrac12c^2$, agreeing again.  Finally the graph Laplacian: with the
sign convention of \Cref{chap:model-calculations}, in which $\ddc$ at a
point is the sum of the outgoing slopes, $\phi$ has slope $-c$ on
$[0,1]$ and slope $0$ beyond, so
$\ddc\phi=-c\,\delta_{v_{\triv}}+c\,\delta_{\ord_p}$ and
$\omega+\ddc\phi=(d-c)\delta_{v_{\triv}}+c\,\delta_{\ord_p}$, which is
\eqref{eq:bj-curve-laplacian}.
\end{example}

\begin{example}[The two-corner function]
\label{ex:ch46-p1-two-corner}
Keep $X=\mathbf P^1$ and take the data of \Cref{ex:ch34-two-corner}:
$\cY'=\operatorname{Bl}_\xi\cY$ with $\cY'_0=\widetilde P'+E'+2G$,
$D=cG+eE'$ with $c=-2$, $e=-\tfrac32$, and $\theta=c_1(\cO(2))$,
so $d=V=2$.  By \eqref{eq:ch34-two-corner-function} the vertex values are
$\phi_D(v_{\triv})=0$, $\phi_D(\tfrac12\ord_p)=-1$,
$\phi_D(\ord_p)=-\tfrac32$, and the slopes on
$[0,\tfrac12],[\tfrac12,1],[1,\infty)$ are $-2,-1,0$.

The intersection numbers on $\cY'$ recorded in
\Cref{ex:ch34-psh-check} give
\[
 \begin{aligned}
  (\theta_{\cY'}+D)\cdot\widetilde P'&=d+c=0,\\
  (\theta_{\cY'}+D)\cdot G&=-c+e=\tfrac12,\\
  (\theta_{\cY'}+D)\cdot E'&=c-2e=1,
 \end{aligned}
\]
so by \Cref{prop:ch46-pl-measure}, remembering $b_G=2$,
\[
 \MA^{\mathrm u}_\omega(\phi_D)
 =0\cdot\delta_{v_{\triv}}
  +2\cdot\tfrac12\,\delta_{\frac12\ord_p}
  +1\cdot\delta_{\ord_p}
 =\delta_{\frac12\ord_p}+\delta_{\ord_p},
\]
of total mass $2=V$.  The Laplacian check: the outgoing-slope sums are
$-2$ at $v_{\triv}$, $-1-(-2)=1$ at $\tfrac12\ord_p$, and
$0-(-1)=1$ at $\ord_p$; adding $\omega=2\delta_{v_{\triv}}$ gives the
same measure.  The energy is
\[
 E^{\mathrm u}_\omega(\phi_D)
 =\tfrac12\Bigl(\int\phi_D\,\omega+\int\phi_D\,\MA^{\mathrm u}(\phi_D)\Bigr)
 =\tfrac12\bigl(0+(-1)+(-\tfrac32)\bigr)=-\tfrac54,
\]
so that $E_\omega(\phi_D)=-\tfrac58$.
The multiplicity $b_G=2$ is essential: dropping it would give the atom
$\tfrac12$ at $\tfrac12\ord_p$, a total mass of $\tfrac32\ne V$, and
an energy of $-\tfrac12(1/2+3/2)=-1$.  This is the failure mode against
which \Cref{warn:ch46-two-indexings} warns.
\end{example}

\section{Extension by decreasing approximation}
\label{sec:ch46-extension}

For each $\omega_i\in\Amp(X)$, every
$\phi_i\in\PSH(\omega_i)$ is the limit of a decreasing net in
$\cH_{\mathrm{dom}}(\omega_i)$; a sequence is available only after the countable
regularization theorem.  This leads to the following definition-free
characterization.

\subsection{Three preparatory statements}
\label{sec:ch46-preparation}

The extension formula \eqref{eq:bj-energy-infimum} is an infimum over a
family of competitors, and three things have to be known about that family
before the formula can be used: that it is not empty, that its members come
arbitrarily close to $\phi$ where it matters, and that it interacts
correctly with decreasing limits.  We prove the three statements
separately, because they consume different inputs.

\begin{lemma}
\label{lem:ch46-competitors-nonempty}
Let $\omega\in\Amp(X)$ and $\phi\in\PSH(\omega)$.  Then
\[
 \mathcal F_{\mathrm{dom}}(\phi):=
 \{\psi\in\cH_{\mathrm{dom}}(\omega)\ :\ \psi\ge\phi\}
\]
contains every rational constant $c\ge\sup_{X^{\An}}\phi$, and in
particular is nonempty.  Moreover
$\mathcal F_{\mathrm{dom}}(\phi)+\mathbf Q_{\ge0}\subset\mathcal F_{\mathrm{dom}}(\phi)$.
\end{lemma}

\begin{proof}
By \eqref{eq:normalized-psh-space} the number $\sup_{X^{\An}}\phi$ is
finite, so rational constants $c$ above it exist.  A rational constant is
the PL function $\phi_{c\,\cX_0}$ on the trivial test configuration
$\cX=X\times\mathbf A^1$, by the computation $\phi_{\cX_0}=1$ carried
out in the proof of \Cref{prop:ch46-pairing-wellposed}; and
$\omega_{\cX}+c\,\cX_0$ is relatively ample over $\mathbf A^1$ because
$\omega$ is ample on $X$ and $\cX_0$ is a fibre, hence relatively
numerically trivial.  So $c\in\cH_{\mathrm{dom}}(\omega)$ by
\Cref{def:ch34-hdom}, and $c\ge\phi$.  The same computation shows that
adding a nonnegative rational constant to a member of
$\cH_{\mathrm{dom}}(\omega)$ keeps it there and only increases it.
\end{proof}

The next lemma is the technical heart of the extension.  It says that PL
psh competitors in the \emph{fixed} class $\omega$ can be found that
approximate $\phi$ at any prescribed finite set of divisorial points.
This is not a formal consequence of \Cref{def:bj-psh}, whose approximants
live in classes that merely converge to $\omega$; see
\Cref{rem:ch46-fixed-class-is-the-point}.

\begin{lemma}
\label{lem:ch46-pl-domination}
Let $\omega\in\Amp(X)$ be a real ample class, let
$u\in\PSH(\omega)$, let $S\subset X^{\divis}$ be finite, and let
$\varepsilon>0$.  Then there is
$\psi\in\cH_{\mathrm{dom}}(\omega)$ with
\[
 \psi\ge u\ \text{ on }X^{\An},
 \qquad
 \psi(v)\le u(v)+\varepsilon\ \text{ for every }v\in S .
\]
Consequently
$\inf\{\psi(v):\psi\in\mathcal F_{\mathrm{dom}}(u)\}=u(v)$ for every
$v\in X^{\divis}$.
\end{lemma}

\begin{proof}
By \Cref{thm:ch34-fs-regularization} there is a decreasing net
$(\psi_i)\subset\cH_{\mathrm{dom}}(\omega)$ with
$\psi_i\downarrow u$ pointwise on $X^{\An}$.  Thus every
$\psi_i\ge u$.  Since $S$ is finite, directedness supplies one index
$i$ for which
$\psi_i(v)\le u(v)+\varepsilon$ simultaneously for all $v\in S$.
Take $\psi=\psi_i$.

The last assertion of the lemma follows by letting
$S=\{v\}$ and $\varepsilon\downarrow0$, the reverse inequality being
the definition of $\mathcal F_{\mathrm{dom}}$.
\end{proof}

\begin{lemma}
\label{lem:ch46-dini}
Let $(f_j)_{j\in J}$ be a downward directed net of upper semicontinuous
functions $X^{\An}\to\mathbf R\cup\{-\infty\}$ with pointwise infimum
$f$, and let $g\in C^0(X^{\An})$ satisfy $f\le g$.  Then there is
$j\in J$ with $f_j\le g$ on all of $X^{\An}$, provided $f<g$
pointwise; in general, for every $\varepsilon>0$ there is $j$ with
$f_j\le g+\varepsilon$.
\end{lemma}

\begin{proof}
Fix $\varepsilon>0$ and set $U_j:=\{f_j<g+\varepsilon\}$.  Each
$U_j$ is open, because $g+\varepsilon-f_j$ is lower semicontinuous.
The family $(U_j)$ is upward directed, since $(f_j)$ is downward
directed.  It covers $X^{\An}$: for $x\in X^{\An}$,
$\inf_jf_j(x)=f(x)\le g(x)<g(x)+\varepsilon$, so some $j$ has
$f_j(x)<g(x)+\varepsilon$.  The space $X^{\An}$ is compact by
\Cref{thm:triv-berkovich-compactness}, so finitely many $U_j$ cover it,
and by directedness a single one does.  The variant with $f<g$ and no
$\varepsilon$ is the same argument with $U_j=\{f_j<g\}$.
\end{proof}

\subsection{The extension theorem}
\label{sec:ch46-extension-theorem}

\begin{theorem}\label{thm:bj-extended-energy-pairing}
For ample $\omega_0,\dots,\omega_n$, there is a unique map

\[
 \prod_{i=0}^n\PSH(\omega_i)\longrightarrow\mathbf R\cup\{-\infty\},
 \quad
 (\phi_i)\longmapsto(\omega_0,\phi_0)\cdots(\omega_n,\phi_n),
\tag{46.11}\label{eq:bj-extended-energy-map}
\]

which extends the PL pairing, is increasing in every variable, upper
semicontinuous for pointwise convergence on $X^{\divis}$, and continuous
along decreasing nets.  Explicitly,

\[
 (\omega_0,\phi_0)\cdots(\omega_n,\phi_n)
 =\inf_{\substack{\psi_i\in\cH_{\mathrm{dom}}(\omega_i)\\
                   \psi_i\geq\phi_i}}
      (\omega_0,\psi_0)\cdots(\omega_n,\psi_n).
\tag{46.12}\label{eq:bj-energy-infimum}
\]
\end{theorem}

\begin{proof}[Proof of the foundational extension]
Write $I(\phi_0,\dots,\phi_n)$ for the right-hand side of
\eqref{eq:bj-energy-infimum} and $P$ for the PL pairing of
\Cref{def:bj-pl-energy-pairing}.  The classes $\omega_i$ may be arbitrary
real ample classes: \Cref{lem:ch46-pl-domination} has exactly that scope,
by the fixed-class regularization of
\Cref{thm:ch34-fs-regularization}.

\emph{Step 1: $I$ is well defined and increasing.}  The competitor family
is nonempty by \Cref{lem:ch46-competitors-nonempty}, so the infimum is
taken over a nonempty set and lies in $\mathbf R\cup\{-\infty\}$.  If
$\phi_0\le\phi_0'$, every competitor for $\phi_0'$ is one for
$\phi_0$, so the infimum for $\phi_0$ is over a larger set and is
therefore smaller.  This is monotonicity in each variable, and it uses no
property of $P$ at all.

\emph{Step 2: $I$ extends the PL pairing.}  Let
$\phi_i\in\PL_{\mathbf R}(X)\cap\PSH(\omega_i)$.  If
$\psi_i\in\cH_{\mathrm{dom}}(\omega_i)$ satisfies $\psi_i\ge\phi_i$,
then $P(\psi)\ge P(\phi)$ by
\Cref{prop:bj-model-energy-monotonicity} applied in each slot in turn; so
$I(\phi)\ge P(\phi)$. Conversely, fixed-class regularization gives
a decreasing net in $\cH_{\mathrm{dom}}(\omega_i)$ above each
$\phi_i$. Since $\phi_i$ is continuous, \Cref{lem:ch46-dini} gives,
for every $\varepsilon>0$, a member $\psi_i$ with
$\phi_i\le\psi_i\le\phi_i+\varepsilon$ on all of $X^{\An}$. By
\Cref{cor:ch46-pl-lipschitz}, applied in each of the $n+1$ slots,
$P(\psi)\le P(\phi)+\varepsilon\sum_i(\omega_0\cdots\widehat{\omega_i}
\cdots\omega_n)$.  Letting $\varepsilon\downarrow0$ gives
$I(\phi)\le P(\phi)$. The same approximation shows that the infimum
is unchanged if the competitors are allowed to be all real PL psh
majorants, rather than just members of
$\cH_{\mathrm{dom}}(\omega_i)$. This enlarged family is useful for
real translations and simultaneous scaling of classes and potentials.

\emph{Step 3: finiteness on bounded tuples.}  If $m_i\le\phi_i\le M_i$
with $m_i,M_i\in\mathbf R$, monotonicity and Step~2 give
$I(\phi)\le I(M_0,\dots,M_n)=P(M_0,\dots,M_n)<\infty$ and
$I(\phi)\ge I(m_0,\dots,m_n)=P(m_0,\dots,m_n)>-\infty$. Thus
$I$ is real valued on tuples of bounded psh functions.  For unbounded
$\phi_i$ the value $-\infty$ does occur, and
\Cref{sec:ch46-e1} is devoted to the locus where it does not.

\emph{Step 4: continuity along decreasing nets.}  Let
$(\phi_{i,j})_{j\in J}$ be downward directed nets in $\PSH(\omega_i)$
with $\phi_{i,j}\downarrow\phi_i$, indexed by a common directed set.  By
Step~1 the net $\bigl(I(\phi_{0,j},\dots,\phi_{n,j})\bigr)_j$ is
decreasing and bounded below by $I(\phi_0,\dots,\phi_n)$, so it has a
limit $\ell\ge I(\phi)$.  For the reverse inequality, fix
$\varepsilon>0$ and competitors $\psi_i\in\cH_{\mathrm{dom}}(\omega_i)$
with $\psi_i\ge\phi_i$ and $P(\psi)\le I(\phi)+\varepsilon$, which
exist by the definition of the infimum unless $I(\phi)=-\infty$, in which
case one takes $P(\psi)\le-1/\varepsilon$.  Each $\psi_i$ is
continuous, and $\phi_i\le\psi_i$; by \Cref{lem:ch46-dini} there is, for
each $i$, an index beyond which $\phi_{i,j}\le\psi_i+\varepsilon$, and
by directedness of $J$ a single index $j_0$ works for all $i$.  For
$j\ge j_0$ the functions $\psi_i+\varepsilon'$, with
$\varepsilon'\in\mathbf Q\cap[\varepsilon,2\varepsilon]$, are competitors
for $\phi_{i,j}$, whence by
\eqref{eq:bj-energy-translation-pairing}
\[
 I(\phi_{0,j},\dots,\phi_{n,j})
 \le P(\psi_0+\varepsilon',\dots,\psi_n+\varepsilon')
 =P(\psi)+\varepsilon'\sum_{i=0}^n
   (\omega_0\cdots\widehat{\omega_i}\cdots\omega_n).
\]
Hence $\ell\le I(\phi)+\varepsilon+2\varepsilon C$ with $C$ the
displayed sum, or $\ell\le-1/\varepsilon+2\varepsilon C$ in the infinite
case.  Letting $\varepsilon\downarrow0$ gives $\ell\le I(\phi)$,
respectively $\ell=-\infty$.

\emph{Step 5: joint upper semicontinuity.}
Fix a tuple $\phi$ and $b\in\mathbf R$ with $I(\phi)<b$.
Choose competitors $\psi_i\in\cH_{\mathrm{dom}}(\omega_i)$ and
$\delta>0$ such that $P(\psi)<b-\delta$. Put
\[
 C:=\sum_{i=0}^n(\omega_0\cdots\widehat{\omega_i}\cdots\omega_n)>0.
\]
For each PL function $\psi_i$, \cite[Lemma~4.26]{BJ22} supplies a finite
set $\Sigma_i\subset X^{\divis}$ such that
\[
 \sup_{X^{\An}}(u-\psi_i)=\max_{v\in\Sigma_i}(u(v)-\psi_i(v))
 \qquad\text{for every }u\in\PSH(\omega_i).
\]
This is a continuous function of $u$ for the divisorial pointwise topology.
Since $\phi_i\le\psi_i$, there is an open neighbourhood $U_i$ of
$\phi_i$ on which that supremum is less than $\delta/C$.
For $\phi'_i\in U_i$, monotonicity, the PL compatibility of Step~2, and
translation of the pairing give
\[
 I(\phi'_0,\ldots,\phi'_n)
 \le P(\psi_0+\delta/C,\ldots,\psi_n+\delta/C)
 =P(\psi)+\delta<b.
\]
Thus $\prod_iU_i$ lies in $\{I<b\}$, proving joint upper
semicontinuity. The finite-set comparison is the precise topological input;
separate upper semicontinuity alone would not prove this conclusion.

\emph{Step 6: uniqueness.}
For each $\phi_i$, choose a decreasing net in
$\cH_{\mathrm{dom}}(\omega_i)$ converging to $\phi_i$, and index their
finite product by one directed set. Any competing extension $G$ that
agrees with $P$ and is continuous along decreasing nets satisfies
\[
 G(\phi)=\lim_jP(\psi_{0,j},\ldots,\psi_{n,j})=I(\phi),
\]
the last equality by Step~4. Hence $G=I$.

This is the proof of \cite[Theorem~7.1]{BJ22}, with nets
retained throughout.
\end{proof}

\begin{remark}[Why the fixed class is the whole point]
\label{rem:ch46-fixed-class-is-the-point}
\Cref{def:bj-psh}\,(3) approximates a $\theta$-psh function by
Fubini--Study functions in classes $c_1(L_\iota)$ that merely
\emph{converge} to $\theta$, and \Cref{rem:ch45-real-again} explained why
the varying-class convention cannot be dispensed with in general.  For the
present chapter that convention is useless: the pairing
\eqref{eq:bj-pl-energy-intersection} depends on the classes as well as on
the potentials, so a varying-class approximation would produce a family of
numbers $\bigl(c_1(L_{0,\iota}),\phi_{0,\iota}\bigr)\cdots$ whose limit
would have to be shown independent of the approximating classes before it
could be called the pairing of $(\omega_i,\phi_i)$.  What
\Cref{thm:ch34-fs-regularization} supplies, and what the bare definition
does not, is a monotone PL approximation in the fixed class $\omega$
itself.  This holds for every real ample class, whether or not it is the
class of a $\mathbf Q$-line bundle, and requires no envelope property.
\Cref{thm:ch45-continuous-regularization} supplies a different,
obstacle-indexed net that is useful for envelope calculus but is not needed
to define the extended pairing.
\end{remark}

\begin{remark}[Real ample classes]
\label{rem:ch46-real-classes}
The set $\cH_{\mathrm{dom}}(\omega)$ and the regularization theorem are
defined for every real $\omega\in\Amp(X)$; rationality is required only
of the vertical divisor determining each PL approximant, not of $\omega$
itself.  Consequently \Cref{lem:ch46-pl-domination} and the proof of
\Cref{thm:bj-extended-energy-pairing} apply verbatim to arbitrary real
ample classes, as in \cite[Theorem~7.1]{BJ22}.  The same observation gives
the real-class form of \Cref{prop:ch46-energy-usc}, whose proof of
\cref{itm:ch46-Eusc:rep} uses only
\Cref{lem:ch46-pl-domination}.  The larger obstacle family
$\{w\in C^0\cap\PSH(\omega):w\ge u\}$ from
\Cref{lem:ch46-continuous-competitors} remains useful for envelope
arguments, but it is not needed to bridge rational and real ample classes.
\end{remark}

\begin{lemma}
\label{lem:ch46-continuous-competitors}
Let $\omega\in\Amp(X)$ be a real class and let $u\in\PSH(\omega)$.
Then
\[
 \mathcal D_u:=\{w\in C^0(X^{\An})\cap\PSH(\omega)\ :\ w\ge u\}
\]
is nonempty, downward directed, and satisfies $\inf\mathcal D_u=u$ pointwise on
$X^{\An}$.
\end{lemma}

\begin{proof}
Nonemptiness: the constant $\sup_{X^{\An}}u$ lies in $\mathcal D_u$, being
$\omega$-psh by \Cref{lem:ch34-psh-nonempty}.

Directedness: let $w_1,w_2\in\mathcal D_u$ and put
$f:=\min\{w_1,w_2\}\in C^0(X^{\An})$.  Then $f$ is an admissible
obstacle in the sense of \Cref{def:ch45-admissible-obstacle}: it is
continuous, hence bounded and upper semicontinuous, and its competitor
family contains $u$.  By \Cref{thm:bj-final-obstacle-psh} the function
$P_\omega(f)$ is the greatest $\omega$-psh function $\le f$, so
$u\le P_\omega(f)\le f\le w_1,w_2$; and $P_\omega(f)$ is continuous by
\Cref{cor:ch45-continuity-envelopes}, $\omega$ being ample.  Hence
$P_\omega(f)\in\mathcal D_u$ lies below both.

Infimum: by \Cref{thm:ch45-continuous-regularization} there is a downward
directed net $(u_g)\subset C^0(X^{\An})\cap\PSH(\omega)$ with
$u_g\downarrow u$; each $u_g\ge u$, so $u_g\in\mathcal D_u$ and
$\inf\mathcal D_u\le\inf_gu_g=u$.  The reverse inequality is the definition of
$\mathcal D_u$.
\end{proof}

\begin{remark}[The topological input in Step~5]
\label{rem:ch46-usc-scope}
The finite-set formula of \cite[Lemma~4.26]{BJ22} makes
$u\mapsto\sup(u-\psi)$ continuous for every fixed PL test function
$\psi$. Uniform approximation of a continuous $\psi$ by PL functions
extends this to continuous tests, as in \cite[Proposition~4.28]{BJ22}.
This controls a varying psh function against a fixed continuous majorant
on the whole analytification, even though the topology records only
its divisorial values. It supplies joint upper semicontinuity in Step~5.

The mixed pairing is not jointly concave. Already for $n=1$, its second
variation along a tuple of two affine paths is the mixed vertical
intersection $(0,\chi_0-\psi_0)(0,\chi_1-\psi_1)$, which can have either
sign. Concavity of the pure energy therefore gives a different argument
for its upper semicontinuity in \Cref{prop:ch46-energy-usc}.
\end{remark}

\begin{definition}[Monge--Amp\`ere energy]
\label{def:bj-ma-energy}
For $\phi\in\PSH(\omega)$, define

\[
                   E_\omega(\phi)
       :=\frac{(\omega,\phi)^{n+1}}{(n+1)V}
       \in\mathbf R\cup\{-\infty\}.
\tag{46.13}\label{eq:bj-energy-definition}
\]
\end{definition}

The normalization has three immediate consequences:

\[
 E_\omega(0)=0,\qquad E_\omega(\phi+c)=E_\omega(\phi)+c,
 \qquad E_{t\omega}(t\phi)=tE_\omega(\phi).
\tag{46.14}\label{eq:bj-energy-basic-identities}
\]

The functional is increasing, concave, upper semicontinuous, and continuous
along decreasing nets; these are \cite[Proposition~7.7]{BJ22}.  Concavity
is the Hodge-index input, first proved for PL functions and then passed to
decreasing limits.

We prove these four properties, since the argument is short and since the
representation formula it produces is used again in
\Cref{chap:bj-finite-energy}.

\begin{proposition}
\label{prop:ch46-energy-usc}
Let $\omega\in\Amp(X)$ be a real ample class.  Then, on
$\PSH(\omega)$,
\begin{enumerate}
\item\label{itm:ch46-Eusc:rep}
      $\displaystyle
       E_\omega(\phi)
       =\inf_{\psi\in\cH_{\mathrm{dom}}(\omega)}
        \Bigl[E_\omega(\psi)
              +\int_{X^{\An}}(\phi-\psi)\,\MA_\omega(\psi)\Bigr];$
\item\label{itm:ch46-Eusc:props}
      $E_\omega$ is increasing, concave, and upper semicontinuous for
      pointwise convergence on $X^{\divis}$;
\item\label{itm:ch46-Eusc:decreasing}
      $E_\omega$ is continuous along decreasing nets;
\item\label{itm:ch46-Eusc:identities}
      \eqref{eq:bj-energy-basic-identities} holds, and
      $E_\omega(\phi)\le\sup_{X^{\An}}\phi$, with equality when
      $\phi$ is constant;
\item\label{itm:ch46-Eusc:strict}
      if $\phi\in\cH_{\mathrm{dom}}(\omega)$ is nonconstant, then
      $E_\omega(\phi)<\sup_{X^{\An}}\phi$.
\end{enumerate}
\end{proposition}

\begin{proof}
\cref{itm:ch46-Eusc:rep}.  Write $T_\psi(\phi)$ for the bracket.  For
$\psi\in\cH_{\mathrm{dom}}(\omega)$ and a PL competitor
$\chi\ge\phi$, $\chi\in\cH_{\mathrm{dom}}(\omega)$, the tangent
inequality \Cref{cor:ch46-pl-tangent}\,\cref{itm:ch46-tan:tangent} gives
$E_\omega(\chi)\le T_\psi(\chi)$.  The measure
$\MA_\omega(\psi)$ is carried by a finite set
$S\subset X^{\divis}$, so
\Cref{lem:ch46-pl-domination} gives
$\inf_\chi\int\chi\,\MA_\omega(\psi)=\int\phi\,\MA_\omega(\psi)$, the
infimum over competitors $\chi$ for $\phi$.  Taking that infimum in the
displayed inequality and using \eqref{eq:bj-energy-infimum} on the left
yields $E_\omega(\phi)\le T_\psi(\phi)$, for every $\psi$.  Conversely,
for a competitor $\chi\ge\phi$ one has
$T_\chi(\phi)=E_\omega(\chi)+\int(\phi-\chi)\MA_\omega(\chi)
\le E_\omega(\chi)$, because $\phi\le\chi$ and $\MA_\omega(\chi)\ge0$;
taking the infimum over $\chi$ gives
$\inf_\psi T_\psi(\phi)\le E_\omega(\phi)$.

\cref{itm:ch46-Eusc:props}.  Each $T_\psi$ is an affine function of
$\phi$, and it is continuous for pointwise convergence on
$X^{\divis}$, being an integral against a measure carried by a finite
subset of $X^{\divis}$.  An infimum of affine continuous functions is
concave and upper semicontinuous.  Monotonicity is Step~1 of the proof of
\Cref{thm:bj-extended-energy-pairing}, or again the fact that each
$T_\psi$ is increasing.

\cref{itm:ch46-Eusc:decreasing} is Step~4 of the same proof, specialized to
equal arguments.

\cref{itm:ch46-Eusc:identities}.  The three identities
\eqref{eq:bj-energy-basic-identities} hold on PL data by
\eqref{eq:bj-trivial-energy-zero},
\eqref{eq:bj-energy-translation-pairing} and multilinearity, and pass to
the infimum \eqref{eq:bj-energy-infimum} because the operations involved
--- adding a constant, multiplying a class and a potential by $t>0$ ---
are bijections of the real PL competitor families justified in Step~2
of the extension proof. For the inequality, monotonicity gives
$E_\omega(\phi)\le E_\omega(\sup\phi)=\sup\phi$;
constants have their own value by PL compatibility and translation.

\cref{itm:ch46-Eusc:strict}.
Normalize $u:=\phi-\sup\phi\le0$, with $u(v_{\triv})=0$.
Choose a common integrally closed test configuration determining $u$
such that $\cL=\omega_{\cX}+D_u$ is relatively ample. Every coefficient
of $D_u$ is nonpositive by the effectivity criterion, and at least one
is negative because $u$ is nonconstant. The atomic formula gives
\[
 \int u\,\MA_\omega(u)
 =V^{-1}\sum_E b_Eu(v_E)(\cL|_E)^n<0:
\]
every $(\cL|_E)^n$ is strictly positive by ampleness.
In the PL cocycle comparing $u$ with $0$, each summand is nonpositive,
since $u\le0$ and the mixed measures are positive. The summand with all
measure slots equal to $u$ is strictly negative by the display.
Thus $E_\omega(u)<0$, and translation gives
$E_\omega(\phi)<\sup\phi$. This proof uses relative ampleness and the
atomic formula, with no equality case of a Hodge index theorem on a
possibly singular model.
\end{proof}

\begin{remark}[Equality for unbounded potentials]
\label{rem:ch46-equality-scope}
\Cref{prop:ch46-energy-usc}\,\cref{itm:ch46-Eusc:strict} is stated for
$\cH_{\mathrm{dom}}(\omega)$ because its proof uses strictly positive
component weights on a relatively ample model.  The corresponding statement
for an arbitrary $\phi\in\PSH(\omega)$ --- that
$E_\omega(\phi)=\sup\phi$ forces $\phi$ constant --- is true, and
follows from the domination principle of
\Cref{chap:bj-finite-energy}; it is not proved here, and no statement of
this chapter uses it.  What the cocycle does give, and what is used, is the
weaker assertion of \Cref{prop:ch46-cocycle}\,\cref{itm:ch46-coc:mono}:
equality in a monotonicity comparison forces equality almost everywhere for
the relevant Monge--Amp\`ere measure.
\end{remark}

\begin{warning}{No undefined subtraction}
\label{warn:ch46-no-subtraction}
The mixed pairing of arbitrary psh functions may be $-\infty$.  We never
polarize an identity by subtracting two infinite values.  Multilinear
extensions and cocycle formulae below are asserted only on the finite-energy
locus, where every mixed term is finite.

Concretely, three formulae in this chapter would be meaningless without the
restriction.  The cocycle \eqref{eq:bj-energy-cocycle} contains the
difference $\phi-\psi$, which is undefined at a proper-support point
where both are $-\infty$; it is asserted for
$\phi,\psi\in\cE^1(\omega)$, and even there the integrand is defined
almost everywhere using the integrability in \Cref{thm:bj-mixed-ma}.  The multilinear extension
\eqref{eq:bj-e1-polarized-pairing} is asserted on the vector space
$\vec\cE^{\,1}(X)$ and not on $\PSH$, which is not a vector space.  And
the integration-by-parts formula \eqref{eq:bj-integration-by-parts} is
asserted for $f,g\in\vec\cE^{\,1}(X)$, for the same reason.  The general
principle is the one stated in the front matter: a displayed identity
between extended-real quantities is asserted only where both sides are
defined, and here that locus has been isolated rather than left implicit.
\end{warning}

\subsection{Where the envelope theorem is used}
\label{sec:ch46-where-envelope}

It is worth listing, once, the places in this chapter where the results of
\Cref{chap:bj-envelope-consequences} are genuinely consumed, so that the
reader can see which parts of the theory are older than the envelope
theorem and which are not.

\begin{enumerate}
\item\label{itm:ch46-env:directed}
      \emph{Directedness of the continuous competitor family.}
      \Cref{lem:ch46-continuous-competitors} uses
      \Cref{thm:bj-final-obstacle-psh} and
      \Cref{cor:ch45-continuity-envelopes} for the lower bound of two
      continuous competitors. It packages continuous majorants as an
      obstacle-indexed decreasing net.  The foundational pairing and the
      mixed-measure construction do not require this package: their
      monotone PL approximants come directly from
      \Cref{thm:ch34-fs-regularization}.
\item\label{itm:ch46-env:real}
      \emph{Real ample classes do not consume the envelope theorem.}
      \Cref{thm:ch34-fs-regularization} regularizes in
      $\cH_{\mathrm{dom}}(\omega)$ for every real ample $\omega$.
      This is the input used by \Cref{lem:ch46-pl-domination} and the
      extension theorem; see \Cref{rem:ch46-real-classes}.
\item\label{itm:ch46-env:compactness}
      \emph{Compactness.}  \Cref{cor:ch45-compactness} together with the
      upper semicontinuity of \Cref{prop:ch46-energy-usc} is what makes
      $E_\omega$ attain its supremum on nonempty weakly closed subsets of
      $\PSH^{\sup}(\omega)$.  Neither statement alone suffices, and the
      compactness half is exactly the envelope theorem
      (\Cref{rem:ch45-compactness-equivalent}).  This is the input to the
      variational method of \Cref{chap:bj-finite-energy}.
\item\label{itm:ch46-env:obstacles}
      \emph{Extremal functions and obstacle calculus.} The operator
      $P_\theta$ and the extremal function $V_\theta=P_\theta(0)$
      are psh under the hypotheses of Chapter 45.
      Locality and comparison in \Cref{sec:ch46-ibp} instead use finite
      maxima and the published mixed-measure theory; they consume no
      envelope theorem.
\end{enumerate}

\begin{remark}[What survives without the envelope theorem]
\label{rem:ch46-without-envelope}
Everything in \Cref{sec:ch46-model,sec:ch46-pl-measure,sec:ch46-p1}
survives, for an arbitrary integral projective $X$ over the standing
algebraically closed trivially valued field: those sections are intersection theory on a
projective $(n+1)$-fold.  So do the extension
\Cref{thm:bj-extended-energy-pairing} for arbitrary real ample classes and
the mixed measures and finite-energy formalism built from it: their
fixed-class PL approximation is \Cref{thm:ch34-fs-regularization}, which
does not require the envelope property.

What does not survive is the compact variational and envelope package:
without the envelope property, the compactness and psh-valued envelope
inputs used by the variational argument are unavailable. This does not
exclude particular maximizers or solutions proved by another method.
The raw extremal supremum remains defined, but the envelope theorem is
what guarantees its psh membership when competitors exist.
The continuous obstacle-indexed regularization invoked here is an
ample-class result; no such approximation outside the ample regime is
asserted. This is the
sense in which \Cref{chap:bj-envelope-consequences} is the hinge of
\Cref{part:bj}: the finitary theory and its monotone ample-class extension
were already available, while compactness, extremal envelopes, and the
variational existence theory had to wait.
\end{remark}

\section{Finite energy and polarization}
\label{sec:ch46-e1}

The extended pairing takes the value $-\infty$, and the mechanism is
already visible on a curve: a psh function whose graph on a divisorial ray
decays too fast has infinite Dirichlet energy.  The locus on which this does
not happen is a convex class that behaves like a vector space, and it is the
domain on which every multilinear identity of this chapter is asserted.

\begin{definition}[Potential-side finite energy]
\label{def:bj-e1}
Set

\[
 \cE^1(\omega)
 :=\{\phi\in\PSH(\omega)\mid E_\omega(\phi)>-\infty\}.
\tag{46.15}\label{eq:bj-e1-definition}
\]
\end{definition}

This is a definition about pointwise psh functions.  It is not the
measure-side synthetic completion discussed in
\Cref{chap:bfj-finite-energy-measures}.

\begin{theorem}\label{thm:bj-e1-basic}
The class $\cE^1(\omega)$ is convex, contains every bounded
$\omega$-psh function, is stable under constants, and is upward closed
inside $\PSH(\omega)$.  If $\omega\leq\omega'$, then
$\cE^1(\omega)\subset\cE^1(\omega')$, and

\[
 \cE^1(\omega)+\cE^1(\omega')
       \subset\cE^1(\omega+\omega').
\tag{46.16}\label{eq:bj-e1-addition}
\]

Every mixed pairing of $n+1$ psh finite-energy pairs is finite.
\end{theorem}

\begin{proof}[Proof architecture]
Convexity and upward closure follow from concavity and monotonicity of
energy.  For the change of class, normalize the potentials to be
nonpositive and compare ample classes by
$\omega_i\leq t\omega_j$.  Expanding the energy of sums and applying
the Hodge-type inequalities gives a uniform lower bound for every mixed
term in terms of the pure energies.  This is
\cite[Theorem~7.8, Lemmas~7.9--7.10, and Corollary~7.11]{BJ22}.
The estimates, rather than formal polarization with infinities, prove the
last assertion.

We give the first three assertions in full, since they are immediate from
\Cref{prop:ch46-energy-usc} and are used constantly, and indicate precisely
what the quoted estimates supply.

\emph{Bounded functions.}  If $\phi\in\PSH(\omega)$ is bounded, then
$m\le\phi\le M$ for real $m,M$, and monotonicity together with
$E_\omega(m)=m$ gives $m\le E_\omega(\phi)\le M$; in particular
$E_\omega(\phi)$ is finite.  This is Step~3 of the proof of
\Cref{thm:bj-extended-energy-pairing}, specialized to equal arguments.

\emph{Constants and upward closure.}
$E_\omega(\phi+c)=E_\omega(\phi)+c$ by
\eqref{eq:bj-energy-basic-identities}, so $\cE^1(\omega)+\mathbf R
=\cE^1(\omega)$; and if $\phi\le\psi$ with
$\phi\in\cE^1(\omega)$ and $\psi\in\PSH(\omega)$, then
$E_\omega(\psi)\ge E_\omega(\phi)>-\infty$, so
$\psi\in\cE^1(\omega)$.

\emph{Convexity.}  $\PSH(\omega)$ is convex by
\Cref{thm:bj-basic-closure}\,(3), and $E_\omega$ is concave by
\Cref{prop:ch46-energy-usc}\,\cref{itm:ch46-Eusc:props}; a concave function
is $>-\infty$ on the segment joining two points at which it is
$>-\infty$.

\emph{Change of class and addition.}  For \eqref{eq:bj-e1-addition} and the
inclusion $\cE^1(\omega)\subset\cE^1(\omega')$ one compares energies in
different classes, and this is where the Hodge-type estimates enter.  The
mechanism is the following elementary case, which we record because it
already shows the shape of the argument.  If $\omega\le\omega'$ are both
ample and $\phi\le0$ is $\omega$-psh, then $\phi$ is
$\omega'$-psh by \Cref{thm:bj-basic-closure}\,(3) applied to
$0\in\PSH(\omega'-\omega)$, and, with
$V=(\omega^n)$, $V'=(\omega'^n)$,
\[
 (n+1)V'E_{\omega'}(\phi)=(\omega',\phi)^{n+1}
 \ \le\ (\omega,\phi)^{n+1}=(n+1)VE_\omega(\phi),
\]
because expanding $(\omega',\phi)^{n+1}$ in
$(\omega,\phi)+(\omega'-\omega,0)$ produces, besides
$(\omega,\phi)^{n+1}$, only terms of the form
$(\omega'-\omega,0)^{k}(\omega,\phi)^{n+1-k}$ with $k\ge1$, each of
which is $\le0$ for $\phi\le0$: such a term is either $0$, or the
integral of the nonpositive $\phi$ against a positive mixed measure, by
\Cref{prop:ch46-pl-measure}\,\cref{itm:ch46-pl-ma:pos} at PL level and by
monotone approximation in general.  The direction is easy to check on
$X=\mathbf P^1$ with $\omega=c_1(\cO(1))$, $\omega'=c_1(\cO(2))$ and
$\phi\equiv-1$: there $E_\omega(\phi)=E_{\omega'}(\phi)=-1$ by
\eqref{eq:bj-energy-basic-identities}, so $(\omega,\phi)^2=-2$ while
$(\omega',\phi)^2=-4$.  Note that this inequality runs the \emph{wrong}
way to give the inclusion $\cE^1(\omega)\subseteq\cE^1(\omega')$, which
is why that inclusion is not proved this way.  The general statements
\eqref{eq:bj-e1-addition} and the finiteness of every mixed term require
the two-sided estimates of \cite[Theorem~7.8, Lemmas~7.9--7.10, and
Corollary~7.11]{BJ22}, which bound each mixed pairing of finite-energy
pairs from below by a fixed positive combination of the pure energies; they
are not deduced here.
\end{proof}

\begin{example}
\label{ex:ch46-e1-curve}
Let $X=\mathbf P^1$, $\omega=c_1(\cO(1))$, $V=1$, and $p\in X(k)$.
Let $\phi$ vanish off the branch of $p$, with
$\phi(t\,\ord_p)=h(t)$, where $h:[0,\infty)\to\mathbf R$ is continuous,
convex and decreasing, and $h(0)=0$.
At the proper-support endpoint set
$\phi(v_{p,\triv})=h(\infty):=\lim_{t\to\infty}h(t)$.
The curve criterion \cite[Example~4.13]{BJ22}, recalled in
\Cref{ex:ch34-psh-check}, says that $\phi$ is $\omega$-psh exactly when
$h'_+(0)\geq-1$.  Assume this condition.

Write $D^2h$ for the distributional second derivative on $(0,\infty)$,
viewed as a measure on the branch.  It is a positive measure; at a
corner its mass is the upward jump of the slope.  Put
$\ell=\lim_{t\to\infty}h'_+(t)\in[-1,0]$.
The graph measure attached to $h$ is
\[
 \nu_h=(1+h'_+(0))\delta_{v_{\triv}}
          +D^2h-\ell\,\delta_{v_{p,\triv}}.
\]
Its total mass is
$(1+h'_+(0))+(\ell-h'_+(0))-\ell=1$.
Writing $h''(t)\,\rd t$ would lose the corner atoms unless this
second derivative were understood distributionally.
For finite-energy $\phi$, the graph measure agrees with
$\MA_\omega(\phi)$; this is \cite[Example~7.20]{BJ22}, or the
decreasing model construction in \Cref{thm:bj-mixed-ma}.

We prove the energy formula using only finite model calculations.
For an integer $R\geq1$, define $h_R(t)=h(\min\{t,R\})$.
On $[0,R]$, interpolate $h$ linearly on the dyadic subdivisions and
extend each interpolant constantly beyond $R$.  Call these functions
$h_{R,j}$.  Convexity gives
$h_{R,j}\downarrow h_R$ uniformly; their slopes lie in $[-1,0]$.
They have rational breakpoints and real coefficients, so their branch
potentials $\phi_{R,j}$ belong to
$\PL(X)_{\mathbf R}\cap\PSH(\omega)$.
The model pairing and integration by parts on this finite interval give
\[
 E_\omega(\phi_{R,j})
    =\frac12\int_{(0,\infty)}h_{R,j}\,\rd(D^2h_{R,j})
    =-\frac12\int_0^R|h'_{R,j}(t)|^2\,\rd t.
\]
The jump to zero slope at $R$ is included in $D^2h_{R,j}$.
There is no boundary term at infinity, and the boundary term at zero
vanishes because $h_{R,j}(0)=0$.

The slopes converge to $h'$ almost everywhere on $[0,R]$ and are
bounded in absolute value by one.  Dominated convergence and decreasing
continuity of energy therefore give
$E_\omega(\phi_R)=-\tfrac12\int_0^R|h'|^2\,\rd t$.
Now $\phi_R\downarrow\phi$ on the entire analytification, including the
endpoint.  Applying
\Cref{prop:ch46-energy-usc}\,\cref{itm:ch46-Eusc:decreasing} once more
proves
\[
 E_\omega(\phi)=-\frac12\int_0^\infty|h'(t)|^2\,\rd t
       \in[-\infty,0],
 \qquad
 \phi\in\cE^1(\omega)\ \Longleftrightarrow\ h'\in L^2(0,\infty).
\]
No product of a limiting slope with the possibly infinite endpoint
value has been used.  The same proof on a smooth curve with
$\deg\omega=d>0$ allows initial slope at least $-d$ and replaces
the coefficient $1/2$ by $1/(2d)$.

For $h(t)=-\log(1+t)$, the squared derivative has integral one, so
$E_\omega(\phi)=-1/2$ although $\phi(v_{p,\triv})=-\infty$.
For $h(t)=-t$, the integral diverges and $E_\omega(\phi)=-\infty$;
the graph measure is $\nu_h=\delta_{v_{p,\triv}}$.
The latter is a graph measure outside the finite-energy domain of
\Cref{thm:bj-mixed-ma}.
The two profiles show that $\cE^1(\omega)$ lies strictly between
the bounded $\omega$-psh functions and $\PSH(\omega)$, as in
\Cref{ex:ch26-unbounded-finite-energy,ex:ch26-infinite-energy}.

Vanishing of the endpoint atom alone is insufficient for finite energy.
For $h(t)=-2(\sqrt{1+t}-1)$ the limiting slope is zero, while
$\int_0^\infty|h'|^2\,\rd t=\int_0^\infty(1+t)^{-1}\,\rd t=\infty$.
\end{example}

Let $\vec{\cE}^{\,1}(X)$ be the vector space of functions on
$X^{\lin}$ generated by all $\cE^1(\omega)$, with $\omega$ ample.
These functions are real-valued on $X^{\lin}$, so differences are
defined there.  The finite-energy estimates give
$\cE^1(\omega)=\vec{\cE}^{\,1}(X)\cap\PSH(\omega)$ for ample $\omega$;
this is \cite[Lemma~7.13]{BJ22}.  We use the real-valued symmetric
multilinear extension of \cite[Theorem~7.14]{BJ22}:
\[
 \bigl(N^1(X)\times\vec{\cE}^{\,1}(X)\bigr)^{n+1}
       \longrightarrow\mathbf R.
\tag{46.17}\label{eq:bj-e1-polarized-pairing}
\]
It agrees with the model intersection pairing and with the pairing on
ample finite-energy psh tuples.  For any fixed ample $\omega$ and fixed
classes $\theta_0,\dots,\theta_n$, it is continuous along coordinatewise
decreasing nets in $\cE^1(\omega)^{n+1}$ whose limits lie in that space.

\begin{remark}
\label{rem:ch46-vector-space}
A psh function can take the value $-\infty$ at proper-support points.
On $X^{\lin}$ every psh function is finite, by
\cite[Corollary~4.17]{BJ22}.  Restriction to this locus therefore
makes differences meaningful, as required by
\Cref{warn:ch46-no-subtraction}.  Finiteness there does not characterize
Fubini--Study functions: the potential $h(t)=-\log(1+t)$ in
\Cref{ex:ch46-e1-curve} is finite at every valuation but is not piecewise
affine.

The class $\cE^1(\omega)$ is convex but is generally not a vector space.
For example, on $\mathbf P^1$ the bounded potential with profile
$h(t)=-c\min\{t,1\}$, $0<c<1$, belongs to $\cE^1(c_1(\cO(1)))$;
its negative fails convexity along the ray and hence is not psh.
Negating a potential does not produce a psh potential in the negative
background class.

For the extension, the relevant cone consists of \emph{pairs}:
\[
 \mathcal C
 =\bigcup_{\omega\in\operatorname{Amp}(X)}
       \{\omega\}\times\cE^1(\omega).
\]
The addition and change-of-class properties in \Cref{thm:bj-e1-basic},
together with positive homogeneity, make $\mathcal C$ a convex cone
spanning $N^1(X)\times\vec{\cE}^{\,1}(X)$.  The finite pairing on this
cone is additive and positively homogeneous in each slot.  Its extension
by differences is therefore independent of the decompositions: an
equality $a-b=a'-b'$ of cone differences gives $a+b'=a'+b$, to which
additivity applies.  This is the algebraic step in
\cite[Theorem~7.14]{BJ22}; finiteness alone would not establish it.
\end{remark}

\section{Mixed Monge--Amp\`ere measures}
\label{sec:ch46-mixed}

We construct the measure by testing the finite pairing against model
functions.  The point requiring care is positivity for a nonnegative
test function whose curvature need not be positive.

\begin{theorem}\label{thm:bj-mixed-ma}
For $(\theta_i,\phi_i)\in
N^1(X)\times\vec{\cE}^{\,1}(X)$, $1\leq i\leq n$, there is a
unique signed Radon measure
\[
               \bigwedge_{i=1}^n(\theta_i+\ddc\phi_i)
\tag{46.18}\label{eq:bj-mixed-ma}
\]
such that, for all $f\in\PL(X)_{\mathbf R}$,
\[
 \int_{X^{\An}}f\bigwedge_{i=1}^n(\theta_i+\ddc\phi_i)
  =(0,f)\prod_{i=1}^n(\theta_i,\phi_i).
\tag{46.19}\label{eq:bj-mixed-ma-duality}
\]
It is symmetric and multilinear, has total mass
$(\theta_1\cdots\theta_n)$, and is positive when every $\theta_i$ is
nef and $\phi_i\in\PSH(\theta_i)\cap\vec{\cE}^{\,1}(X)$.
For fixed $\theta_i$ and a fixed ample $\omega$, the map
$\cE^1(\omega)^n\to C^0(X^{\An})^\vee$ so defined is continuous along
coordinatewise decreasing nets with limits in $\cE^1(\omega)^n$.
Every $f\in\vec{\cE}^{\,1}(X)$ has an integrable representative against
this measure, obtained by taking a difference of its finite-energy psh
representatives; \eqref{eq:bj-mixed-ma-duality} remains valid and is
independent of that representation.
\end{theorem}

\begin{proof}[Construction]
We use the finite pairing \eqref{eq:bj-e1-polarized-pairing}.  The
construction is \cite[Theorem~7.18 and Proposition~7.19]{BJ22};
we give the positivity and limiting arguments explicitly.

\emph{Step 1: a positive functional.}
Fix nef $\theta_i$ and
$\phi_i\in\PSH(\theta_i)\cap\vec{\cE}^{\,1}(X)$, and put
$\Gamma=\prod_{i=1}^n(\theta_i,\phi_i)$.  Define
\[
 \Lambda(f)=(0,f)\Gamma,\qquad f\in\PL(X)_{\mathbf R}.
\]
This is a finite real linear functional by
\eqref{eq:bj-e1-polarized-pairing}.  To prove positivity, write
\[
 f=f^+-f^-,
 \qquad f^\pm\in\PL(X)_{\mathbf R}\cap\PSH(\omega)
\]
for a common ample class $\omega$.  Such a decomposition follows from
\Cref{lem:ch34-pl-lattice}\,\cref{itm:ch34-pl:difference} and
\Cref{cor:ch34-pl-plus-union}, by taking real positive combinations and
then enlarging the backgrounds to a common ample class.  The two
functions are bounded, hence have finite energy.  Thus
\[
 \Lambda(f)=(\omega,f^+)\Gamma-(\omega,f^-)\Gamma.
\]
This expression uses two psh potentials; it does not require $f$ itself
to be psh.

Suppose $f\geq0$, so $f^+\geq f^-$.  Choose an ample class $\eta$ and
set $\Gamma_\varepsilon=\prod_{i=1}^n
(\theta_i+\varepsilon\eta,\phi_i)$ for $\varepsilon>0$.
Each $\phi_i$ belongs to $\cE^1(\theta_i+\varepsilon\eta)$ by
\cite[Lemma~7.13]{BJ22}.  Monotonicity of the ample pairing
(\Cref{thm:bj-extended-energy-pairing}) gives
\[
 (\omega,f^+)\Gamma_\varepsilon
       -(\omega,f^-)\Gamma_\varepsilon\geq0.
\]
By finite multilinearity this difference is a polynomial in
$\varepsilon$, whose value at zero is $\Lambda(f)$.  Letting
$\varepsilon\downarrow0$ proves $\Lambda(f)\geq0$.
This is the nef-boundary argument of \cite[Lemma~7.15]{BJ22}.

\emph{Step 2: representation by a measure.}
The translation identity gives
$\Lambda(1)=(\theta_1\cdots\theta_n)=:M\geq0$.  Positivity implies
\[
 |\Lambda(f)|\leq M\lVert f\rVert_\infty
 \qquad(f\in\PL(X)_{\mathbf R}).
\]
By \Cref{thm:ch34-pl-density}, $\Lambda$ extends uniquely to a bounded
positive functional on $C^0(X^{\An})$.  Since $X^{\An}$ is compact
Hausdorff, the Riesz representation theorem gives a positive Radon
measure satisfying \eqref{eq:bj-mixed-ma-duality}, of mass $M$.
This also covers $M=0$, when the measure is zero.

\emph{Step 3: arbitrary pairs.}
Every pair $(\theta,\phi)\in N^1(X)\times\vec{\cE}^{\,1}(X)$
is a difference of two pairs in $\mathcal C$.
Indeed, write $\phi=u-v$ with $u,v\in\cE^1(\eta)$ for one ample
$\eta$, and choose $A$ so positive that $A-\eta$ and $A-\theta-\eta$
are ample.  The change-of-class property then gives
\[
 (\theta,\phi)=(A,u)-(A-\theta,v),
 \quad
 u\in\cE^1(A),\quad v\in\cE^1(A-\theta).
\]
Applying this decomposition in each slot defines a finite signed sum
of the positive measures already constructed.  Its integrals against
every PL test function equal the prescribed finite pairing.  Density
therefore proves independence of all decompositions, uniqueness,
symmetry and multilinearity.  Testing with $1$ gives the total mass
$(\theta_1\cdots\theta_n)$ also for arbitrary pairs.

\emph{Step 4: limits and finite-energy tests.}
Let $\phi_{i,j}\downarrow\phi_i$ in a fixed $\cE^1(\omega)$.
For a PL test function, convergence of the integrals follows from
\cite[Theorem~7.14(iii)]{BJ22}, enlarging $\omega$ to accommodate the
two fixed psh functions in the test decomposition from Step~1.
To pass from PL tests to all continuous tests, a bound on total
variation is needed.  For each $i$, choose a fixed ample $A_i$ with
$A_i-\omega$ and $A_i-\theta_i$ ample.  Then
\[
 (\theta_i,\phi_{i,j})
   =(A_i,\phi_{i,j})-(A_i-\theta_i,0)
\]
is a difference of positive pairs for every $j$.  The $2^n$ positive
mixed measures in the expansion have fixed masses, given by
intersections of their fixed background classes.  Their sum bounds
the total variation uniformly in $j$.  Density now proves weak
convergence, including for signed measures.

For integrability, first let $\mu$ be a positive mixed measure of
ample finite-energy psh pairs and let $u\in\cE^1(\eta)$.
Choose a decreasing net $u_j\downarrow u$ in
$\PL(X)_{\mathbf R}\cap\PSH(\eta)$, using
\cite[Theorem~4.15]{BJ22}.
The defining identity reads
\[
 \int u_j\,\rd\mu
   =(\eta,u_j)\Gamma-(\eta,0)\Gamma
   \longrightarrow(\eta,u)\Gamma-(\eta,0)\Gamma\in\mathbf R.
\]
Monotone convergence for Radon measures and decreasing nets of upper
semicontinuous functions, \cite[Lemma~7.17]{BJ22}, applies here.
After restricting to a tail, the $u_j$ have a common finite upper
bound.  The finite limit therefore proves $u\in L^1(\mu)$ and the
duality identity for $u$.  Taking differences of such tests and the
finite signed sums from Step~3 proves the assertion for arbitrary
finite-energy pairs and tests.

There is no ambiguity from the initial definition on $X^{\lin}$.
If $u-v=u'-v'$ there, then $u+v'$ and $u'+v$ agree on divisorial
valuations.  They are psh for a common sufficiently large ample class,
so \Cref{cor:ch35-divisorial-determination} identifies them on $X^{\An}$.
All four functions are finite almost everywhere for the total
variation of the measures just constructed, by the preceding
integrability argument.  Their differences consequently agree almost
everywhere.  This completes the deduction of
\cite[Proposition~7.19(iv)--(v)]{BJ22} from the finite pairing.
\end{proof}

\begin{remark}
\label{rem:ch46-zeroth-slot}
The decomposition in Step~1 is essential.  An arbitrary PL function
need not become psh after increasing a pulled-back ample background.
For example, let
$\mu:\cX=\operatorname{Bl}_{(p,0)}(\mathbf P^1\times\mathbf A^1)
\to\mathbf P^1\times\mathbf A^1$, let $E$ be the exceptional curve,
and put $\omega=c_1(\cO(1))$.  For every $C>0$,
\[
 (C\omega_{\cX}+E)\cdot E=E^2=-1.
\]
Thus $\phi_E$ is not $C\omega$-psh for any $C$, by
\eqref{eq:pl-relative-nefness}.  The positive cone $\PL^+(X)$ in
\Cref{cor:ch34-pl-plus-union} cannot be replaced by all of $\PL(X)$.

Nor can one obtain the functional by taking an infimum of test
integrals over dominating psh models.  On the same curve choose
$c\in\mathbf Q$, $0<c<1$, and set $\phi=\phi_{-cE}$ and $f=-\phi\geq0$.
The model calculation in \Cref{ex:ch46-p1-energy} gives
\[
 \MA^{\mathrm u}_\omega(\phi)
     =(1-c)\delta_{v_{\triv}}+c\delta_{\ord_p},
 \qquad
 \int f\,\MA^{\mathrm u}_\omega(\phi)=c^2.
\]
But $0\geq\phi$ is an admissible psh model competitor and
$\int f\,\MA^{\mathrm u}_\omega(0)=0$.  Since every such integral is
nonnegative, their infimum is zero, not $c^2$.
The infimum formula for a psh energy pairing does not commute with
subtracting the two finite pairings used to define $\Lambda$.
\end{remark}

On a common integrally closed test configuration, if
$\phi_i=\phi_{D_i}$, then

\[
 \bigwedge_{i=1}^n(\theta_i+\ddc\phi_i)
   =\sum_{E\subset\cX_0}
      b_E\bigl((\theta_{1,\cX}+D_1)|_E\cdots
               (\theta_{n,\cX}+D_n)|_E\bigr)\delta_{v_E},
\tag{46.20}\label{eq:bj-model-ma-atomic}
\]

where $b_E=\ord_E(\cX_0)$.  Thus model products are finite atomic
measures on divisorial valuations.  This is
\Cref{prop:ch46-pl-measure}, and the reader should note that the formula
was proved there \emph{before} the extension theorem: the model case is
finitary and is the input, not a special case, of
\Cref{thm:bj-mixed-ma}.

\begin{proposition}
\label{prop:ch46-no-mass-outside-lin}
Let $\theta_i$ be nef,
$\phi_i\in\PSH(\theta_i)\cap\vec{\cE}^{\,1}(X)$, and
$\mu=\bigwedge_{i=1}^n(\theta_i+\ddc\phi_i)$.
Every $u\in\cE^1(\omega)$, with $\omega$ ample, is $\mu$-integrable.
In particular,
\[
 \mu\bigl(\{u=-\infty\}\bigr)=0
 \qquad\text{for every }u\in\cE^1(\omega).
\]
\end{proposition}

\begin{proof}
By \Cref{thm:bj-mixed-ma}, $u$ is integrable against $\mu$, and hence
finite $\mu$-almost everywhere.
\end{proof}

\begin{remark}
\label{rem:ch46-mass-location}
The proposition removes the polar sets of each finite collection of
finite-energy psh functions, so their differences are defined almost
everywhere.  This is the measure-theoretic fact used in the cocycle and
integration-by-parts formulae.
The stronger nonpluripolarity conclusion uses the additional fact that
every pluripolar set is contained in the pole set of a finite-energy
psh function; see \cite[Corollaries~7.48--7.49]{BJ22}.
For instance, the profile $h(t)=-\log(1+t)$ in
\Cref{ex:ch46-e1-curve} is a finite-energy witness for the
proper-support endpoint of its ray.

These facts do not by themselves assert that a measure is carried by
$X^{\lin}$.  In particular, taking a union of the exceptional null sets
over all finite-energy tests is not justified by countable additivity.
Nor does weak convergence of measures carried by $X^{\divis}$ force
their limit to be carried by $X^{\lin}$: mass can escape to an endpoint,
as in \Cref{ex:ch46-homogeneous-mass-escapes}.
The construction above uses integrable representatives, so it requires
no such support inference.
\end{remark}

\begin{example}
\label{ex:ch46-curve-laplacian}
Let $X$ be a smooth projective curve, $\deg\omega=d>0$, and
$\phi\in\cE^1(\omega)$.  With the graph-Laplacian sign of
\Cref{chap:model-calculations}, the unnormalized measure is
\[
             \omega+\ddc\phi=d\,\delta_{v_{\triv}}+\Delta\phi.
\tag{46.21}\label{eq:bj-curve-laplacian}
\]
The second derivative on each branch is a measure, including its
corner atoms, as in \Cref{ex:ch46-e1-curve}.  The equality is the model
intersection formula followed by decreasing approximation;
see \cite[Example~7.20]{BJ22}.
Every model Laplacian has total mass zero.  Weak convergence, tested
against the constant function $1$, preserves this equality, so the
measure in \eqref{eq:bj-curve-laplacian} has mass $d$.
A general finite-energy profile need not be eventually constant.

For $X=\mathbf P^1$ and $\phi=\phi_{-cE}$ with $0<c\leq d$,
\Cref{ex:ch46-p1-energy} gives
\[
 \Delta\phi=-c\,\delta_{v_{\triv}}+c\,\delta_{\ord_p},
 \qquad
 \omega+\ddc\phi=(d-c)\delta_{v_{\triv}}+c\,\delta_{\ord_p}.
\]
The upward jump from slope $-c$ to zero at $\ord_p$ accounts for the
positive atom there.  Reversing the sign of the Laplacian would give
a negative atom at $\ord_p$ and violate positivity.
For finite-energy $\phi$, the proper-support endpoints have no mass
by integrability; they remain points of $X^{\An}$ and are not removed
from the graph.
\end{example}

\begin{remark}[Comparison with the discretely valued mixed measure]
\label{rem:ch46-mixed-part3}
\Cref{def:mixed-model-measure} and \Cref{prop:model-ma-basic} construct the
mixed Chambert--Loir measure of model metrics over a discretely valued
base, and \Cref{thm:ch24-mass} computes its mass; the formula
\eqref{eq:model-mixed-measure} is \eqref{eq:bj-model-ma-atomic} with
$\cX_s$ in place of $\cX_0$, and the mass formula
\eqref{eq:mixed-total-mass} is
\Cref{prop:ch46-pl-measure}\,\cref{itm:ch46-pl-ma:mass}.  The extension to
bounded psh functions is \Cref{thm:bounded-ma-extension} there and Step~4
above here, and the two proofs use the same monotone-convergence device.
The place where the two theories diverge is the domain of the extension:
over a discretely valued base one extends to \emph{bounded} psh metrics and
then to the finite-energy class of \Cref{def:e1} through
\Cref{thm:e1-mixed-products}, whereas here the extension to
$\vec\cE^{\,1}(X)$ is made in one step because
\eqref{eq:bj-energy-infimum} is available for unbounded arguments from the
outset.  This uses fixed-class regularization and finite-energy estimates; it does
not consume the compactness input of \Cref{sec:ch46-where-envelope}.
\end{remark}

\section{Locality, monotonicity and integration by parts}
\label{sec:ch46-ibp}

\subsection{The cocycle formula and its consequences}
\label{sec:ch46-cocycle}

For $\phi,\psi\in\cE^1(\omega)$, multilinearity yields

\[
 E(\phi)-E(\psi)
  =\frac1{n+1}\sum_{j=0}^n
       \int(\phi-\psi)
       V^{-1}(\omega+\ddc\phi)^j
       \wedge(\omega+\ddc\psi)^{n-j}.
\tag{46.22}\label{eq:bj-energy-cocycle}
\]

In particular,

\[
 \int(\phi-\psi)\MA(\phi)
 \leq E(\phi)-E(\psi)
 \leq\int(\phi-\psi)\MA(\psi).
\tag{46.23}\label{eq:bj-energy-tangent-inequalities}
\]

The finite multilinear pairing proves the cocycle identity directly;
concavity then gives the tangent inequalities.

\begin{proposition}
\label{prop:ch46-cocycle}
Let $\omega\in\Amp(X)$ and $\phi,\psi\in\cE^1(\omega)$.  Then
\eqref{eq:bj-energy-cocycle} and
\eqref{eq:bj-energy-tangent-inequalities} hold, all the integrals being
finite; consequently
\begin{enumerate}
\item\label{itm:ch46-coc:mono}
      $E_\omega$ is increasing on $\cE^1(\omega)$, and strictly
      increasing in the sense that $\phi\le\psi$ with
      $E_\omega(\phi)=E_\omega(\psi)$ forces $\phi=\psi$
      $\MA_\omega(\psi)$-almost everywhere;
\item\label{itm:ch46-coc:cocycle}
      the three-term cocycle identity
      \[
       \bigl(E_\omega(\phi)-E_\omega(\psi)\bigr)
       +\bigl(E_\omega(\psi)-E_\omega(\chi)\bigr)
       +\bigl(E_\omega(\chi)-E_\omega(\phi)\bigr)=0
      \]
      holds for $\phi,\psi,\chi\in\cE^1(\omega)$, which is what the name
      records;
\item\label{itm:ch46-coc:lipschitz}
      $|E_\omega(\phi)-E_\omega(\psi)|
        \le\sup_{X^{\divis}}|\phi-\psi|$.
\end{enumerate}
\end{proposition}

\begin{proof}
Put $A=(\omega,\phi)$ and $B=(\omega,\psi)$ in the vector space of
\eqref{eq:bj-e1-polarized-pairing}.  Every pairing below is finite, and
multilinearity gives
\[
 A^{n+1}-B^{n+1}
   =(A-B)\sum_{j=0}^n A^jB^{n-j}
   =(0,\phi-\psi)\sum_{j=0}^n A^jB^{n-j}.
\]
Divide by $(n+1)V$ and use the finite-energy duality in
\Cref{thm:bj-mixed-ma} to obtain \eqref{eq:bj-energy-cocycle}.
For the inequalities, the function
$F(t)=E_\omega((1-t)\psi+t\phi)$ is concave on $[0,1]$ and is a
polynomial there by the same finite pairing.  Its endpoint derivatives
are
\[
 F'(0)=\int(\phi-\psi)\,\MA_\omega(\psi),
 \qquad F'(1)=\int(\phi-\psi)\,\MA_\omega(\phi).
\]
The concavity inequalities $F'(1)\leq F(1)-F(0)\leq F'(0)$ give
\eqref{eq:bj-energy-tangent-inequalities}.
All integrands are defined almost everywhere by finite-energy
integrability.  This argument does not interchange a varying unbounded
test function with a weakly convergent sequence of measures.

\cref{itm:ch46-coc:mono}.  Monotonicity is
\eqref{eq:bj-energy-cocycle} with a nonnegative integrand and positive
measures.  If $E_\omega(\phi)=E_\omega(\psi)$ with $\phi\le\psi$, then
the $j=0$ term of \eqref{eq:bj-energy-cocycle}, namely
$\int(\phi-\psi)\MA_\omega(\psi)$, vanishes; the integrand is
$\le0$ and the measure is positive, so $\phi=\psi$ almost everywhere
for it.

\cref{itm:ch46-coc:cocycle} is immediate from the fact that each bracket is
a difference of two finite numbers.

\cref{itm:ch46-coc:lipschitz}.  Set
$C=\sup_{X^{\divis}}|\phi-\psi|$.  The case $C=\infty$ is immediate.
Otherwise, $\phi\leq\psi+C$ and $\psi\leq\phi+C$ on divisorial
valuations imply the same inequalities on $X^{\An}$ by
\Cref{cor:ch35-divisorial-determination}.  Thus
$|\phi-\psi|\leq C$ almost everywhere for both probability measures
in \eqref{eq:bj-energy-tangent-inequalities}.  These inequalities give
the claimed energy bound.
\end{proof}

\begin{remark}[The same identity in \Cref{part:bfj}]
\label{rem:ch46-cocycle-part3}
\Cref{prop:energy-cocycle} and \Cref{cor:energy-monotone-concave} are
\eqref{eq:bj-energy-cocycle} and
\eqref{eq:bj-energy-tangent-inequalities} over a discretely valued base,
in the unnormalized convention: there the sandwich
\eqref{eq:energy-sandwich} reads with $\MA$ of mass $V$, and the
factor $V^{-1}$ that appears in \eqref{eq:bj-energy-cocycle} is exactly
the conversion of \Cref{warn:ch46-normalization}.  The scopes differ:
\Cref{prop:ch26-ibp-regimes} distinguishes three regimes over a discretely
valued base --- model, bounded, finite energy --- because the extension
there is built in those three stages, whereas here
\eqref{eq:bj-energy-infimum} covers the unbounded case at once.  The
compensating cost is that here one must check integrability separately,
which is \Cref{prop:ch46-no-mass-outside-lin}.
\end{remark}

\subsection{Integration by parts}
\label{sec:ch46-ibp-proper}

\begin{proposition}\label{prop:bj-integration-by-parts}
For $f,g\in\vec{\cE}^{\,1}(X)$ and finite-energy pairs
$(\theta_i,\phi_i)$,

\[
 \int f\,\ddc g\wedge\bigwedge_{i=1}^{n-1}(\theta_i+\ddc\phi_i)
 =\int g\,\ddc f\wedge\bigwedge_{i=1}^{n-1}(\theta_i+\ddc\phi_i).
\tag{46.24}\label{eq:bj-integration-by-parts}
\]
\end{proposition}

\begin{proof}
Apply the finite-energy duality of \Cref{thm:bj-mixed-ma} to the signed
measure with first pair $(0,g)$ and remaining pairs
$(\theta_i,\phi_i)$.  Its integral against $f$ is
\[
 (0,f)(0,g)\prod_{i=1}^{n-1}(\theta_i,\phi_i).
\]
Interchanging $f$ and $g$ gives the other side of
\eqref{eq:bj-integration-by-parts}, with the same value by symmetry of
\eqref{eq:bj-e1-polarized-pairing}.
Both tests are integrable against the corresponding signed measures.
Their almost-everywhere representatives and the independence of these
representatives were established in \Cref{thm:bj-mixed-ma}.
\end{proof}

\begin{corollary}
\label{cor:ch46-dirichlet-e1}
For $f,g\in\vec\cE^{\,1}(X)$ and finite-energy pairs
$(\theta_i,\phi_i)$, $1\le i\le n-1$, with all $\theta_i$ nef and
$\phi_i$ psh, the pairing
\[
 \langle f,g\rangle_{\mathrm D}
 :=-\int_{X^{\An}}f\;\ddc g\wedge\bigwedge_{i=1}^{n-1}(\theta_i+\ddc\phi_i)
\]
is symmetric, positive semidefinite, and satisfies the Cauchy--Schwarz
inequality
$|\langle f,g\rangle_{\mathrm D}|^2\le\langle f,f\rangle_{\mathrm D}
\langle g,g\rangle_{\mathrm D}$.
\end{corollary}

\begin{proof}
Symmetry follows from \Cref{prop:bj-integration-by-parts}.
Write $f$ and $g$ as differences of finite-energy psh functions in a
common ample class $\eta$.
First suppose the other backgrounds $\theta_i$ are ample.
Approximate each psh constituent of $f,g$ and each $\phi_i$
decreasingly by real PL psh functions in its own background.
The model Dirichlet form is positive semidefinite by
\Cref{prop:ch46-dirichlet-vertical}\,\cref{itm:ch46-dir:pos}.
After expanding the differences, each entry of its two-by-two Gram
matrix converges by the finite-pairing continuity in
\cite[Theorem~7.14(iii)]{BJ22}, using a common larger ample class.
Its limiting Gram matrix is therefore positive semidefinite.

For nef $\theta_i$, replace them by $\theta_i+\varepsilon\eta$.
The potentials belong to $\cE^1(\theta_i+\varepsilon\eta)$ by
\cite[Lemma~7.13]{BJ22}, so the previous argument applies.
The entries of the resulting Gram matrix are polynomials in
$\varepsilon$, by finite multilinearity.
Their limits as $\varepsilon\downarrow0$ give the desired positive
semidefinite matrix.  Its determinant is nonnegative, which is the
Cauchy--Schwarz inequality.
\end{proof}

\subsection{Locality and the comparison principle}
\label{sec:ch46-locality}

\begin{theorem}\label{thm:bj-locality-comparison}
For $\phi,\psi\in\cE^1(\omega)$,

\[
 \mathbf1_{\{\phi>\psi\}}\MA(\max\{\phi,\psi\})
 =\mathbf1_{\{\phi>\psi\}}\MA(\phi),
\tag{46.25}\label{eq:bj-locality}
\]

and

\[
 \int_{\{\phi<\psi\}}\MA(\psi)
 \leq\int_{\{\phi<\psi\}}\MA(\phi).
\tag{46.26}\label{eq:bj-comparison-principle}
\]
\end{theorem}

\begin{proof}
\emph{Locality for rational PL data.}
Let $\phi,\psi\in\PL(X)\cap\PSH(\omega)$ and
$u=\max\{\phi,\psi\}$. Choose an integrally closed test configuration
determining all three functions by rational vertical divisors.
The crossing lemma \cite[Lemma~7.43]{BJ22} says that if
$\phi(v_E)>\psi(v_E)$, the divisors determining $u$ and $\phi$ agree
in a neighbourhood of $E$. Their restrictions to $E$ therefore agree,
so the atomic formula gives the same coefficient at $v_E$.
Applying this to each atom proves locality for these PL data.
This common-model argument does not require resolution of singularities.

\emph{The finite-energy limit.}
Choose decreasing nets $\phi_j,\psi_j\in\cH_{\mathrm{dom}}(\omega)$
converging to $\phi,\psi$, indexed by one product directed set. Put
$u_j=\max\{\phi_j,\psi_j\}\downarrow u=\max\{\phi,\psi\}$.
The limit $u$ has finite energy by upward closure.
For each $j$, the PL identity implies
\[
 (u_j-\psi_j)\MA_\omega(u_j)
       =(u_j-\psi_j)\MA_\omega(\phi_j).
\]
To pass to the limit, use the weighted convergence theorem
\cite[Proposition~7.21(ii)]{BJ22}: for decreasing finite-energy
potentials $a_j\downarrow a$ and measure slots decreasing to a
finite-energy tuple, the signed measures $a_j\mu_j$ converge weakly
to $a\mu$, provided $a$ is integrable for $\mu$.
Here $a_j$ is successively $u_j$ and $\psi_j$, and $\mu_j$ is
successively $\MA_\omega(u_j)$ and $\MA_\omega(\phi_j)$.
Integrability holds by \Cref{thm:bj-mixed-ma}.
Subtracting the two weighted limits gives
\[
 f\MA_\omega(u)=f\MA_\omega(\phi),
 \qquad f=u-\psi\ge0
\]
almost everywhere for these measures. On $\{\phi>\psi\}$ both
potentials are finite almost everywhere and $f>0$.
Restrict the identity there and multiply by $f^{-1}$ to obtain
\eqref{eq:bj-locality}. This is the argument of
\cite[Theorem~7.40]{BJ22}; ordinary weak convergence without the
weighted convergence input would not suffice.

\emph{Comparison from locality.}
Fix $\varepsilon>0$ and put $u=\max\{\phi,\psi-\varepsilon\}$.
Finite maxima preserve psh membership and upward closure gives finite
energy. Locality on the two disjoint Borel sets yields
\[
 \begin{aligned}
 \mathbf1_{\{\phi>\psi-\varepsilon\}}\MA(u)
   &=\mathbf1_{\{\phi>\psi-\varepsilon\}}\MA(\phi),\\
 \mathbf1_{\{\phi<\psi-\varepsilon\}}\MA(u)
   &=\mathbf1_{\{\phi<\psi-\varepsilon\}}\MA(\psi).
 \end{aligned}
\]
since constants do not change the measure. Their total masses are one,
so
\[
 \int_{\{\phi<\psi-\varepsilon\}}\MA(\psi)
 \le 1-\int_{\{\phi>\psi-\varepsilon\}}\MA(\phi)
 =\int_{\{\phi\le\psi-\varepsilon\}}\MA(\phi).
\]
As $\varepsilon\downarrow0$, both sets on the two sides increase to
$\{\phi<\psi\}$, outside the union of the two pole sets. That union is
null for both measures by \Cref{thm:bj-mixed-ma}. Continuity from below
therefore gives \eqref{eq:bj-comparison-principle}.
The comparison loci are defined by the extended-real order on all of
$X^{\An}$ and are Borel. No assertion that the measures are carried by
$X^{\lin}$ is needed.
\end{proof}

\begin{remark}[Why weighted convergence is needed]
\label{rem:ch46-locality-obstruction}
For unbounded psh functions, $\{\phi>\psi\}$ is generally a Borel set,
not an open set in the Berkovich topology. Weak convergence of measures
controls continuous tests and does not allow one simply to pass a
restriction to such a set through a limit. The weighted convergence of
\cite[Proposition~7.21(ii)]{BJ22} resolves this problem in the proof above.
Its limit is justified by finite-energy integrability and convergence of
the corresponding pairings.

At the model level, agreement of values at one divisorial point alone
would not determine an intersection number on its component. The
crossing lemma controls the neighbouring components as well, giving
agreement of the restricted line bundles. This is why a common model
for the two functions and their maximum is required. On a curve the
same conclusion can be read directly from the outgoing slopes of the
graph Laplacian; the finite-energy passage still needs the limiting
argument just given.
\end{remark}

\begin{corollary}
\label{cor:ch46-contact-mass}
For $\phi,\psi\in\cE^1(\omega)$ and $u=\max\{\phi,\psi\}$,
\[
 \MA(u)\ \ge\ \mathbf1_{\{\phi>\psi\}}\MA(\phi)
              +\mathbf1_{\{\psi>\phi\}}\MA(\psi)
\]
as measures, and consequently
\[
 \int_{\{\phi=\psi\}}\MA(u)
 =1-\int_{\{\phi>\psi\}}\MA(\phi)-\int_{\{\psi>\phi\}}\MA(\psi).
\]
\end{corollary}

\begin{proof}
The two Borel sets $\{\phi>\psi\}$ and $\{\psi>\phi\}$ are disjoint, and
\eqref{eq:bj-locality} identifies $\MA(u)$ with $\MA(\phi)$ on the
first and with $\MA(\psi)$ on the second; adding the two restrictions and
discarding the nonnegative remainder on $\{\phi=\psi\}$ gives the
inequality.  The second display follows by partitioning $X^{\An}$ into these
two sets and the equality set, using that $\MA(u)$ is a probability
measure.
\end{proof}

\begin{remark}[What is deferred to the variational chapter]
\label{rem:ch46-not-proved-locality}
The domination principle and normalized uniqueness are treated in
\Cref{chap:bj-finite-energy}, together with the energy estimates that
support the variational method. The domination principle has the
following orientation: if $\phi\in\PSH(\omega)$, $\psi\in\cE^1(\omega)$,
and $\phi\le\psi$ holds $\MA_\omega(\psi)$-almost everywhere, then it
holds everywhere; this is \cite[Corollary~10.6]{BJ22}.

Orthogonality and differentiability of energy composed with an obstacle
envelope also belong there. This division does not assert that all
these results require unbounded obstacles or the envelope property.
The finite-energy measure and comparison results above are available
without it; the compact variational existence argument uses the
additional envelope inputs identified in \Cref{sec:ch46-where-envelope}.
\end{remark}

\section{Homogeneous compatibility and the scaling action}
\label{sec:ch46-homogeneous}

The preceding constructions have counterparts over a discretely valued
base. This section uses the scaling action that preserves the fixed
trivially valued ground field.  This section is
where the trivially valued theory acquires a structure that
\Cref{chap:bfj-energy} does not have.  The scaling action of
\eqref{eq:valuation-scaling} and \eqref{eq:function-scaling-action} makes
$X^{\An}$ a cone, makes $\PSH(\theta)$ an $\mathbf R_{>0}$-set, and
makes the energy homogeneous of degree one.

\subsection{Energy is homogeneous of degree one}
\label{sec:ch46-scaling-energy}

\begin{theorem}
\label{thm:ch46-energy-homogeneous}
For $t\in\mathbf Q_{>0}$ and pairs $(\theta_i,\phi_i)$ that are psh at
PL level,
\[
 (\theta_0,t\cdot\phi_0)\cdots(\theta_n,t\cdot\phi_n)
 =t\,(\theta_0,\phi_0)\cdots(\theta_n,\phi_n),
\]
where $t\cdot\phi$ is the scaling action
\eqref{eq:function-scaling-action}.  For every real $t>0$, the pure energy
satisfies
\[
 E_\omega(t\cdot\phi)=t\,E_\omega(\phi)
 \qquad\text{for every }\phi\in\PSH(\omega),
\]
and $\cE^1(\omega)$ is stable under the scaling action.
\end{theorem}

\begin{proof}
\emph{Step 1: integers.}  Let $t=d\in\mathbf Z_{>0}$.  By
\eqref{eq:scaling-fs-weights} the action multiplies all weights by $d$,
and by the discussion following
\Cref{thm:ch33-intersection-invariance} this is the effect on the
associated function of the degree-$d$ base change $\cX_d\to\cX$ along
$\varpi\mapsto\varpi^d$.  Formula
\eqref{eq:ch33-intersection-base-change} then gives
$(\overline{\cL_{0,d}}\cdots\overline{\cL_{n,d}})
 =d\,(\overline{\cL_0}\cdots\overline{\cL_n})$; the mixed form follows
from the top self-intersection form by polarization, both sides being
symmetric multilinear.

\emph{Step 2: rationals.}  For $t=p/q$ apply Step~1 twice:
$q\cdot(t\cdot\phi_i)=p\cdot\phi_i$ by
\Cref{lem:ch35-action-basic}\,(1), so
$q\,(\theta_0,t\cdot\phi_0)\cdots=(\theta_0,p\cdot\phi_0)\cdots
 =p\,(\theta_0,\phi_0)\cdots$.

\emph{Step 3: reals, at the level of functions.}  For irrational $t$ the
function $t\cdot\phi$ need \emph{not} lie in $\PL_{\mathbf R}(X)$: by
\Cref{lem:ch34-pl-lattice}\,(1) the scaling action preserves
$\PL^+(X)$ only for $t\in\mathbf Q_{>0}$, and indeed
$(t\cdot\phi_c)(s\,\ord_p)=-c\min\{s,t\}$ of
\Cref{ex:ch46-homogeneity-p1} has its corner at the quasimonomial point
$t\,\ord_p$, which is divisorial only for rational $t$.  The first
display of the theorem is therefore asserted at PL level for
$t\in\mathbf Q_{>0}$, where Steps~1 and~2 prove it.

What survives for every real $t$ is the estimate.  Writing
$\phi=\sum_kc_k\phi_{\mathfrak a_k}$ and using
\eqref{eq:ch34-flag-weight-form} together with
\Cref{lem:ch35-action-basic}\,(6),
\[
 t\cdot\phi_{\mathfrak a}
 =\max_{\lambda}\{\log|\mathfrak a_\lambda|+t\lambda\},
\]
a finite maximum of functions depending affinely on $t$ with bounded
coefficients; hence
$\lVert t\cdot\phi-t'\cdot\phi\rVert_\infty
 \le|t-t'|\sum_k|c_k|\max_\lambda|\lambda|$, so $t\mapsto t\cdot\phi$ is
continuous into $C^0(X^{\An})$.  Moreover
$t\cdot\phi\in\PSH(\omega)$ when $\phi\in\PL^{\omega}(X)$,
by \Cref{lem:ch35-action-basic}\,(6), and it is continuous and bounded.
No presentation in a single line bundle is asserted for a general real
ample background.

\emph{Step 4: psh functions.}  The action preserves $\PSH(\theta)$, by
\Cref{lem:ch35-action-basic}\,(6) applied to an approximating net together
with \Cref{lem:ch35-action-basic}\,(4).  For $\phi\in\PL^{\omega}(X)$ the
map $t\mapsto E_\omega(t\cdot\phi)$ is continuous, $E_\omega$ being
$1$-Lipschitz on bounded psh functions by monotonicity and translation,
and $t\mapsto t\cdot\phi$ continuous
into $C^0(X^{\An})$ by Step~3; so
$E_\omega(t\cdot\phi)=tE_\omega(\phi)$ for every $t>0$, from the
rational case.  For general $\phi$, choose by
\Cref{thm:ch34-fs-regularization} a decreasing net
$\phi_i\in\cH_{\mathrm{dom}}(\omega)$ with $\phi_i\downarrow\phi$.
The action is an order isomorphism by
\Cref{lem:ch35-action-basic}\,(2), and directly from its definition
$t\cdot\phi_i\downarrow t\cdot\phi$.  The bounded case gives
$E_\omega(t\cdot\phi_i)=tE_\omega(\phi_i)$, while continuity of the
energy along decreasing nets gives
\[
 E_\omega(t\cdot\phi)
 =\lim_iE_\omega(t\cdot\phi_i)
 =t\lim_iE_\omega(\phi_i)
 =tE_\omega(\phi).
\]
Thus no envelope input is hidden in the passage to arbitrary $\phi$.
Since $t>0$, the value $-\infty$ is preserved, so
$\cE^1(\omega)$ is stable.
\end{proof}

\begin{example}[Homogeneity checked on $\mathbf P^1$]
\label{ex:ch46-homogeneity-p1}
Keep the data of \Cref{ex:ch46-p1-energy} with $d=1$ and $0<c\le1$, so that
$\phi=\phi_{-cE}$ has $\phi(s\,\ord_p)=-c\min\{s,1\}$ and
$E_\omega(\phi)=-c^2/2$.  By \eqref{eq:function-scaling-action},
\[
 (t\cdot\phi)(s\,\ord_p)=t\,\phi(t^{-1}s\,\ord_p)
 =-tc\min\{s/t,1\}=-c\min\{s,t\},
\]
so $t\cdot\phi=\max\{c\log|\mathfrak m_p|,-ct\}$.  Its graph Laplacian
has outgoing slope $-c$ at $v_{\triv}$ and $+c$ at $t\,\ord_p$, so
\[
 \MA^{\mathrm u}_\omega(t\cdot\phi)
 =(1-c)\,\delta_{v_{\triv}}+c\,\delta_{t\ord_p}
 =t_\star\bigl(\MA^{\mathrm u}_\omega(\phi)\bigr),
\]
since $t_\star$ fixes $v_{\triv}$ and carries $\ord_p$ to
$t\,\ord_p$; this is \eqref{eq:bj-homogeneous-ma-scaling}.  The energy is
\[
 E_\omega(t\cdot\phi)
 =\tfrac12\Bigl[(1-c)\cdot0+c\cdot(-ct)\Bigr]
 =-\tfrac{tc^2}{2}=t\,E_\omega(\phi),
\]
as \Cref{thm:ch46-energy-homogeneous} requires.  Notice what the
computation does \emph{not} say: the probability measure
$\MA_\omega(\phi)$ is not invariant under $t_\star$, it is
\emph{transported} by it; the atom moves from $\ord_p$ to $t\,\ord_p$
while its mass $c$ is unchanged.
\end{example}

Let $\widehat\phi$ be the homogeneous extension of a psh function from
\Cref{chap:bj-homogeneity-bdivisors}.  The scaling action $t_\star$ on
valuations satisfies

\[
 \widehat\phi(t v)=t\widehat\phi(v),\qquad
 \bigwedge_i(\theta_i+\ddc(t\cdot\phi_i))
       =t_\star\bigwedge_i(\theta_i+\ddc\phi_i).
\tag{46.27}\label{eq:bj-homogeneous-ma-scaling}
\]

The second identity is \cite[Proposition~7.19(iii)]{BJ22}.  It says that
homogenization is compatible with the cone geometry of $X^{\An}$; it does
not say that a probability measure is invariant under scaling.  Its pushout
by $t_\star$ is generally different.

\subsection{Homogeneous potentials have infinite energy}
\label{sec:ch46-homogeneous-energy}

The first consequence of \Cref{thm:ch46-energy-homogeneous} is a negative
one: homogeneous potentials generally leave the finite-energy domain.
The coefficient dictionary and the measure construction must therefore
be compared through bounded truncations, with the distinct domains
recorded in \Cref{sec:ch46-bdivisor-bridge}.

\begin{corollary}
\label{cor:ch46-homogeneous-infinite-energy}
Let $\omega\in\Amp(X)$ and let $\psi\in\PSH(\omega)$ be homogeneous
in the sense of \Cref{def:homogeneous-function}, that is,
$t\cdot\psi=\psi$ for every $t>0$.  Then:
\begin{enumerate}
\item\label{itm:ch46-hom:dichotomy}
      $E_\omega(\psi)\in\{0,-\infty\}$;
\item\label{itm:ch46-hom:bounded}
      $\psi$ is bounded if and only if $\psi\equiv0$; a nonzero
      homogeneous $\omega$-psh function is unbounded below;
\item\label{itm:ch46-hom:curve}
      if $\dim X=1$, then $E_\omega(\psi)=-\infty$ for every
      $\psi\ne0$; so $\cE^1(\omega)$ contains no nonzero
      homogeneous function.
\end{enumerate}
\end{corollary}

\begin{proof}
\cref{itm:ch46-hom:dichotomy}.  By \Cref{thm:ch46-energy-homogeneous},
$E_\omega(\psi)=E_\omega(t\cdot\psi)=t\,E_\omega(\psi)$ for every
$t>0$; a real number equal to all its positive multiples is $0$, so
$E_\omega(\psi)\in\{0,-\infty\}$.

\cref{itm:ch46-hom:bounded}.  A homogeneous $\omega$-psh function
satisfies $\psi\le0$ and $\psi(v_{\triv})=0$, by
\Cref{prop:ch35-homogeneous-normalization}.  Suppose $\psi\ge-M$ for a
real $M$ and let $v\in X^{\An}$ be a point with $\psi(v)<0$.
Homogeneity gives $\psi(tv)=t\,\psi(v)\to-\infty$ as $t\to+\infty$,
contradicting the lower bound.  Hence $\psi\ge0$, so $\psi\equiv0$.
The converse is trivial.

\cref{itm:ch46-hom:curve}. Let $\dim X=1$, $d=\deg\omega$, and let
$\psi\ne0$ be homogeneous $\omega$-psh. We may work on the smooth
normalization: pullback preserves homogeneity and psh membership, and
energy is unchanged by \cite[Proposition~7.7(iv)]{BJ22}.
We use the same notation for the pulled-back data.  By
\Cref{cor:ch35-divisorial-determination} there is a closed point $q$ and a
$t_0>0$ with $\psi(t_0\ord_q)<0$; homogeneity then forces
$\psi(t\ord_q)=t\,\psi(\ord_q)=-a_qt$ with
$a_q:=-\psi(\ord_q)>0$.  The curve criterion
\cite[Example~4.13]{BJ22}, recalled at the end of
\Cref{ex:ch34-psh-check}, gives
$d+\sum_{q'}\psi'_{q'}(0^+)\ge0$, that is $\sum_{q'}a_{q'}\le d$; in
particular $a_q\le d$.  The function
$\psi_q$ equal to $-a_qt$ on the ray of $q$ and to $0$
elsewhere is therefore $\omega$-psh, being convex on every ray with
$d-a_q\ge0$; and $\psi\le\psi_q$, the two
agreeing on the ray of $q$ while $\psi\le0=\psi_q$ elsewhere.
Monotonicity of $E_\omega$ and the Dirichlet-integral computation of
\Cref{ex:ch46-e1-curve}, which for $\deg\omega=d$ carries the extra
factor $V^{-1}=d^{-1}$, give for the profile $h(t)=-a_qt$
\[
 E_\omega(\psi)\le E_\omega(\psi_q)
 =-\frac1{2d}\int_0^\infty a_q^2\,\rd t=-\infty .
\]
\end{proof}

\begin{remark}[The higher-dimensional statement]
\label{rem:ch46-homogeneous-higher}
The exclusion of nonzero homogeneous functions from $\cE^1(\omega)$
holds in every dimension once the domination principle of
\Cref{chap:bj-finite-energy} is available. If the dichotomy gives
$E_\omega(\psi)=0$, then $\psi\in\cE^1(\omega)$ and $\psi\le0$.
The cocycle comparing $\psi$ with $0$ has nonpositive terms and sum zero.
Its term with all measure slots equal to $\psi$ gives
$\int\psi\,\MA_\omega(\psi)=0$, hence $\psi=0$ almost everywhere
for $\MA_\omega(\psi)$. Apply domination to $0$ and $\psi$ to obtain
$0\le\psi$ everywhere. Thus $\psi=0$.
This is a consequence deferred to that chapter, not a conjectural
higher-dimensional extension of the curve computation.
\end{remark}

\begin{example}[Where the mass goes]
\label{ex:ch46-homogeneous-mass-escapes}
Let $X=\mathbf P^1$, $\omega=c_1(\cO(1))$, $V=1$, $p\in X(k)$, and
$\psi:=\log|\mathfrak m_p|$, which is homogeneous, $\omega$-psh
($h(t)=-t$ has $h'(0^+)=-1$), and has $E_\omega(\psi)=-\infty$ by
\Cref{ex:ch46-e1-curve}, or by
\Cref{cor:ch46-homogeneous-infinite-energy}\,\cref{itm:ch46-hom:curve}.  Its truncations
$\psi^{(t)}:=\max\{\psi,-t\}=t\cdot\phi_{\{p\}}$ are bounded and
$\omega$-psh --- PL exactly when $t\in\mathbf Q_{>0}$, which suffices,
only $t\to+\infty$ along rationals being used ---, and \Cref{ex:ch46-p1-energy} with $d=1$, $c=1$ and
\Cref{ex:ch46-homogeneity-p1} give
\[
 \MA^{\mathrm u}_\omega(\psi^{(t)})=\delta_{t\ord_p},
 \qquad
 E_\omega(\psi^{(t)})=-\frac t2 .
\]
As $t\to+\infty$ the point $t\,\ord_p$ converges in $X^{\An}$ to the
endpoint of the ray, which is the proper-support semivaluation
$v_{p,\triv}$ of \eqref{eq:trivial-endpoints}; hence
\[
 \MA^{\mathrm u}_\omega(\psi^{(t)})\longrightarrow\delta_{v_{p,\triv}}
 \quad\text{weakly},
 \qquad
 E_\omega(\psi^{(t)})\longrightarrow-\infty .
\]
The total mass $1=V$ is conserved, but in the limit it sits at a
proper-support point.  Two lessons.  First, a computation of
$\omega+\ddc\psi$ that uses only the graph Laplacian on the ray gives
$1\cdot\delta_{v_{\triv}}-1\cdot\delta_{v_{\triv}}=0$ and \emph{loses}
the mass: the endpoint of an infinite ray must be included, and it is a
point of $X^{\An}$, not a boundary artefact.  Second, the integrability statement
\Cref{prop:ch46-no-mass-outside-lin} applies to finite-energy tuples.
Here the energy is $-\infty$, and the graph measure charges the pole of
the finite-energy test in \Cref{ex:ch46-e1-curve}. It therefore falls
outside that theorem's domain.
\end{example}

\subsection{Invariance and the support of a homogeneous mixed measure}
\label{sec:ch46-homogeneous-support}

\begin{proposition}[Conditional limit of homogeneous truncations]
\label{prop:ch46-homogeneous-support}
Let $\theta_1,\ldots,\theta_n$ be nef and let
$\psi_i\in\PSH(\theta_i)$ be homogeneous. For $R>0$ set
\[
 \psi_{i,R}=\max\{\psi_i,-R\},\qquad
 \mu_R=\bigwedge_{i=1}^n(\theta_i+\ddc\psi_{i,R}).
\]
Each $\mu_R$ is a positive Radon measure of mass
$(\theta_1\cdots\theta_n)$. Suppose the whole family $\mu_R$ converges
weakly as $R\to+\infty$ to a Radon measure $\mu$. Then:
\begin{enumerate}
\item\label{itm:ch46-homsupp:inv}
      $t_\star\mu=\mu$ for every $t>0$;
\item\label{itm:ch46-homsupp:supp}
      $\mu$ is carried by the fixed points of the scaling action,
      namely $\{v_{Y,\triv}:Y\subset X\text{ irreducible closed}\}$.
\end{enumerate}
No existence of this full weak limit, or extension of the finite-energy
operator to the unbounded tuple, is asserted here.
\end{proposition}

\begin{proof}
The functions $\psi_{i,R}$ are bounded and $\theta_i$-psh. After enlarging
$\theta_i$ to an ample class they have finite energy, so
\Cref{thm:bj-mixed-ma} defines the positive measure $\mu_R$ with the
stated mass. Homogeneity gives
$t\cdot\psi_{i,R}=\psi_{i,tR}$. Equivariance on these bounded tuples
therefore gives $t_\star\mu_R=\mu_{tR}$.
For a fixed $t>0$, both $R$ and $tR$ tend to infinity. Pushforward by
the continuous scaling map preserves weak convergence, so the assumed
full limit yields $t_\star\mu=\mu$.
This argument would not identify the limits of two independently chosen
subnets, which is why the hypothesis concerns the whole family.

\cref{itm:ch46-homsupp:supp}.  Let $F\subset X^{\An}$ be the fixed-point
set of the action, which is closed because the action is continuous
(\Cref{prop:ch32-scaling-action}) and which equals the set of trivial
semivaluations $v_{Y,\triv}$ by \eqref{eq:trivial-endpoints}.  Fix a
nonzero coherent ideal $\mathfrak b\subset\cO_X$ and let
$g:=\max\{\log|\mathfrak b|,-1\}\in\PL^+(X)$, the function
\eqref{eq:ch34-phi-Z}, so that $-1\le g\le0$ and, for $t>0$,
\[
 g(t\,v)=\max\{-t\,v(\mathfrak b),-1\}.
\]
By invariance, $\int g\,\rd\mu=\int g(t\,v)\,\rd\mu(v)$ for every
$t>0$.  Letting $t\to+\infty$ and using dominated convergence, with the
constant $1$ as dominating function and $\mu$ finite,
\[
 \int g\,\rd\mu=-\mu\bigl(\{v:v(\mathfrak b)>0\}\bigr).
\]
Letting instead $t\to0^+$ gives
$\int g\,\rd\mu=-\mu\bigl(\{v:v(\mathfrak b)=+\infty\}\bigr)$.
Subtracting,
\[
 \mu\bigl(\{v:0<v(\mathfrak b)<+\infty\}\bigr)=0 .
\]
Now let $K\subset X^{\An}\setminus F$ be compact.  Each $v\in K$ is not
a trivial semivaluation, so there are a nonzero ideal $\mathfrak b$ and
reals $0<\varepsilon<M$ with $\varepsilon<v(\mathfrak b)<M$; the set
$\{-M<\log|\mathfrak b|<-\varepsilon\}$ is open, contains $v$, and is
contained in $\{0<v(\mathfrak b)<\infty\}$, hence is $\mu$-null.
Finitely many such open sets cover $K$, so $\mu(K)=0$.  As $\mu$ is
Radon and $X^{\An}\setminus F$ is open, $\mu(X^{\An}\setminus F)=0$.
\end{proof}

\begin{example}[A homogeneous curve profile and its truncation limit]
\label{ex:ch46-homogeneous-curve-bdivisor}
Keep $X=\mathbf P^1$, $\omega=c_1(\cO(1))$, and $p\in X(k)$.
For $a\in(0,1]$ put $\psi_a=a\log|\mathfrak m_p|$, and set
$\psi_0=0$ separately. The profile $-at$ is homogeneous and
$\omega$-psh by the degree-one slope budget.
For $a>0$, the bounded truncations
$\psi_a^{(R)}=\max\{\psi_a,-aR\}$ satisfy
\[
 \MA^{\mathrm u}_\omega(\psi_a^{(R)})
    =(1-a)\delta_{v_{\triv}}+a\delta_{R\ord_p}.
\]
They have the explicit full weak limit
\[
 \mu_a=(1-a)\delta_{v_{\triv}}+a\delta_{v_{p,\triv}},
\]
with the same formula for $a=0$. It is the graph measure of this
homogeneous profile and is carried by fixed points. For $a>0$, its
potential has infinite energy; the notation $\mu_a$ does not assert a
value of the finite-energy operator at $\psi_a$.

The Cartier coefficient dictionary \Cref{prop:cartier-homogeneous-pl}
associates $\psi_a$ with $B=-a\,\overline{[p]}$, determined on $X$.
The class $\alpha=\overline{\omega+[B]}$ has degree $1-a$, and hence
\[
 \mu_a(\{v_{X,\triv}\})=1-a=\deg\alpha_X=(\alpha)_{\mathfrak X}.
\]
The remaining mass $a$ is at the proper-support fixed point.
This is an explicit curve identity. It is not an application of the
normalized nef correspondence, whose homogeneous representative vanishes
on all prime divisors of $X$; here $\psi_a(\ord_p)=-a$.
Nor does this computation establish the analogous mass identity for
arbitrary nef $b$-classes in higher dimension.
\end{example}

\subsection{The \texorpdfstring{$b$}{b}-divisor bridge}
\label{sec:ch46-bdivisor-bridge}

For the $b$-divisor comparison, work on the normal variety $X^\nu$.
The coefficient dictionary and the finite-energy measure construction
have different domains:
\[
 \begin{aligned}
  \text{Cartier $b$-divisors}
    &\longleftrightarrow\text{homogeneous divisorial coefficients},\\
  \text{bounded psh truncations}
    &\longrightarrow\text{mixed Radon measures}.
 \end{aligned}
\tag{46.28}\label{eq:bj-energy-homogeneous-dictionary}
\]
The first row is \Cref{prop:cartier-homogeneous-pl}; its homogeneous PL
functions live on the divisorial cone and need not be bounded PL functions
on the whole analytification. For nef non-Cartier $b$-classes, the
correspondence in \Cref{thm:bj-nef-bdivisor-correspondence} retains its
smooth characteristic-zero hypotheses and its normalization along prime
divisors. The second row is \Cref{thm:bj-mixed-ma} applied to bounded
truncations in a nef background.

The two rows are linked by truncating a homogeneous psh coefficient
function. The resulting measures are defined at every finite truncation
level. A full weak limit, if it exists, is scaling invariant by
\Cref{prop:ch46-homogeneous-support}. The energy may tend to $-\infty$,
so the finite-energy continuity theorem does not construct that limit.
\Cref{ex:ch46-homogeneous-curve-bdivisor} computes both a truncation limit
and a $b$-degree on a curve. No general bijection with invariant measures,
or general identity between the mass at $v_{X,\triv}$ and a nef
$b$-intersection number, is proved by this chapter.

\begin{warning}{Energy is not a $b$-intersection number}
\label{warn:ch46-energy-not-bdivisor}
A $b$-divisor on an $n$-dimensional variety has $n$-fold intersection
numbers under the hypotheses of \Cref{app:intersection-bdivisors}; there
is no $(n+1)$-fold product on those $n$-dimensional determinations to
identify with the compactified energy pairing. In addition, a homogeneous
psh potential has energy $0$ or $-\infty$ by
\Cref{cor:ch46-homogeneous-infinite-energy}. The nonzero case has infinite
energy on curves, and the domination argument in
\Cref{rem:ch46-homogeneous-higher} gives the same conclusion in general.
For the curve profile above, the $b$-degree is $1-a$ while the energies
of its truncations are $-Ra^2/2$. These are different invariants on
different spaces, even before the infinite-energy limit is taken.
\end{warning}

\begin{remark}[Legendre slices and where energy is defined]
\label{rem:ch46-legendre}
\Cref{thm:homogeneous-legendre-decomposition} reconstructs an arbitrary
psh function $\phi$ from its homogeneous slices
$\widehat\phi_\lambda$ through \eqref{eq:legendre-reconstruction-bj}.
The slices are generally unbounded; the theorem does not make them bounded
finite-energy functions. Their bounded truncations do have finite energy
in an ample background, but energy need not remain finite in the limit.

If $\phi\in\cE^1(\omega)$, its scaling orbit stays in that class and
$E_\omega(t\cdot\phi)=tE_\omega(\phi)$. This is the energy statement
available for the Legendre ray of \Cref{sec:ch35-legendre}.
The moment formula \Cref{prop:ch46-energy-dh} concerns the probability
Duistermaat--Heckman measure of a relatively ample test configuration.
It is proved by its weight/intersection identity, not by assigning finite
energies to homogeneous slices.
\end{remark}

\section{Worked examples}
\label{sec:ch46-examples}

Three computations close the chapter: a toric surface example in which
energy and mixed measures are obtained from the polytope and independently
from intersection theory; a comparison with \Cref{part:bfj} in which the
same number is produced twice, once over a trivially valued base and once
over a discretely valued one; and an example showing that energy is not
continuous for the weak topology, which is the reason
\Cref{chap:bj-strong-topologies} exists.

\subsection{A toric computation on \texorpdfstring{$\mathbf P^2$}{P2}}
\label{sec:ch46-toric}

\begin{example}[Energy and mixed measures on $\mathbf P^2$, three ways]
\label{ex:ch46-toric-p2}
Let $X=\mathbf P^2_k$ with homogeneous coordinates $z_0,z_1,z_2$, let
$T=\mathbf G_m^2$ act in the standard way with character lattice $M$
and cocharacter lattice $N$, and let $0:=[1:0:0]$ be the torus-fixed
point of the chart $\Spec k[x,y]$, $x=z_1/z_0$, $y=z_2/z_0$.  Take
$\omega=c_1(\cO(1))$, so $n=2$ and $V=(\omega^2)=1$.

\emph{The potential.}  Let
$\cX:=\operatorname{Bl}_{(0,0)}(\mathbf P^2\times\mathbf A^1)$ with
exceptional divisor $E\cong\mathbf P^2$, a smooth test configuration
whose central fibre is $\cX_0=\widetilde P+E$ with
$b_{\widetilde P}=b_E=1$, where $\widetilde P$ is the strict transform
of $\mathbf P^2\times\{0\}$.  As in \Cref{ex:ch34-normal-cone-segment} one
has $v_{\widetilde P}=v_{X,\triv}$ and
$v_E=\ord_0$, the multiplicity valuation at $0$, which in toric terms
is the monomial valuation $v_{(1,1)}$.  For $c\in[0,1]$ put
$D:=-cE$ and
\[
 \phi_c:=\phi_D=c\max\{\log|\mathfrak m_0|,-1\}
.
\]

\emph{Positivity range.}  On $\overline\cX
=\operatorname{Bl}_{(0,0)}(\mathbf P^2\times\mathbf P^1)$ write $H$ for
the pullback of the hyperplane class of $\mathbf P^2$, $F$ for the
fibre class $\overline\pi^*\{0\}$, and $E$ for the exceptional divisor.
The only nonzero triple products are
\[
 H^2F=1,\qquad E^3=1,
\]
all products involving both $E$ and one of $H,F$ vanishing because
$E$ is contracted by both projections; and $E^3=1$ because
$\cO_E(E)=\cO_{\mathbf P^2}(-1)$.  With $\cL=H-cE$, the curves in the
central fibre are the lines of $E$, the exceptional curve
$e=E\cap\widetilde P$ of $\widetilde P\cong\operatorname{Bl}_0\mathbf P^2$,
and the strict transforms of curves of $\mathbf P^2$; testing on a line
of $E$ gives $c\ge0$, and on the strict transform of a line through
$0$ gives $1-c\ge0$.  Conversely, for $0\le c\le1$ its restriction to $E$ is
$c\,\cO_E(1)$ and its restriction to $\widetilde P$ is
$(1-c)H+c(H-e)$, a nonnegative combination of nef classes.
All other fibres carry $\cO(1)$. Thus $\phi_c$ is $\omega$-psh
exactly for $0\le c\le1$.

\emph{First computation: intersection theory.}
\[
 (\omega,\phi_c)^3=(H-cE)^3
 =H^3-3cH^2E+3c^2HE^2-c^3E^3=-c^3,
\]
\[
 E_\omega(\phi_c)=\frac{(\omega,\phi_c)^3}{(n+1)V}=-\frac{c^3}{3}.
\]
For the measures, use $\widetilde P=F-E$ and
\Cref{prop:ch46-pl-measure}:
\[
 (H-cE)^2\cdot\widetilde P=1-c^2,
 \qquad
 (H-cE)^2\cdot E=c^2,
\]
so that
\[
 \MA^{\mathrm u}_\omega(\phi_c)
 =(1-c^2)\,\delta_{v_{\triv}}+c^2\,\delta_{\ord_0},
\]
of total mass $1=V$ and positive throughout the psh range $c\in[0,1]$. The mixed
measure of two such potentials is
\[
 (H-cE)(H-c'E)\cdot\widetilde P=1-cc',
 \qquad
 (H-cE)(H-c'E)\cdot E=cc',
\]
that is
\[
 (\omega+\ddc\phi_c)\wedge(\omega+\ddc\phi_{c'})
 =(1-cc')\,\delta_{v_{\triv}}+cc'\,\delta_{\ord_0},
\]
visibly symmetric in $c,c'$, positive for $c,c'\in[0,1]$, of mass
$1$, and reducing to the previous display when $c=c'$.

\emph{Second computation: from the polytope.}  With the conventions of
\Cref{thm:bj-toric-psh-convex} --- support function
$f_P(u)=\sup_{m\in P_{\mathrm{BJ}}}\langle u,m\rangle$, convex function
$f_\phi=(\phi+f_P)|_{N_{\mathbf R}}$, and
$P_{\mathrm{BJ}}=-P$ --- the polytope of $\cO(1)$ is
$P=\{m_1,m_2\ge0,\ m_1+m_2\le1\}$, so
$V_P=\operatorname{vol}(P_{\mathrm{BJ}})=\tfrac12=(\omega^2)/2!$ and
\[
 f_P(u)=\max\{0,-u_1,-u_2\}.
\]
For a monomial valuation $v_u$, $\log|\mathfrak m_0|(v_u)
=-\max\{0,\min(u_1,u_2)\}$, the outer maximum accounting for the fact that
$\mathfrak m_0$ is the unit ideal at the centre of $v_u$ when
$u\notin\sigma_0=\langle e_1,e_2\rangle$.  Hence
\[
 \phi_c(v_u)=-c\min\bigl\{1,\max\{0,\min(u_1,u_2)\}\bigr\},
 \qquad
 f_{\phi_c}=f_P+\phi_c .
\]
Its Legendre transform on $P_{\mathrm{BJ}}$ is computed by splitting the
supremum: on $\sigma_0$ the support function vanishes and the optimum has
$u_1=u_2=s\in[0,1]$, while off $\sigma_0$ the supremum is attained at
$u=0$; writing $m=-p$ with $p\in P$,
\[
 f_{\phi_c}^\vee(m)
 =\sup_{u}\bigl(\langle u,m\rangle-f_{\phi_c}(u)\bigr)
 =\max\{0,\ c-p_1-p_2\}
 =(c+m_1+m_2)_+ .
\]
Now \eqref{eq:bj-toric-energy} gives
\[
 E_\omega(\phi_c)
 =-\frac1{V_P}\int_{P_{\mathrm{BJ}}}f_{\phi_c}^\vee\,\rd m
 =-2\int_P(c-p_1-p_2)_+\,\rd p
 =-2\int_0^{c}(c-t)\,t\,\rd t
 =-\frac{c^3}{3},
\]
where $\int_Ph(p_1+p_2)\rd p=\int_0^1h(t)\,t\,\rd t$ because the area of
$\{p_1+p_2\le t\}\cap P$ is $t^2/2$.  This is the number found by
intersection theory.

For the measure, \eqref{eq:bj-toric-ma} says
$\MA^{\mathrm u}_\omega(\phi_c)=n!\,\iota_*\operatorname{MA}_{\mathbf R}
(f_{\phi_c})$, and the Alexandrov measure of a convex function is the
pushforward of Lebesgue measure under the gradient of its conjugate.  Here
$\nabla f_{\phi_c}^\vee=(1,1)$ on $\{m_1+m_2>-c\}$ and $=(0,0)$ on
$\{m_1+m_2<-c\}$, so
\[
\begin{aligned}
 \operatorname{MA}_{\mathbf R}(f_{\phi_c})
 &=\operatorname{vol}\{p\in P:p_1+p_2<c\}\ \delta_{(1,1)}
  +\operatorname{vol}\{p\in P:p_1+p_2>c\}\ \delta_{(0,0)}\\
 &=\frac{c^2}{2}\delta_{(1,1)}+\frac{1-c^2}{2}\delta_{(0,0)} .
\end{aligned}
\]
Multiplying by $n!=2$ and using $v_{(1,1)}=\ord_0$,
$v_{(0,0)}=v_{X,\triv}$ recovers the measure
\[
 \MA^{\mathrm u}_\omega(\phi_c)
  =c^2\,\delta_{\ord_0}+(1-c^2)\,\delta_{v_{\triv}}
\]
computed above by intersection theory.

\emph{Third computation: polarization on the polytope.}  The mixed measure
can also be read off convexly.  The conjugate of $f_{\phi_c}+f_{\phi_{c'}}$
is the infimal convolution
$f_{\phi_c}^\vee\square f_{\phi_{c'}}^\vee$, which on $2P_{\mathrm{BJ}}$
equals $(m_1+m_2+c+c')_+$. To see this, put $m=-p$ and
$s=p_1+p_2$. The two summands of $p$ have total sizes
$s_1,s_2\in[0,1]$ with $s_1+s_2=s$; every such split is feasible by
splitting $p$ proportionally. Minimizing
$(c-s_1)_++(c'-s_2)_+$ gives $(c+c'-s)_+$, since $c,c'\le1$.
Its gradient is $(1,1)$ on a set of volume
$(c+c')^2/2$.  Hence
\[
 \operatorname{MA}_{\mathbf R}(f_{\phi_c}+f_{\phi_{c'}})
 =\frac{(c+c')^2}{2}\delta_{(1,1)}+\Bigl(2-\frac{(c+c')^2}{2}\Bigr)
   \delta_{(0,0)},
\]
and the polarization identity
$\operatorname{MA}_{\mathbf R}(f+g)=\operatorname{MA}_{\mathbf R}(f)
+2\operatorname{MA}_{\mathbf R}(f,g)+\operatorname{MA}_{\mathbf R}(g)$
gives $\operatorname{MA}_{\mathbf R}(f_{\phi_c},f_{\phi_{c'}})
=\tfrac12\bigl[\tfrac{(c+c')^2}{2}-\tfrac{c^2}{2}-\tfrac{c'^2}{2}\bigr]
\delta_{(1,1)}+\cdots=\tfrac{cc'}{2}\delta_{(1,1)}+\cdots$.  Multiplying
by $2$ gives the atom $cc'$ at $\ord_0$, which is the mixed
intersection number computed above.

\emph{Unit tests.}  Setting $c=0$ gives $E_\omega=0$ and
$\MA^{\mathrm u}=\delta_{v_{\triv}}$, which is
\eqref{eq:bj-trivial-energy-zero} and the trivial-metric measure.  Setting
$c=1$ gives $E_\omega=-\tfrac13$ and
$\MA^{\mathrm u}=\delta_{\ord_0}$: all the mass has moved to the single
divisorial point, which is the boundary of the psh range.  Adding a
constant $e$ multiplies nothing and adds $e$, as
\eqref{eq:bj-energy-basic-identities} requires; on the polytope side it
translates $f_{\phi}^\vee$ by $-e$ and changes the integral by
$eV_P$, giving the same $+e$.  Finally the scaling action.  The support
function $f_P$ is positively homogeneous of degree one, so
\eqref{eq:function-scaling-action} gives
\[
 f_{t\cdot\phi}(u)=f_P(u)+t\,\phi(v_{u/t})
 =t\bigl[f_P(u/t)+\phi(v_{u/t})\bigr]
 =t\,f_\phi(u/t),
 \qquad
 f_{t\cdot\phi}^\vee=t\,f_\phi^\vee,
\]
the second identity by the substitution $u=tw$ in the definition of the
Legendre transform.  Formula \eqref{eq:bj-toric-energy} then gives
$E_\omega(t\cdot\phi)=t\,E_\omega(\phi)$ at once, which is a second, and
purely convex-analytic, proof of \Cref{thm:ch46-energy-homogeneous} in the
toric case.  Concretely
$(t\cdot\phi_c)(v_u)=-c\min\{t,\max\{0,\min(u_1,u_2)\}\}$, the same
potential with the truncation level $1$ replaced by $t$, and
$E_\omega(t\cdot\phi_c)=-tc^3/3$.
\end{example}

\begin{remark}[Comparison with the toric computation of \Cref{part:bfj}]
\label{rem:ch46-toric-part3}
\Cref{thm:ch26-toric-energy} proves the toric energy formula over a
discretely valued base, in the form
$E(F)-E(F_0)=n!\int_P(g_F-g_{F_0})\rd\lambda$ with $g_F$ the roof
function of \Cref{def:ch26-roof}, and
\Cref{ex:ch26-toric-surface} computes it on a toric surface; the mixed
measure is \Cref{ex:ch24-toric-surface-mixed}.  The two formulae are the
same up to the normalization of \Cref{warn:ch46-normalization} and the
reflection $P_{\mathrm{BJ}}=-P$ of
\eqref{eq:bj-toric-sign-conversion}: the roof function $g_F$ and the
Legendre transform $f^\vee$ differ by that reflection and a sign, which
is why one formula carries a plus and the other a minus.  A reader
comparing the two should carry out the substitution explicitly; mixing them
produces an energy of the wrong sign, which is detected immediately by the
trivial metric, for which one formula must give $0$.
\end{remark}

\subsection{The same number twice}
\label{sec:ch46-same-number}

\begin{example}[One energy, computed over a trivial and over a discrete
                valuation]
\label{ex:ch46-same-number-twice}
Let $d\in\mathbf Z_{>0}$ and $c\in[0,d]$.

\emph{Over a trivially valued field.}  This is
\Cref{ex:ch46-p1-energy}: $X=\mathbf P^1_k$, $\omega=c_1(\cO(d))$,
$V=d$; the test configuration
$\cY=\operatorname{Bl}_{(p,0)}(\mathbf P^1\times\mathbf A^1)$ with
$\cY_0=\widetilde P+E$; the vertical divisor $D=-cE$; and
\[
 \begin{aligned}
  (\overline\cL^{\,2})&=(dH-cE)^2=-c^2,\\
  \MA^{\mathrm u}_\omega(\phi_D)
    &=(d-c)\delta_{v_{\triv}}+c\,\delta_{\ord_p},\\
  E_\omega(\phi_D)&=-\frac{c^2}{2d}.
 \end{aligned}
\]

\emph{Over a discretely valued field.}  Let $R=k[[\varpi]]$ and $K=k((\varpi))$, with the complete discrete
valuation normalized by $v_K(\varpi)=1$,
and let $X_K=\mathbf P^1_K$ with $L=\cO(d)$, so that $V=(L)=d$.  Let
$\cX:=\operatorname{Bl}_z\mathbf P^1_R$ be the blowup at a closed point
$z$ of the special fibre, with exceptional curve $E$ and strict
transform $\widetilde P$ of $\mathbf P^1_k$; both have multiplicity
$1$, so the dual graph is a segment of length one with vertices
$x_0=v_{\widetilde P}$, the Gauss point, and $x_1=v_E$.  The trivial
model metric has
$\omega=\MA(\varphi_{\mathrm{triv}})=d\,\delta_{x_0}$, since
$\mu^*\cO(d)$ has degree $d$ on $\widetilde P$ and degree $0$ on
$E$; in the notation of \Cref{ex:ch26-normalization-audit}, $m_0=d$ and
$m_1=0$.  Take $\varphi:=\varphi_{-cE}$, so that
$a=\varphi(x_0)=0$ and $b=\varphi(x_1)=-c$.  Then
\eqref{eq:ch26-two-vertex-energy} gives
\[
 E^{\mathrm{BFJ}}_\omega(\varphi)-E^{\mathrm{BFJ}}_\omega(0)
 =m_0a+m_1b-\tfrac12(a-b)^2
 =-\frac{c^2}{2},
\]
and \eqref{eq:model-energy}, read through the vertical-factor pairing
\eqref{eq:ch20-pairing-as-integral}, gives
\[
 2\bigl(E^{\mathrm{BFJ}}_\omega(\varphi)
        -E^{\mathrm{BFJ}}_\omega(0)\bigr)
 =\bigl((-cE)\cdot\omega_\varphi\bigr)
  +\bigl((-cE)\cdot\omega\bigr)
 =-c^2 .
\]
Here the first slot in each term is the actual vertical divisor $-cE$;
the last equality uses $E^2=-1$ and $(\omega\cdot E)=0$.
The Monge--Amp\`ere measure is
$\omega+\ddc\varphi=(d-c)\delta_{x_0}+c\,\delta_{x_1}$, by the
two-vertex Laplacian of \Cref{ex:ch24-good-reduction} and
\Cref{ex:ch26-normalization-audit}.

\emph{The comparison.}  The three displayed objects match term by term:
the compactified self-intersection on the trivially valued side and the
mixed vertical sum on the discretely valued side are both $-c^2$, the
measures are both
$(d-c)\delta_{\text{Gauss}}+c\,\delta_{\text{exceptional}}$, and the
energies agree after the conversion of \Cref{warn:ch46-normalization},
\[
 E_\omega(\phi_D)=\frac1V
 \bigl(E^{\mathrm{BFJ}}_\omega(\varphi)-E^{\mathrm{BFJ}}_\omega(0)\bigr)
 =-\frac{c^2}{2d}.
\]
This is not a coincidence and the reason is
\Cref{rem:ch46-fibre-class}: in both computations the ambient scheme is a
regular relative surface over its respective one-dimensional base,
the divisor $-cE$ is
supported in one fibre, the class $\theta$ restricts to $\cO(d)$ on the
general fibre, and the only intersection numbers used are
$E^2=-1$, $\theta\cdot E=0$, $\theta\cdot(\text{fibre})=d$.  Those
three numbers are the same on $\operatorname{Bl}_{(p,0)}
(\mathbf P^1\times\mathbf P^1)$ and on $\operatorname{Bl}_z\mathbf P^1_R$.

\emph{Where the two computations diverge.}  Three places, and they are the
content of the last two rows of the table in \Cref{sec:ch46-part3}.  First,
over the discretely valued base the special fibre is the divisor of a
uniformizer, so $\varphi+e$ and $\varphi$ differ by $e\,\cX_s$; over
the trivially valued base the fibre class $\overline\pi^*\{0\}$ is not
principal, and the corresponding statement is
\eqref{eq:ch33-intersection-twist}.  The two nevertheless produce the same
translation rule, because \Cref{sec:ch20-pairing} gives $\cX_s$ the same
pairing with the horizontal classes.  Second, the trivially valued side
carries the scaling action: $t\cdot\phi_D$ has energy $-tc^2/(2d)$ by
\Cref{ex:ch46-homogeneity-p1}, and there is no operation on the discretely
valued side that does this --- degree-$d$ base change exists there too,
but it changes the ground field, whereas here it is an automorphism of the
situation.  Third, the displayed two-vertex model calculation describes one finite
segment. The mass-escape example instead keeps the initial slope fixed
and lets the truncation length tend to infinity through the scaling
parameter $t$. Its energies tend to $-\infty$ and its atoms approach the
proper-support endpoint. The parameter $c$ remains in $[0,d]$ throughout;
it is not sent to infinity. Proper-support points also exist over a
nontrivially valued field, so their existence alone does not distinguish
the two theories.
\end{example}

\subsection{Energy is not weakly continuous}
\label{sec:ch46-not-continuous}

\begin{example}[A weakly convergent sequence whose energies do not converge]
\label{ex:ch46-energy-not-weakly-continuous}
Let $X=\mathbf P^1_k$ and $\omega=c_1(\cO(1))$, so $V=1$.  Since
$k$ is algebraically closed it is infinite; choose pairwise distinct
closed points $p_1,p_2,\dots\in X(k)$ and set
\[
 \phi_j:=\max\{\log|\mathfrak m_{p_j}|,-1\}\in\PL^{\omega}(X),
 \qquad j\ge1 .
\]
By \Cref{ex:ch46-p1-energy} with $d=1$ and $c=1$,
\[
 E_\omega(\phi_j)=-\tfrac12,
 \qquad
 \MA^{\mathrm u}_\omega(\phi_j)=\delta_{\ord_{p_j}} .
\]

\emph{Weak convergence of the potentials.}  Let $v\in X^{\An}$.  If
$v$ lies on the ray of a closed point $q$, or is a proper-support point
$v_{q,\triv}$, then $\phi_j(v)=0$ unless $q=p_j$; since the $p_j$
are distinct, this happens for at most one $j$.  Hence
$\phi_j\to0$ pointwise on all of $X^{\An}$, in particular on
$X^{\divis}$; so $\phi_j\to0$ weakly in $\PSH(\omega)$.

\emph{Weak convergence of the measures.}  Let
$f\in\PL_{\mathbf R}(X)$.  Writing $f$ as a real combination of
flag-ideal functions $\phi_{\mathfrak a}
=\max_\lambda\{\log|\mathfrak a_\lambda|+\lambda\}$ and using that each
nonzero $\mathfrak a_\lambda$ has finite zero locus on the curve $X$,
one has $f(\ord_q)=f(v_{\triv})$ for all but finitely many closed points
$q$.  Hence $\int f\,\rd\MA^{\mathrm u}_\omega(\phi_j)=f(\ord_{p_j})
=f(v_{\triv})=\int f\,\rd\MA^{\mathrm u}_\omega(0)$ for $j$ large, and by
\Cref{thm:ch34-pl-density} and the uniform mass bound,
$\MA^{\mathrm u}_\omega(\phi_j)\to\MA^{\mathrm u}_\omega(0)
=\delta_{v_{\triv}}$ weakly.

\emph{The conclusion.}  Both the potentials and their Monge--Amp\`ere
measures converge weakly to those of $\phi=0$, and yet
\[
 \lim_jE_\omega(\phi_j)=-\tfrac12\ \ne\ 0=E_\omega(0).
\]
So $E_\omega$ is not continuous for the weak topology on
$\PSH(\omega)$, and it is not continuous on $\cE^1(\omega)$ either,
every $\phi_j$ being bounded.  Upper semicontinuity does hold, as
\Cref{prop:ch46-energy-usc} asserts: $-\tfrac12\le0$, and the inequality
is strict.
\end{example}

\begin{remark}[What the example rules out, and what it motivates]
\label{rem:ch46-strong-topology-motivation}
The example rules out three plausible-sounding statements.  It is
\emph{not} true that weak convergence of psh functions implies convergence
of energies; it is not true that weak convergence of potentials together
with weak convergence of their Monge--Amp\`ere measures implies convergence
of energies; and it is not true that $E_\omega$ is continuous on the
bounded part of $\cE^1(\omega)$ for the weak topology, even along
sequences.  The failure is not exotic: the sequence $(\phi_j)$ consists
of the simplest PL psh functions on $\mathbf P^1$, all with the same
energy.

What it motivates is the definition of the \emph{strong} topology on
$\cE^1(\omega)$ in \Cref{chap:bj-strong-topologies}: the coarsest
refinement of the weak topology for which $E_\omega$ is continuous.  In
that topology the sequence above does not converge, which is the intended
behaviour, and the theory recovers the continuity statements that the weak
topology cannot support.  The corresponding phenomenon over a discretely
valued base is \Cref{ex:ch26-usc-strict} and
\Cref{prop:ch26-no-strong-compactness}, and the strong topology introduced
in \Cref{sec:ch26-strong} solves the same problem there; the reader will
find that the two constructions are formally identical and that the present
example transports verbatim, with the points $p_j$ replaced by closed
points of the special fibre.

One further caution.  \Cref{prop:ch46-energy-usc} shows that $E_\omega$
is upper semicontinuous, and \Cref{cor:ch45-compactness} that
$\PSH^{\sup}(\omega)$ is weakly compact.  Together these give existence
of maximizers of $E_\omega$ under a weak constraint --- the variational
method --- and this is not affected by the discontinuity: an upper
semicontinuous function on a compact space attains its supremum whether or
not it is continuous.  The discontinuity is an obstruction to identifying
limits, not to producing them.
\end{remark}

\section{Scope and hypotheses}
\label{sec:ch46-scope}

\subsection{Where each hypothesis is used}
\label{sec:ch46-hypothesis-use}

\begin{center}
\begin{tabular}{p{.30\textwidth}p{.62\textwidth}}
\toprule
hypothesis & where it is used\\
\midrule
$X$ integral projective &
throughout; projectivity for the compactification
\Cref{thm:ch33-compactification} and for the existence of degrees,
integrality for the identification of the pairing with an intersection
number on the integral $\overline\cX$
(\Cref{thm:ch33-intersection-invariance}\,(1))\\
$k$ algebraically closed and trivially valued &
in the correspondence between test configurations and PL functions
(\Cref{thm:tc-pl-correspondence}), and in every statement about the
scaling action\\
$X$ unibranch &
only through the compactness and obstacle-envelope parts of
\Cref{imp:ch46-envelope-imported}; see
\Cref{sec:ch46-where-envelope} for the uses and
\Cref{rem:ch46-without-envelope} for what survives without it\\
$\omega$ ample &
in \eqref{eq:bj-volume} so that $V>0$ and the probability normalization
is legitimate; in \Cref{lem:ch46-competitors-nonempty} so that constants
are competitors; in \Cref{cor:ch45-continuity-envelopes}, hence in
\Cref{lem:ch46-continuous-competitors}\\
$\omega$ the class of a $\mathbf Q$-line bundle &
only when a member of $\cH_{\mathrm{dom}}(\omega)$ is to be interpreted
as a Fubini--Study function of one line bundle; neither
\Cref{lem:ch46-pl-domination} nor the extended pairing requires this
rationality; see \Cref{rem:ch46-real-classes}\\
$\theta_i$ nef and $\phi_i$ psh &
in the positivity clauses
\Cref{prop:ch46-pl-measure}\,\cref{itm:ch46-pl-ma:pos},
\Cref{lem:ch46-vertical-negativity}, and
\Cref{prop:bj-model-energy-monotonicity}; multilinearity removes the
hypothesis at the cost of signed measures\\
finite energy &
in every identity involving a difference: \eqref{eq:bj-energy-cocycle},
\eqref{eq:bj-integration-by-parts},
\eqref{eq:bj-e1-polarized-pairing}; see
\Cref{warn:ch46-no-subtraction}\\
\bottomrule
\end{tabular}
\end{center}

\subsection{What is not claimed}
\label{sec:ch46-not-claimed}

\begin{enumerate}
\item No statement here identifies the class $\PSH(\theta)$ with a class
      defined by a local differential-geometric condition.  The comparison
      with the Chambert--Loir--Ducros product, for continuous semipositive
      metrics, is recorded in \cite[Remark~7.23]{BJ22} and is taken up in
      \Cref{chap:compare-ma}; nothing in this chapter proves it.
\item No regularity of $P_\omega$ beyond continuity in an ample class is
      claimed, and none is used; \Cref{ex:ch45-no-regularity} shows that
      higher regularity is not guaranteed in general.
\item Joint upper semicontinuity uses the finite-set majorant comparison
      of \cite[Lemma~4.26]{BJ22}; it is not inferred from separate upper
      semicontinuity or from concavity of the mixed pairing
      (\Cref{rem:ch46-usc-scope}).
\item The domination principle, the uniqueness of solutions of the
      Monge--Amp\`ere equation, and orthogonality are deferred to
      \Cref{chap:bj-finite-energy} (\Cref{rem:ch46-not-proved-locality}).
\item The signed pairing and mixed measures are defined for arbitrary
      background classes with potentials in $\vec\cE^{\,1}(X)$, including
      zero-volume backgrounds. The probability operator $\MA_\omega$,
      normalized energy $E_\omega$, and class $\cE^1(\omega)$ are used
      here only for ample $\omega$. The extremal function $V_\theta$ by
      itself does not define a finite-energy theory in a merely
      pseudoeffective class.
\end{enumerate}

\section{Chapter dictionary}
\label{sec:ch46-dictionary}

\begingroup
\begin{longtable}{@{}p{.23\textwidth}p{.72\textwidth}@{}}
\toprule
object & statement and reference\\
\midrule
\endfirsthead
\toprule
object & statement and reference\\
\midrule
\endhead
PL pairing &
$(\theta_0,\phi_0)\cdots(\theta_n,\phi_n)
=(\cL_0\cdots\cL_n)_{\overline\cX}$; symmetric and multilinear,
monotone in a metric slot when the other slots are psh; Lipschitz with
the mixed intersection mass as coefficient
\par\smallskip \Cref{def:bj-pl-energy-pairing},
\Cref{prop:ch46-pairing-wellposed},
\Cref{prop:bj-model-energy-monotonicity},
\Cref{cor:ch46-pl-lipschitz}\\
vertical negativity &
$(F^2\cdot\cL_1\cdots\cL_{n-1})\le0$ for $F$ in one fibre and
$\cL_i$ relatively nef
\par\smallskip \Cref{lem:ch46-vertical-negativity}; compare
\Cref{thm:ch20-mixed-hodge}\\
model measure &
$\sum_Eb_E(\cL_1|_E\cdots\cL_n|_E)\delta_{v_E}$; mass
$(\theta_1\cdots\theta_n)$; positive for nef psh data
\par\smallskip \Cref{prop:ch46-pl-measure}, \eqref{eq:bj-model-ma-atomic}\\
energy &
$E_\omega(\phi)=(\omega,\phi)^{n+1}/((n+1)V)$; normalized so that
$E(\phi+c)=E(\phi)+c$
\par\smallskip \Cref{def:bj-ma-energy}, \Cref{warn:ch46-normalization}\\
first moment &
$E_\omega(\phi_{\cL})=\int\lambda\,\rd\DHm(\mathcal F)$
\par\smallskip \Cref{prop:ch46-energy-dh}\\
extension &
$\inf$ over PL psh competitors above $\phi$; unique with the four
listed properties
\par\smallskip \Cref{thm:bj-extended-energy-pairing},
\Cref{lem:ch46-pl-domination}, \Cref{lem:ch46-dini}\\
usc and concavity &
$E_\omega=\inf_\psi[E_\omega(\psi)+\int(\phi-\psi)\MA_\omega(\psi)]$,
an infimum of weakly continuous affine functions
\par\smallskip \Cref{prop:ch46-energy-usc}\\
finite energy &
$\cE^1(\omega)=\{E_\omega>-\infty\}$; convex and upward closed,
containing bounded psh functions; the curve examples make both inclusions
strict
\par\smallskip \Cref{def:bj-e1}, \Cref{thm:bj-e1-basic},
\Cref{ex:ch46-e1-curve}\\
mixed measures &
Riesz representation of $f\mapsto(0,f)\prod(\theta_i,\phi_i)$; symmetric,
mass $(\theta_1\cdots\theta_n)$; positive for nef psh data
\par\smallskip \Cref{thm:bj-mixed-ma}\\
integration by parts &
symmetry of one intersection number, extended by monotone convergence
\par\smallskip \Cref{cor:ch46-pl-ibp},
\Cref{prop:bj-integration-by-parts},
\Cref{cor:ch46-dirichlet-e1}\\
locality and comparison &
$\mathbf1_{\{\phi>\psi\}}\MA(\max)=\mathbf1_{\{\phi>\psi\}}\MA(\phi)$;
$\int_{\{\phi<\psi\}}\MA(\psi)\le\int_{\{\phi<\psi\}}\MA(\phi)$
\par\smallskip \Cref{thm:bj-locality-comparison},
\Cref{rem:ch46-locality-obstruction}\\
homogeneity &
$E_\omega(t\cdot\phi)=tE_\omega(\phi)$; mixed measures are
$t_\star$-equivariant
\par\smallskip \Cref{thm:ch46-energy-homogeneous},
\eqref{eq:bj-homogeneous-ma-scaling}\\
homogeneous potentials &
energy $0$ or $-\infty$, and $-\infty$ unless zero on a curve;
unbounded unless zero; a full weak limit of their bounded-truncation
measures, when it exists, is invariant and carried by trivial semivaluations
\par\smallskip \Cref{cor:ch46-homogeneous-infinite-energy},
\Cref{prop:ch46-homogeneous-support}\\
$b$-divisor bridge &
Cartier $b$-divisors correspond to homogeneous divisorial coefficients
on the normal variety; bounded psh truncations have mixed measures;
no general bijection with invariant measures is asserted
\par\smallskip \eqref{eq:bj-energy-homogeneous-dictionary},
\Cref{warn:ch46-energy-not-bdivisor}\\
\bottomrule
\end{longtable}
\endgroup

Four separations should be carried forward.

\emph{Finitary versus analytic.}  Everything through
\Cref{sec:ch46-p1} is intersection theory on a projective
$(n+1)$-fold and holds for any integral projective $X$; everything
after it is analysis.  The monotone extension to singular potentials still
uses no envelope theorem; compactness, extremal functions, obstacle
calculus, and the resulting variational applications do
(\Cref{sec:ch46-where-envelope}).

\emph{Normalized versus unnormalized.}  $\MA_\omega$ is a probability
measure here and has mass $V$ in \Cref{part:bfj}; $E_\omega$ translates
by $c$ here and by $cV$ there.  Every displayed comparison between the
two parts carries the conversion explicitly
(\Cref{warn:ch46-normalization},
\Cref{ex:ch46-same-number-twice}).

\emph{Energy versus mass.}  The energy is an $(n+1)$-fold number and is
homogeneous of degree one for the scaling action; the mass of a mixed
measure is an $n$-fold number and is scaling invariant.  Confusing them
produces \Cref{warn:ch46-energy-not-bdivisor}.

\emph{Weak versus strong.}  $E_\omega$ is upper semicontinuous but not
continuous for the weak topology
(\Cref{prop:ch46-energy-usc},
\Cref{ex:ch46-energy-not-weakly-continuous}); the repair is
\Cref{chap:bj-strong-topologies}, not a sharper argument here.

\section{Source and provenance}
\label{sec:ch46-source}

\subsection{Source note}
\label{sec:ch46-source-note}

The PL pairing is \cite[Definition~3.12 and Proposition~3.14]{BJ22}; its
singular extension, energy, finite-energy calculus, and mixed measures are
\cite[Theorem~7.1, Definition~7.6, Theorems~7.8 and~7.18, and
Proposition~7.19]{BJ22}.  The agreement with the Chambert--Loir--Ducros
product for continuous semipositive metrics is recorded in
\cite[Remark~7.23]{BJ22}.  Locality and comparison are
\cite[Theorems~7.40 and~7.44]{BJ22}.  None of these constructions requires
the envelope property.

That last sentence has to be read carefully, and the care is the subject of
\Cref{sec:ch46-where-envelope}.  It is true of the \emph{constructions}:
the pairing, the measures, their formal properties, and the fixed-real-ample
PL regularization used to extend them are available without the envelope
property.  What the envelope theorem adds is
\Cref{cor:ch45-compactness}, the extremal functions and obstacle operators,
and hence the variational applications that require attainment.  The
correct summary is that monotone analytic extension predates the envelope
theorem, whereas compact variational theory does not.

The Duistermaat--Heckman interpretation \Cref{prop:ch46-energy-dh} rests on
\Cref{thm:ch33-dh-weight-count}\,(3), quoted there from
\cite[Corollary~3.4]{BHJ17}.  The toric formulae used in
\Cref{ex:ch46-toric-p2} are \Cref{thm:bj-toric-ma-energy}, quoted in
\Cref{chap:bj-linear-growth} from \cite[Appendix~B]{BJ22}; the reflection
\eqref{eq:bj-toric-sign-conversion} between the polytope conventions of
\Cref{chap:model-calculations} and of that appendix is applied there and
must be applied by anyone comparing the two.  The vertical negativity
inequality \Cref{lem:ch46-vertical-negativity} is
\cite[Lemma~3.24]{BJ22}; its counterpart over a discretely valued base is
\Cref{thm:ch20-mixed-hodge}, and the proof recorded there is the one
sketched in \Cref{lem:ch46-vertical-negativity}.

\subsection{What is supplied here and is not in the source}
\label{sec:ch46-additions}

The following material is not in \cite{BJ22} in this form and is supplied
here.

\Cref{prop:ch46-pairing-wellposed}, \Cref{lem:ch46-vertical-negativity},
\Cref{cor:ch46-pl-lipschitz}, \Cref{lem:ch46-telescope}: the formal
properties of the PL pairing, assembled from
\Cref{chap:bj-test-configurations} and
\Cref{app:intersection-bdivisors} so that the reader can verify each of
them on a model without consulting the source.  The Lipschitz estimate is
the one used four times below and is not stated separately in
\cite{BJ22}.

\Cref{prop:ch46-pl-measure}, \Cref{cor:ch46-pl-ibp},
\Cref{prop:ch46-dirichlet-vertical},
\Cref{prop:ch46-first-variation-pl},
\Cref{cor:ch46-pl-tangent}: the PL-level measure theory, integration by
parts, Dirichlet form and first variation.  These are the trivially valued
counterparts of \Cref{prop:ch20-intersection-pairing},
\Cref{cor:ch20-integration-by-parts} and
\Cref{prop:ch26-dirichlet-vertical}, and they are written out here because
the comparison with \Cref{part:bfj} is one of the stated purposes of this
chapter.

\Cref{prop:ch46-energy-dh}: the identification of the energy of a PL psh
metric with the first moment of its Duistermaat--Heckman probability
measure.  The ingredients are in \cite{BHJ17} and in
\Cref{chap:bj-test-configurations}; the statement in this form is supplied
here, and it is what fixes the normalization of
\Cref{warn:ch46-normalization}.

\Cref{lem:ch46-competitors-nonempty},
\Cref{lem:ch46-pl-domination},
\Cref{lem:ch46-dini},
\Cref{lem:ch46-continuous-competitors}: the three preparatory statements
for the extension, followed by the envelope-dependent directedness of the
larger continuous competitor family.  The short proof of
\Cref{lem:ch46-pl-domination} now exposes its exact input: the
fixed-real-ample regularization
\Cref{thm:ch34-fs-regularization} already supplies the required
dominating net in $\cH_{\mathrm{dom}}(\omega)$.

\Cref{prop:ch46-energy-usc}: the representation of $E_\omega$ as an
infimum of weakly continuous affine functions, from which concavity and
upper semicontinuity follow at once.  The properties are
\cite[Proposition~7.7]{BJ22}; the representation is supplied here and is
reused in \Cref{chap:bj-finite-energy}.

\Cref{thm:ch46-energy-homogeneous},
\Cref{cor:ch46-homogeneous-infinite-energy}, and
\Cref{prop:ch46-homogeneous-support} explain the scaling action.
Energy homogeneity is \cite[Proposition~7.7(iii)]{BJ22}; equivariance of
finite-energy mixed measures is \cite[Proposition~7.19(iii)]{BJ22}.
The direct proofs, homogeneous energy dichotomy, and fixed-point support
argument for an assumed full truncation limit are supplied here.
The latter statement does not extend the operator beyond finite energy.

\Cref{ex:ch46-p1-energy},
\Cref{ex:ch46-p1-two-corner},
\Cref{ex:ch46-e1-curve},
\Cref{ex:ch46-homogeneity-p1},
\Cref{ex:ch46-homogeneous-mass-escapes},
\Cref{ex:ch46-homogeneous-curve-bdivisor},
\Cref{ex:ch46-toric-p2},
\Cref{ex:ch46-same-number-twice},
\Cref{ex:ch46-energy-not-weakly-continuous}: worked presentations supplied here. Basic curve profiles also occur in
\cite[Examples~3.18--3.19 and~7.20]{BJ22}. They are included because
the two claims this chapter makes --- that the theory is intersection
theory, and that the normalizations are forced --- are exactly the claims
that a plausible reading can appear to establish without doing so.  A
reader who has computed $-c^2/(2d)$ three ways on $\mathbf P^1$,
$-c^3/3$ three ways on $\mathbf P^2$, and the same $-c^2$ over both a
trivial and a discrete valuation, will not mix the two normalizations by
accident.

\subsection{Historical note}
\label{sec:ch46-history}

The energy functional of this chapter has three ancestries, and the
trivially valued theory is where they meet.

The oldest is the Aubin--Mabuchi energy of Kähler geometry, the primitive
of the complex Monge--Amp\`ere operator, whose non-Archimedean shadow is
\eqref{eq:appJ-energy}.  Its characteristic properties --- monotone,
concave, with first variation the Monge--Amp\`ere measure --- are the ones
verified here in \Cref{prop:ch46-first-variation-pl} and
\Cref{cor:ch46-pl-tangent}, and they are what any candidate definition must
reproduce.

The second is the intersection-theoretic one.  Over a discretely valued
field this is the model-metric formalism of
\Cref{chap:bfj-model-forms,chap:bfj-energy}, going back to the arithmetic
intersection theory of Gillet--Soulé and Zhang and made into a
pluripotential theory by Boucksom--Favre--Jonsson.  On model functions its
intersection formula is the mixed vertical sum
\[
 E_\omega(\varphi)-E_\omega(\psi)
 =\frac1{n+1}\sum_{j=0}^n
  \bigl(D_{\varphi-\psi}\cdot
        \omega_\varphi^{\,j}\cdot\omega_\psi^{\,n-j}\bigr),
\]
where $D_{\varphi-\psi}$ is an actual vertical divisor; this is
\eqref{eq:model-energy} and \eqref{eq:ch20-pairing-as-integral} in
energy-difference form.  Over a trivially valued field the corresponding
number is the absolute intersection on a compactified test configuration
from \Cref{def:bj-pl-energy-pairing}.  The reader who wants a single
sentence for the relation between the two should take
\Cref{rem:ch46-fibre-class}: the fibre class over the origin plays the role
of the special fibre, and it has the same two properties.

The third is the stability-theoretic one.  For a test configuration the
number $(\overline\cL^{\,n+1})$ determines the leading coefficient
of the total-weight polynomial (\Cref{thm:ch33-dh-weight-count}), and
\Cref{prop:ch46-energy-dh} says that the energy is the first moment of the
Duistermaat--Heckman measure of the associated filtration.  That this
number, arrived at from geometric invariant theory, coincides with the
non-Archimedean Aubin--Mabuchi energy is the observation that made the
valuative approach to K-stability possible, and it is developed in
\Cref{chap:bj-stability}.

The finitary theory --- one test configuration, one intersection number ---
is the content of \Cref{sec:ch46-model,sec:ch46-pl-measure}.  Its extension
to singular potentials uses the unconditional fixed-class PL
regularization of \Cref{thm:ch34-fs-regularization}; this is the moment at
which an algebraic invariant becomes an analytic functional.  What
\Cref{chap:bj-envelope-consequences} subsequently supplies is the
compactness and envelope calculus that make the functional variationally
effective.  Keeping those two transitions separate is essential: the
monotone energy theory does not require unibranchness, while its compact
variational applications do.

A final remark on normalization, since it is the commonest source of error
in this subject.  Three conventions are in circulation: the unnormalized
energy of \eqref{eq:appJ-energy}, which translates by $cV$; the
normalized one used here, which translates by $c$; and the
Okounkov-body one of \Cref{warn:ch33-dh-normalizations}, which carries a
further $n!$.  They differ by positive constants and nothing else, but a
statement about moments, about Donaldson--Futaki invariants, or about
comparison with a complex functional is meaningless until one of the three
has been named.  The comparisons in this chapter carry these translation factors explicitly.
